\documentclass[12pt,a4paper,oneside]{book}
\usepackage[utf8]{inputenc}
\usepackage[T1]{fontenc}
\usepackage[margin=2.5cm,top=3cm,bottom=3cm]{geometry}
\usepackage{amsmath,amssymb,amsfonts}
\usepackage{booktabs}
\usepackage{longtable}
\usepackage{array}
\usepackage{multirow}
\usepackage{xcolor}
\usepackage{colortbl}
\usepackage{tabularx}
\usepackage{enumitem}
\usepackage{titlesec}
\usepackage{fancyhdr}
\usepackage{hyperref}
\usepackage{tcolorbox}
\usepackage{float}
\usepackage{graphicx}
\usepackage{setspace}
\usepackage{indentfirst}
\usepackage{emptypage}

\definecolor{primary}{RGB}{31, 73, 125}
\definecolor{secondary}{RGB}{217, 226, 243}
\definecolor{lightgray}{RGB}{242, 242, 242}
\definecolor{accent}{RGB}{198, 89, 17}
\definecolor{successgreen}{RGB}{40, 120, 70}

\hypersetup{
  colorlinks=true,
  linkcolor=primary,
  urlcolor=primary,
  citecolor=primary,
  pdftitle={Trabajo de Investigacion: Navegacion social robots agricolas},
  pdfauthor={R. J. Betancourt}
}

\titleformat{\chapter}[display]
  {\normalfont\Huge\bfseries\color{primary}}
  {\filright\Large\textsc{Capitulo} \thechapter}
  {1ex}
  {\titlerule\vspace{1ex}\filright}
  [\vspace{1ex}\titlerule]

\titleformat{\section}{\normalfont\Large\bfseries\color{primary}}{\thesection}{1em}{}
\titleformat{\subsection}{\normalfont\large\bfseries\color{primary}}{\thesubsection}{1em}{}
\titleformat{\subsubsection}{\normalfont\bfseries\color{primary}}{\thesubsubsection}{1em}{}

\pagestyle{fancy}
\fancyhf{}
\fancyhead[L]{\small\nouppercase{\leftmark}}
\fancyhead[R]{\small\thepage}
\renewcommand{\headrulewidth}{0.5pt}

\fancypagestyle{plain}{
  \fancyhf{}
  \fancyfoot[C]{\small\thepage}
  \renewcommand{\headrulewidth}{0pt}
}

\onehalfspacing

\newcommand{\figplaceholder}[2]{%
  \begin{figure}[H]
    \centering
    \fbox{
      \begin{minipage}{0.85\textwidth}
        \vspace{2cm}
        \centering
        \textcolor{accent}{\Large\textbf{[FIGURA #1: PENDIENTE]}}\\[0.5em]
        \textit{#2}
        \vspace{2cm}
      \end{minipage}
    }
    \caption{#2}
    \label{fig:#1}
  \end{figure}
}

\begin{document}

\begin{titlepage}
\begin{center}
\vspace*{1cm}
{\large \textbf{TRABAJO DE INVESTIGACION}}\\
\vspace{0.3cm}
{\small Programa de Investigacion en Robotica}\\
\vfill
\rule{\textwidth}{2pt}\\[0.4cm]
{\Huge\bfseries\color{primary}
Navegacion social de robots moviles agricolas en gemelos digitales:
}\\[0.4cm]
{\LARGE\bfseries\color{primary}
Comparativa estadistica de cinco arquitecturas de control en un entorno Unity-MATLAB
}\\[0.2cm]
\rule{\textwidth}{2pt}\\[1cm]
\vspace{0.5cm}
{\Large\textit{Respaldo de calculos, estado del arte y analisis multicriterio}}\\
\vspace{0.3cm}
{\large\textit{184 corridas $\cdot$ 5 metodos $\cdot$ 4 escenarios $\cdot$ 15 metricas}}\\
\vfill
{\large\textbf{Autor:}}\\[0.2cm]
{\Large R. J. Betancourt}\\[0.5cm]
{\large\textbf{Director:}}\\[0.2cm]
{\Large Dr. J. P. Vasconez}\\
\vfill
{\large\textbf{Documento integrado:}}\\
{\small Reune los siete documentos de respaldo del estudio en un unico compendio}\\
{\small con estado del arte, marco metodologico, resultados y bibliografia completa.}\\
\vfill
\rule{0.7\textwidth}{0.5pt}\\[0.4cm]
{\large \today}
\end{center}
\end{titlepage}

\thispagestyle{empty}
\cleardoublepage

\thispagestyle{empty}
\vspace*{4cm}
\begin{flushright}
\textit{``All models are wrong, but some are useful.''}\\[0.3em]
--- George E. P. Box (1976)\\[2em]
\textit{``En la simulacion uno aprende donde esta el problema;\\
en el campo, uno aprende cuan grande es realmente.''}\\[0.3em]
--- Adagio practico de robotica
\end{flushright}
\vfill
\newpage

\thispagestyle{plain}
\chapter*{Resumen}
\addcontentsline{toc}{chapter}{Resumen}

\noindent
\textbf{Problema y motivacion.} La navegacion de robots agricolas en corredores compartidos con agentes humanos requiere combinar planificacion geometrica con comportamiento socialmente aceptable, integrando seguridad social, eficiencia operacional y suavidad cinematica. La evaluacion de estrategias de navegacion social previa al despliegue fisico se realiza tipicamente en gemelos digitales, donde tanto el robot como los agentes humanos son entidades simuladas.

\noindent
\textbf{Metodos.} Se evaluan cinco arquitecturas de control: M0 (NavMesh puro), M1-M2 (NavMesh con supervisores binarios basados en distancia y orientacion), M3 (NavMesh con campo gaussiano anisotropico continuo) y M4 (M3 mas supervisor de modulacion de velocidad con $\alpha_{\min} > 0$ que evita bloqueo operacional). El estudio se construye sobre 184 corridas en un gemelo digital Unity-MATLAB que cubre cuatro escenarios tipicos de vinedo en hileras. Se aplican: (i) estadistica descriptiva exhaustiva sobre 15 metricas escalares, (ii) Kruskal-Wallis omnibus, (iii) Mann-Whitney U pareado con correccion Holm-Bonferroni y tamano de efecto Cliff's $\delta$, (iv) Fisher exacto y $\chi^2$ con Cramer's V, (v) ranking multicriterio TOPSIS con dos analisis paralelos (con y sin penalizacion de fallos), y (vi) analisis de sensibilidad Monte Carlo de los pesos del TOPSIS con 1000 muestras por configuracion.

\noindent
\textbf{Resultados.} Kruskal-Wallis confirma diferencias significativas ($p<0.001$) en 14 de 15 metricas. Las pruebas pareadas con Holm-Bonferroni muestran que M4 domina a todos los demas metodos en suavidad cinematica ($\delta < 0$, large) y supera a M0/M1/M2 en distancia minima preservada. Fisher exacto sobre el escenario critico E4 (humanos estaticos en pasillo congestionado) arroja $p = 7.14 \times 10^{-4}$ para M4 vs M1 y M4 vs M2, con $\chi^2$ global $= 26.09$ ($p<0.001$, Cramer's $V = 0.361$). TOPSIS posiciona a M4 como Rank \#1 con $CC = 0.804$ (sin penalizacion) y $CC = 0.801$ (con penalizacion), seguido muy cerca de M3. El analisis Monte Carlo confirma que M3 y M4 ocupan Top-2 en el 100\% de las 4000 evaluaciones realizadas, y M4 domina Rank \#1 con frecuencia 78.1\% bajo penalizacion de fallos.

\noindent
\textbf{Contribucion.} El estudio aporta una comparativa exhaustiva de cinco arquitecturas con analisis estadistico riguroso, identificacion explicita del bloqueo operacional de los supervisores binarios, y un protocolo multicriterio robusto frente a perturbaciones de pesos. Toda la metodologia se documenta paso a paso para reproducibilidad.

\noindent
\textbf{Palabras clave:} navegacion social, gemelo digital agricola, NavMesh, ORCA, campo gaussiano anisotropico, Kruskal-Wallis, Mann-Whitney, TOPSIS, Monte Carlo, ISO 18497.

\newpage

\chapter*{Lista de simbolos y nomenclatura}
\addcontentsline{toc}{chapter}{Lista de simbolos}

\begin{table}[H]
\centering
\renewcommand{\arraystretch}{1.4}
\small
\begin{tabular}{lll}
\toprule
\rowcolor{primary}
\textcolor{white}{\textbf{Simbolo}} &
\textcolor{white}{\textbf{Descripcion}} &
\textcolor{white}{\textbf{Unidades}} \\
\midrule
$d^\star$        & Distancia minima robot-humano por corrida & m \\
$d_I$            & Radio de zona intima (= 0.45) & m \\
$d_P$            & Radio de zona personal (= 1.20) & m \\
$d_S$            & Radio de zona social (= 3.60) & m \\
$d_{stop}$       & Distancia de parada del supervisor M4 (= 0.60) & m \\
$d_{resume}$     & Distancia de reanudacion del supervisor M4 (= 1.20) & m \\
$\alpha(d,\beta)$ & Funcion de modulacion de velocidad de M4 & --- \\
$\alpha_{\min}$  & Velocidad minima de escape no nula (= 0.15) & --- \\
$\beta$          & Angulo entre direccion robot y vector robot-humano & rad \\
$\gamma(\beta)$  & Factor angular de modulacion & --- \\
$\sigma_x$       & Desviacion gaussiana frontal (= 1.50) & m \\
$\sigma_y$       & Desviacion gaussiana lateral (= 0.90) & m \\
$\sigma_v$       & Desviacion estandar de la velocidad (suavidad) & m/s \\
$T_{personal}$   & Tiempo en zona personal por corrida & s \\
$T_{intima}$     & Tiempo en zona intima por corrida & s \\
$T_{mission}$    & Tiempo total de la corrida & s \\
$N_{stop}$       & Numero de paradas en una corrida & --- \\
$a_{max}$        & Aceleracion maxima (pseudo, numerica) & m/s2 \\
$\eta_s$         & Tasa de exito (success rate) & \% \\
$\eta_t$         & Eficiencia de trayectoria & --- \\
$v_{nom}$        & Velocidad nominal del robot (= 1.0) & m/s \\
$\mathbf{u}_{M4}$ & Vector direccion de salida del campo de M4 & --- \\
$H$              & Estadistico de Kruskal-Wallis & --- \\
$U$              & Estadistico de Mann-Whitney & --- \\
$\delta$         & Cliff's delta (tamano de efecto) & --- \\
$\chi^2$         & Estadistico chi-cuadrado & --- \\
$V$              & Cramer's V & --- \\
$CC_i$           & Closeness coefficient TOPSIS para alternativa $i$ & --- \\
$D^+$, $D^-$     & Distancias TOPSIS al ideal positivo y negativo & --- \\
\bottomrule
\end{tabular}
\end{table}

\newpage

\tableofcontents
\cleardoublepage
\listoftables
\cleardoublepage

\mainmatter

\part{Marco teorico y estado del arte}

\chapter{Introduccion y posicionamiento del estudio}

La integracion de robots moviles en entornos compartidos con personas constituye uno de los desafios centrales de la robotica contemporanea. A diferencia de las aplicaciones industriales tradicionales, donde los robots operan en celdas fisicamente segregadas, las aplicaciones emergentes en agricultura, atencion de salud, logistica y servicios requieren que el robot \emph{coexista, coopere y se anticipe} a la presencia humana. En el contexto particular de la agricultura, esta problematica se intensifica por la naturaleza estructurada pero parcialmente impredecible del entorno: corredores estrechos en hileras de vinedo o huertos, presencia de agentes humanos realizando tareas estaticas (poda, cosecha) o dinamicas (transporte, supervision), velocidad relativamente baja pero coexistencia prolongada, y exigencias normativas crecientes sobre la seguridad funcional de maquinas autonomas.

\section{Contexto del estudio: gemelo digital Unity-MATLAB}

Es fundamental enfatizar desde el inicio que \emph{todo el estudio se desarrolla integramente en simulacion}. No se realizan mediciones de campo. Especificamente:

\begin{itemize}
    \item El robot agricola es un agente simulado en Unity 3D con dinamica configurada en su NavMeshAgent.
    \item Los \emph{agentes humanos} son entidades simuladas con patrones de movimiento parametrizados.
    \item Las distancias, velocidades y aceleraciones se calculan a partir de los \texttt{transforms} de Unity.
    \item La telemetria se transmite via UDP a MATLAB para analisis estadistico off-line.
    \item El estudio cubre 184 corridas distribuidas en cinco metodos de control y cuatro escenarios.
\end{itemize}

Los rangos numericos adoptados para parametrizar el simulador (anchura de corredor 1.8-3.0 m, velocidades de robots agricolas 0.4-1.5 m/s) se inspiran como \emph{referencia descriptiva} en literatura previa de robotica agricola, sin constituir mediciones realizadas en este estudio.

\section{Objetivos}

\subsection{Objetivo general}

Disenar, implementar y evaluar comparativamente cinco arquitecturas de navegacion social para un robot agricola en un gemelo digital, identificando trade-offs entre seguridad social, eficiencia operacional y suavidad cinematica mediante analisis estadistico riguroso.

\subsection{Objetivos especificos}

\begin{enumerate}[leftmargin=*]
    \item Implementar las cinco arquitecturas de control en el gemelo digital Unity-MATLAB.
    \item Disenar el protocolo experimental con cuatro escenarios representativos de vinedo en hileras.
    \item Ejecutar 200 corridas planificadas (10 por celda metodo-escenario) y consolidar la telemetria completa.
    \item Aplicar estadistica descriptiva exhaustiva sobre 15 metricas escalares.
    \item Verificar diferencias significativas mediante Kruskal-Wallis omnibus, Mann-Whitney U pareado con correccion Holm-Bonferroni y Cliff's $\delta$.
    \item Cuantificar el bloqueo operacional con Fisher exacto y $\chi^2$.
    \item Sintetizar los resultados multicriterio con TOPSIS bajo dos analisis (con y sin penalizacion de fallos).
    \item Evaluar la robustez del ranking ante perturbaciones de pesos mediante simulacion Monte Carlo (1000 muestras x 4 configuraciones).
\end{enumerate}

\section{Hipotesis de trabajo}

\begin{table}[H]
\centering
\caption{Hipotesis del estudio y metodos estadisticos asociados.}
\renewcommand{\arraystretch}{1.4}
\small
\begin{tabular}{lp{8cm}l}
\toprule
\rowcolor{primary}
\textcolor{white}{\textbf{ID}} &
\textcolor{white}{\textbf{Hipotesis}} &
\textcolor{white}{\textbf{Test}} \\
\midrule
H1 & M0 (NavMesh puro) presenta corridas con invasion de zona intima ($T_{intima} > 0$). & K-W, M-W \\
H2 & M1 y M2 mejoran la distancia minima en E2 y E3, donde los humanos son dinamicos. & M-W pareado \\
H3 & M3 ofrece mayor distancia minima global que los demas. & K-W + post-hoc \\
H4 & M4 mejora la suavidad cinematica ($\sigma_v$) frente a todos los demas. & K-W + post-hoc \\
H5 & M4 reduce drasticamente el numero de paradas frente a M1 y M2. & K-W + post-hoc \\
H6 & M4 mejora la tasa de exito $\eta_s$ en E4 frente a M1 y M2. & Fisher, $\chi^2$ \\
\bottomrule
\end{tabular}
\end{table}

\chapter{Estado del arte}
\label{chap:estado_arte}




%==================================================================
\newpage
%==================================================================
\section{Proxémica humana y modelos de incomodidad}
\label{doc7:sec:proxemics}
%==================================================================

\subsection{Origen antropológico: Edward T. Hall}

El concepto fundamental de \emph{distancia interpersonal} se debe al antropólogo Edward T. Hall, quien en su obra seminal \emph{The Hidden Dimension}~\cite{Hall1966} introdujo cuatro zonas espaciales que rodean a un individuo y modulan su confort en función de la distancia y la naturaleza de la interacción:

\begin{table}[H]
\centering
\caption{Zonas proxémicas según Hall (1966).}
\renewcommand{\arraystretch}{1.3}
\begin{tabular}{lll}
\toprule
\rowcolor{primary}
\textcolor{white}{\textbf{Zona}} &
\textcolor{white}{\textbf{Rango (m)}} &
\textcolor{white}{\textbf{Contexto típico}} \\
\midrule
Íntima      & $0.00 - 0.45$ & contacto físico, parejas \\
Personal    & $0.45 - 1.20$ & conversación cercana, amigos \\
Social      & $1.20 - 3.60$ & interacciones formales, profesional \\
Pública     & $> 3.60$      & oradores, presentaciones \\
\bottomrule
\end{tabular}
\end{table}

Aunque la propuesta original era cualitativa y centrada en humano-humano, sus rangos numéricos se han adoptado ampliamente como guía operacional para robots móviles~\cite{Mead2017,Pacchierotti2005}. Es importante destacar dos limitaciones del modelo de Hall que la literatura posterior ha subrayado:

\begin{itemize}
    \item \textbf{Variabilidad cultural:} Aiello~\cite{Aiello1987,AielloAiello1974} demostró diferencias significativas en distancias preferidas según cultura (e.g., culturas mediterráneas vs anglosajonas).
    \item \textbf{Variabilidad contextual:} la distancia preferida varía con la velocidad de aproximación, el ángulo, la familiaridad previa, y la naturaleza de la actividad simultánea.
\end{itemize}

\subsection{Aplicación a la robótica social}

La adopción de proxémica de Hall en robótica móvil fue impulsada por trabajos como Pacchierotti et al.~\cite{Pacchierotti2005}, que evaluaron distancias preferidas en encuentros pasillo--humano, y Walters et al.~\cite{Walters2005}, que estudiaron preferencias de aproximación robot--humano. Estos estudios establecieron que, en general, los humanos prefieren que un robot \emph{respete al menos su zona personal} ($d \geq 1.20$~m) cuando éste se aproxima.

El estudio canónico contemporáneo de revisión es Mead \& Matarić~\cite{Mead2017}, que sintetiza tres décadas de investigación en proxémica robótica y propone un \emph{modelo computacional} de distancia preferida en función del contexto y de las características del robot. Rios-Martinez et al.~\cite{RiosMartinez2015} ofrecen una revisión específicamente orientada a navegación social, integrando proxémica con teoría de actividad situada.

\subsection{Modelos cuantitativos de incomodidad: trabajos recientes}

El estudio reciente más exhaustivo sobre incomodidad humana con robots móviles es Neggers et al.~\cite{Neggers2024,Neggers2022}, quienes condujeron un experimento controlado donde participantes calificaron su incomodidad ante un robot que se aproximaba o adelantaba a diferentes velocidades y distancias. Sus principales hallazgos:

\begin{enumerate}[leftmargin=*]
    \item La incomodidad puede aproximarse mediante \emph{gaussianas invertidas} con parámetros $(\sigma_x, \sigma_y)$ dependientes de la velocidad relativa y del ángulo.
    \item La velocidad del robot aumenta la incomodidad de forma no lineal: a mayor velocidad, mayor distancia mínima requerida para igualar el confort.
    \item La aproximación frontal es percibida como más invasiva que la lateral o la posterior.
\end{enumerate}

Su modelo, junto con el código fuente publicado por el grupo de Wengefeld et al.~\cite{Wengefeld2022}, ha sido adoptado como referencia para construir \emph{funciones de costo social} en planificadores recientes.

\subsection{Crítica a la aplicabilidad directa en agricultura}

Un punto importante para el contexto agrícola, que rara vez se discute explícitamente en la literatura proxémica, es que las distancias de Hall fueron \emph{calibradas para encuentros sociales y profesionales} en interiores. Su transferencia directa a entornos agrícolas con corredores estrechos (1.8--3.0~m de ancho) puede ser sub-óptima: en un viñedo, la zona social de Hall ($1.20$--$3.60$~m) literalmente no cabe en el corredor. Esto motiva una \emph{interpretación operacional} en lugar de normativa de la proxémica de Hall, como la que adoptamos en este estudio.

\newpage
%==================================================================
\section{Planificación de trayectoria clásica}
\label{doc7:sec:planning}
%==================================================================

\subsection{Algoritmos basados en grafos}

La búsqueda óptima en grafos discretos es la base teórica de la planificación clásica. Los dos algoritmos canónicos son:

\paragraph{Dijkstra (1959).} El algoritmo de Dijkstra~\cite{Dijkstra1959} resuelve el problema del camino más corto desde un nodo origen hacia todos los demás nodos en un grafo con pesos no negativos. Su complejidad con cola de prioridad es $O((|V|+|E|)\log|V|)$.

\paragraph{A* (Hart, Nilsson, Raphael, 1968).} El algoritmo A*~\cite{Hart1968} extiende Dijkstra introduciendo una función heurística admisible $h(n)$ que estima el costo restante hasta el objetivo. La función de evaluación $f(n)=g(n)+h(n)$ permite expansión informada y, bajo admisibilidad y consistencia de $h$, A* es óptimo y completo.

A* y sus variantes (Theta*~\cite{Daniel2010}, JPS~\cite{Harabor2011}, Hybrid A*~\cite{Dolgov2010}) son la base de la mayoría de planificadores globales en sistemas de navegación de producción.

\subsection{Algoritmos basados en muestreo}

Para configuraciones de alta dimensionalidad o espacios continuos complejos, los métodos de búsqueda determinística sufren la \emph{maldición de la dimensionalidad}. Los algoritmos basados en muestreo abordan esto:

\paragraph{Probabilistic Roadmaps (PRM).} Kavraki et al.~\cite{Kavraki1996} introdujeron PRM, que construye un grafo aproximando el espacio libre mediante muestreo aleatorio y conexión de vecinos cercanos válidos. PRM es asintóticamente completo.

\paragraph{Rapidly-Exploring Random Trees (RRT).} LaValle y Kuffner~\cite{LaValle1998,LaValle2001} propusieron RRT, que construye un árbol expandiéndose aleatoriamente desde la configuración inicial hacia la región objetivo. Particularmente útil en espacios de configuración no convexos o de alta dimensión (e.g., manipuladores).

\paragraph{RRT*.} Karaman y Frazzoli~\cite{Karaman2011} extendieron RRT a RRT*, garantizando \emph{optimalidad asintótica}: a medida que crece el número de muestras, el costo del camino tiende al óptimo. RRT* y sus variantes (Informed RRT*~\cite{Gammell2014}, BIT*~\cite{Gammell2015}) son referencia en planificación de manipulación y navegación móvil de alto rendimiento.

\subsection{NavMesh: navegación en superficies poligonales}

\textbf{Navigation Mesh (NavMesh)} es una representación del espacio libre como una colección de polígonos convexos conectados, originalmente popularizada en la industria de videojuegos por Snook~\cite{Snook2000}. Su atractivo: una vez construida, permite búsquedas A* extremadamente rápidas sobre el grafo dual de los polígonos, con triangulación dinámica para manejar obstáculos móviles.

Mononen~\cite{Mononen2013} desarrolló \texttt{Recast \& Detour}, la implementación open-source más utilizada y la base del NavMesh Agent del motor Unity, que es el sistema empleado como \emph{baseline geométrico} en el presente estudio. Sus características clave:

\begin{itemize}
    \item Construcción off-line del NavMesh estático a partir de geometría 3D.
    \item Búsqueda A* sobre el grafo de polígonos convexos.
    \item Smoothing path con string-pulling (Funnel algorithm).
    \item Evitación local mediante steering behaviors o RVO integrado.
\end{itemize}

NavMesh es ideal para entornos con geometría conocida y mayoritariamente estática, como escenas de juego o entornos agrícolas estructurados (corredores fijos), pero su limitación es que \emph{trata a los humanos como obstáculos cinemáticos puramente geométricos} sin componente social. Esta es precisamente la limitación que motiva los métodos socialmente conscientes (Sec.~\ref{doc7:sec:social}).

\subsection{Comparativa y lecciones aprendidas}

Sánchez-Ibáñez, Pérez-del-Pulgar y García-Cerezo~\cite{SanchezIbanez2021} ofrecen una revisión sistemática de planificación móvil que cubre estos métodos clásicos. Su conclusión es que \emph{ningún método único domina en todas las métricas}; la elección depende del trade-off entre completitud (clave en entornos cluttered), optimalidad asintótica (clave para minimizar tiempo de misión), velocidad de cómputo (clave en hardware embebido) y manejabilidad de obstáculos dinámicos.

Para el contexto agrícola específicamente, Mohanan y Salgoankar~\cite{Mohanan2018} subrayan que la \emph{repetitividad} del entorno (corredores casi idénticos) hace que los métodos basados en muestreo sean innecesariamente costosos: A* sobre NavMesh estático suele ser suficiente para la planificación global, dejando a un planificador local (Sec.~\ref{doc7:sec:reactive}) la tarea de ajuste reactivo.

\newpage
%==================================================================
\section{Métodos reactivos de evitación de obstáculos}
\label{doc7:sec:reactive}
%==================================================================

Los métodos reactivos operan en un \emph{horizonte temporal corto} (típicamente 0.1--2~s) ajustando la velocidad del robot para evitar colisiones con obstáculos detectados en el horizonte sensorial. Se ejecutan a alta frecuencia (10--50 Hz) y complementan al planificador global.

\subsection{Campos potenciales artificiales}

Khatib~\cite{Khatib1986} introdujo el concepto de \emph{Artificial Potential Field (APF)}: el robot se modela como una partícula moviéndose en un potencial donde el objetivo ejerce una atracción y los obstáculos ejercen repulsión. La velocidad del robot se determina como el gradiente negativo del potencial total.

\paragraph{Limitaciones reconocidas:}
\begin{itemize}
    \item \textbf{Mínimos locales:} cuando atracción y repulsión se cancelan en una posición que no es el objetivo, el robot queda atrapado.
    \item \textbf{Oscilaciones:} en pasajes estrechos donde dos repulsiones laterales se equilibran.
    \item \textbf{Inadecuación para humanos:} APF clásico trata humanos como obstáculos circulares isotrópicos, ignorando su orientación, velocidad y zonas proxémicas.
\end{itemize}

A pesar de estas limitaciones, las extensiones de APF con potenciales armónicos~\cite{Connolly1990}, navigation functions~\cite{Rimon1992} y campos sociales gaussianos anisotrópicos~\cite{Kirby2009,Pradalier2003} siguen siendo activas en la literatura actual.

\subsection{Dynamic Window Approach (DWA)}

Fox, Burgard y Thrun~\cite{Fox1997} propusieron el \emph{Dynamic Window Approach}, que introduce una idea fundamental: la búsqueda de velocidad óptima se realiza en el \emph{espacio de velocidades alcanzables} dado el modelo cinemático del robot, evaluando cada candidato $(v, \omega)$ con una función objetivo que combina:
\begin{equation}
J(v,\omega) = \alpha \cdot \text{heading}(v,\omega) + \beta \cdot \text{dist}(v,\omega) + \gamma \cdot \text{velocity}(v,\omega)
\end{equation}

DWA es computacionalmente eficiente (búsqueda exhaustiva sobre rejilla discretizada de velocidades) y ha sido la opción estándar en los sistemas ROS Navigation y Nav2 durante años. Sus desventajas principales son: (i) horizonte muy corto (un solo paso), (ii) no anticipa cambios de velocidad de obstáculos dinámicos, y (iii) sintonización de pesos sensible al entorno.

Variantes recientes incluyen Dynamic Adaptive DWA~\cite{Liang2024} con pesos adaptativos por fuzzy logic, que mejora desempeño en cruces complejos.

\subsection{Timed Elastic Band (TEB)}

Rösmann et al.~\cite{Rosmann2012,Rosmann2017} introdujeron \emph{Timed Elastic Band}, una formulación de optimización local donde la trayectoria se representa como una secuencia de poses con intervalos temporales asociados. La trayectoria óptima minimiza un funcional que penaliza:
\begin{itemize}
    \item Tiempo total de tránsito.
    \item Distancia a obstáculos (con coste exponencial cerca).
    \item Curvatura excesiva, jerk, aceleración.
    \item Restricciones de holonomicidad o no-holonomicidad.
\end{itemize}

TEB es ampliamente reconocido por producir \emph{trayectorias cinemáticamente suaves}, lo que lo hace particularmente adecuado para robots con dinámica significativa (e.g., vehículos no-holonómicos en altas velocidades). Comparativas en el ROS Navigation Stack~\cite{Bermudez2024,Yuan2022} muestran que TEB supera consistentemente a DWA en métricas de tiempo y suavidad, a costa de mayor coste computacional.

\paragraph{Variante consciente de humanos (Human-Aware TEB).} Kollmitz et al.~\cite{Kollmitz2021} extendieron TEB con una función de costo proxémica y predicción de movimiento humano, mostrando reducción significativa de invasiones de zona personal frente al TEB clásico.

\subsection{Velocity Obstacles y ORCA}

\paragraph{Velocity Obstacles (VO).} Fiorini y Shiller~\cite{Fiorini1993,Fiorini1998} introdujeron el concepto de \emph{velocity obstacle}, el conjunto de velocidades del robot que conducirían a colisión con un obstáculo móvil dentro de un horizonte $\tau$. La elección de cualquier velocidad fuera de la unión de VOs garantiza no-colisión durante $\tau$.

\paragraph{Reciprocal Velocity Obstacles (RVO).} Van den Berg, Lin y Manocha~\cite{VandenBerg2008} extendieron VO a multi-agente: si dos agentes son cooperativos, cada uno asume responsabilidad por la mitad de la maniobra evasiva. Esto reduce las oscilaciones del VO clásico.

\paragraph{Optimal Reciprocal Collision Avoidance (ORCA).} Van den Berg et al.~\cite{VandenBerg2011} propusieron ORCA, donde la solución de evitación se obtiene como \emph{programa lineal} sobre semiplanos: cada VO se aproxima por un semiplano, y la velocidad óptima del agente es la solución del LP que minimiza distancia a la velocidad preferida sujeto a todos los semiplanos. ORCA es decentralizado, eficiente en tiempo real y matemáticamente garantiza no-colisión en su horizonte.

ORCA es la base de Unity NavMeshAgent, lo que lo hace especialmente relevante para este estudio. Sus extensiones recientes incluyen ORCA-FL~\cite{London2025} con fuzzy logic, y ORCA con kinematic constraints~\cite{Alonso2018,Snape2011} para vehículos no-holonómicos.

\subsection{Comparativa entre familias reactivas}

\begin{table}[H]
\centering
\caption{Comparativa cualitativa entre familias de métodos reactivos.}
\renewcommand{\arraystretch}{1.3}
\small
\begin{tabular}{lcccc}
\toprule
\rowcolor{primary}
\textcolor{white}{\textbf{Aspecto}} &
\textcolor{white}{\textbf{APF}} &
\textcolor{white}{\textbf{DWA}} &
\textcolor{white}{\textbf{TEB}} &
\textcolor{white}{\textbf{ORCA}} \\
\midrule
Horizonte temporal       & 1 paso     & 1 paso       & multi-paso  & multi-paso \\
Coste computacional      & muy bajo   & bajo         & medio-alto  & medio \\
Suavidad cinemática      & media      & media        & alta        & alta \\
Manejo dinámico          & pobre      & local        & medio       & alto \\
Garantías formales       & ninguna    & ninguna      & opt.\ local & no-coll.\ \\
Multi-agente cooperativo & no         & no           & limitado    & sí \\
Implementación referencia& --         & ROS Nav      & ROS Nav     & Unity, RVO2 \\
\bottomrule
\end{tabular}
\end{table}

Singamaneni et al.~\cite{Singamaneni2024} concluyen que ningún método reactivo aislado es suficiente para navegación social robusta; la tendencia actual es \emph{combinar} múltiples paradigmas (e.g., NavMesh global + ORCA local + supervisor social adicional), que es precisamente la arquitectura híbrida del método M4 propuesto en este estudio.

\newpage
%==================================================================
\section{Modelos socialmente conscientes}
\label{doc7:sec:social}
%==================================================================

Los métodos reactivos clásicos (Sec.~\ref{doc7:sec:reactive}) tratan a los humanos como obstáculos cinemáticos puramente físicos. Los métodos \emph{socialmente conscientes} agregan modelado explícito de las propiedades sociales del humano: zona personal, orientación corporal, intención inferida.

\subsection{Modelo de fuerza social}

Helbing y Molnar~\cite{Helbing1995} propusieron el seminal \emph{Social Force Model (SFM)}, originalmente para modelar dinámica de peatones, que describe el movimiento de cada peatón como sujeto a tres fuerzas:
\begin{equation}
\mathbf{F}_i = \mathbf{F}_i^{\text{desired}} + \sum_{j \neq i} \mathbf{F}_{ij}^{\text{repulsive}} + \sum_W \mathbf{F}_{iW}^{\text{wall}}
\end{equation}

El SFM ha tenido enorme impacto: 
\begin{itemize}
    \item Captura naturalmente fenómenos colectivos (carriles auto-organizados, formación de cuellos de botella).
    \item Es físicamente interpretable (Newton + términos psicológicos).
    \item Es computacionalmente liviano: $O(N)$ por agente con vecindad espacial.
\end{itemize}

Las extensiones más relevantes son: Helbing, Farkas y Vicsek~\cite{Helbing2000,Helbing2002} (comportamiento de pánico), Moussaïd et al.~\cite{Moussaid2009,Moussaid2010} (mecanismos cognitivos subyacentes), Pellegrini et al.~\cite{Pellegrini2009} (modelos de evitación con horizonte futuro).

\paragraph{Aplicación a robots móviles.} Ferrer, Garrell y Sanfeliu~\cite{Ferrer2013,Ferrer2014} fueron pioneros en aplicar SFM a navegación robótica social en entornos urbanos. Luber et al.~\cite{Luber2010} integraron SFM con tracking de peatones por LiDAR.

\subsection{Campos potenciales gaussianos anisotrópicos}

Una variante ampliamente utilizada en robótica social es el \emph{campo gaussiano anisotrópico} centrado en la persona, que en lugar de tratar al humano como obstáculo circular isotrópico, modela su zona personal como una elipse alineada con su orientación corporal:
\begin{equation}
C(\mathbf{p}_r, \mathbf{p}_h, \theta_h) = \exp\left( -\frac{1}{2}\left[ \frac{x'^2}{\sigma_x^2} + \frac{y'^2}{\sigma_y^2} \right] \right)
\label{doc7:eq:gaussian_aniso}
\end{equation}
donde $(x', y')$ es la posición del robot en el sistema local del humano (rotado por $\theta_h$), y $\sigma_x > \sigma_y$ refleja que la zona frontal es típicamente más extensa que la lateral.

\paragraph{Trabajos seminales:} Pradalier, Hermosillo y Laugier~\cite{Pradalier2003} y Sisbot et al.~\cite{Sisbot2007} popularizaron este modelo en planificación humana-aware. Kirby, Simmons y Forlizzi~\cite{Kirby2009} lo formalizaron y demostraron experimentalmente que la asimetría frontal/lateral es predictiva de incomodidad humana.

\paragraph{Calibración empírica:} Como ya se mencionó, Neggers et al.~\cite{Neggers2024} calibraron empíricamente los parámetros $(\sigma_x, \sigma_y)$ a partir de experimentos controlados, validando que la \emph{gaussiana invertida} (i.e., $1-C$ como costo de incomodidad) ajusta bien los datos para diferentes velocidades. Sus valores de referencia: $\sigma_x \approx 1.5$~m, $\sigma_y \approx 0.9$~m, son los adoptados en M4 del presente estudio.

\subsection{Predicción del movimiento humano}

Un componente crucial complementario a los campos sociales es la predicción de la trayectoria humana futura, lo que permite al robot anticiparse en lugar de reaccionar.

\begin{itemize}
    \item \textbf{Procesos Gaussianos (GP).} Trautman, Ma, Murray y Krause~\cite{Trautman2013,Trautman2015} introdujeron \emph{Interacting Gaussian Processes (IGP)}, que modelan trayectorias humanas como GPs interactivos, capturando la cooperación en multitudes.
    \item \textbf{Social-LSTM.} Alahi et al.~\cite{Alahi2016} propusieron Social-LSTM, una red recurrente que aprende interacciones entre peatones implícitamente vía social pooling.
    \item \textbf{Social-GAN, Trajectron++.} Gupta et al.~\cite{Gupta2018} y Salzmann et al.~\cite{Salzmann2020} introdujeron predicción multi-modal con GAN y conditional VAE respectivamente.
    \item \textbf{Modelos de difusión.} Trabajos muy recientes~\cite{Mao2024} aplican diffusion models a predicción de trayectorias, capturando incertidumbre multimodal de forma natural.
\end{itemize}

Para aplicaciones agrícolas con humanos relativamente estáticos o moviéndose en patrones repetitivos (patrullaje), los modelos sofisticados de predicción son típicamente innecesarios; basta con asumir velocidad constante o un modelo cinemático simple~\cite{Schulz2003}.

\subsection{Modelos basados en reglas y máquinas de estado}

Una alternativa pragmática a los campos potenciales son los \emph{supervisores de regla} o \emph{máquinas de estado finitas} que monitorean condiciones específicas (e.g., distancia $<$ umbral, ángulo de aproximación frontal) y conmutan entre comportamientos discretos (i.e., normal $\to$ frenar $\to$ esperar $\to$ reanudar). Trabajos representativos: Lichtenthäler, Lorenz y Kirsch~\cite{Lichtenthaler2012}, Charalampous, Kostavelis y Gasteratos~\cite{Charalampous2017}.

Estos métodos ofrecen mayor interpretabilidad y depuración pero, como muestran nuestros resultados con M1 y M2, sufren de \emph{bloqueo operacional} ante humanos estáticos en pasillos congestionados: el robot se detiene esperando una señal de "vía libre" que nunca llega.

\newpage
%==================================================================
\section{Métodos avanzados: MPC y aprendizaje profundo}
\label{doc7:sec:advanced}
%==================================================================

\subsection{Model Predictive Control (MPC)}

El \emph{Model Predictive Control} aplicado a navegación social formula el problema como una optimización por horizonte deslizante: en cada instante, se predicen los próximos $N$ pasos del estado del robot y de los humanos, se optimiza la secuencia de control que minimiza un funcional combinando objetivos de seguimiento y restricciones de seguridad social, y se aplica solo el primer paso de control. Trabajos canónicos:

\begin{itemize}
    \item Brito, Floor, Ferranti y Alonso-Mora~\cite{Brito2019} propusieron \emph{Model Predictive Contouring Control}, que separa seguimiento de path y velocidad para navegación en entornos no estructurados.
    \item Schoels et al.~\cite{Schoels2020} extendieron MPC a navegación social con predicción de peatones por GP.
    \item Kamenev et al.~\cite{Kamenev2023} integraron Social-LSTM con MPC en simulaciones de alta densidad de peatones.
    \item Shamsah et al.~\cite{Shamsah2024} propusieron MPC con redes neuronales de zonotopos para robots bípedos en multitudes.
\end{itemize}

\paragraph{Ventajas:} MPC ofrece \emph{garantías de seguridad explícitas} (las restricciones aparecen en la formulación), maneja naturalmente dinámicas no lineales, y se beneficia directamente de cualquier mejora en predicción humana.

\paragraph{Desventajas:} costo computacional alto (típicamente requiere optimización no convexa), sensibilidad a la calidad del modelo predictivo, y dificultad de garantizar feasibilidad bajo perturbaciones grandes.

\subsection{Aprendizaje por refuerzo (RL)}

El RL aprende políticas de navegación social a partir de interacción simulada. Trabajos pioneros:

\begin{itemize}
    \item Chen et al.~\cite{Chen2017,Chen2019} propusieron \emph{CADRL} (Collision Avoidance with Deep RL) y SA-CADRL (Socially Aware), que aprenden políticas decentralizadas de evitación cooperativa.
    \item Everett, Chen y How~\cite{Everett2018} introdujeron \emph{GA3C-CADRL} con redes convolucionales para variable número de agentes.
    \item Liu et al.~\cite{Liu2023}, Chen et al.~\cite{Chen2023} aplicaron Graph Neural Networks (GNN) a navegación social.
\end{itemize}

\paragraph{Crítica reciente.} Karwowski et al.~\cite{Karwowski2024} en su revisión benchmarking de métodos basados en aprendizaje muestran que, en entornos no vistos durante el entrenamiento, los métodos RL pueden colapsar mientras que los métodos clásicos (incluyendo SFM) se degradan suavemente. La conclusión emergente es que \emph{métodos híbridos} (planificadores clásicos con módulos aprendidos) son más robustos que el extremo puro.

\subsection{Posicionamiento del estudio actual}

El método M4 propuesto en el presente estudio se inscribe en la familia de \emph{arquitecturas híbridas}: combina (i) un planificador global geométrico (NavMesh A*), (ii) un evitador local reactivo (ORCA integrado en NavMeshAgent), y (iii) un \emph{supervisor social continuo} basado en campos gaussianos anisotrópicos~\cite{Kirby2009,Neggers2024}. Esta arquitectura busca el balance entre interpretabilidad (todos los componentes son explicables), robustez (no depende del re-entrenamiento ante distribución cambiante) y desempeño (los resultados estadísticos del estudio confirman esta hipótesis).

\newpage
%==================================================================
\section{Robótica agrícola y seguridad funcional}
\label{doc7:sec:agri}
%==================================================================

\subsection{Tendencias generales}

La robótica agrícola es un sector de crecimiento exponencial. Vasconez, Kantor y Auat-Cheein~\cite{Vasconez2019} ofrecen una revisión comprensiva de las aplicaciones (poda, cosecha, monitoreo, fumigación) y las plataformas robóticas (UGV, UAV, brazos manipuladores) en horticultura, viticultura y agricultura extensiva.

Tsolakis et al.~\cite{Tsolakis2024} actualizan esta perspectiva con énfasis en agricultura de precisión y robótica colaborativa, identificando tres barreras principales para la comercialización masiva: (i) coste, (ii) interoperabilidad y, crucialmente, (iii) \emph{seguridad funcional ante presencia humana}.

\subsection{Adaptaciones del Nav2 stack para entornos agrícolas}

Para robots móviles autónomos basados en ROS 2, el stack \emph{Nav2}~\cite{Macenski2020} es la opción de facto. Sin embargo, su uso directo en huertos y viñedos presenta desafíos específicos:

\begin{itemize}
    \item Sani, Sgorbissa y Carpin~\cite{Sani2024} reportaron que la maleza alta y los pastos producen falsos positivos en costmaps de Nav2 con LiDAR, forzando al robot a tomar caminos innecesariamente largos. Su solución: extensiones que clasifican obstáculos como \emph{traversables} basándose en sensores adicionales (cámara, IMU).
    \item Reina et al.~\cite{Reina2017} desarrollaron sistemas de fusión multi-sensor para clasificación de terreno transitable.
\end{itemize}

\subsection{Marco normativo: ISO 18497 y series asociadas}

La seguridad funcional de máquinas agrícolas autónomas está regulada por la familia de normas ISO 18497~\cite{ISO18497}, originalmente publicada en 2018 y revisada en 2024-2025. Aspectos clave:

\begin{itemize}
    \item \textbf{ISO 18497-1:} Principios generales de diseño para máquinas agrícolas autónomas. Establece niveles de autonomía y categorías de riesgo.
    \item \textbf{ISO 18497-2:} Requerimientos para sistemas de detección de personas y obstáculos.
    \item \textbf{ISO 18497-3:} Verificación y validación de sistemas de seguridad.
    \item \textbf{ISO 18497-4:} Requisitos para supervisión humana remota.
\end{itemize}

Bispo et al.~\cite{Bispo2023} y Roozemond et al.~\cite{Roozemond2024} ofrecen revisiones sistemáticas sobre la suficiencia de ISO 18497, concluyendo que las directrices son necesarias pero \emph{no suficientes} para garantizar seguridad: requieren ser complementadas con análisis de riesgo específicos del contexto (FMEA, HAZOP) y con verificación experimental en gemelos digitales.

\paragraph{Implicación para nuestro estudio.} La distancia $d_P = 1.20$~m que adoptamos como zona personal en el estudio se inspira en proxémica de Hall, no en una certificación normativa ISO. Como aclara Spyrosoft~\cite{Spyrosoft2024}, la transferencia de un parámetro de simulación a un sistema certificable bajo ISO 18497 requiere análisis de riesgo adicional, sensórica certificada y validación de campo --- aspectos identificados explícitamente como \emph{trabajo futuro} en nuestro manuscrito.

\subsection{Ejemplos de robots agrícolas comerciales}

Para contextualizar las velocidades operacionales referenciadas en el estudio:
\begin{itemize}
    \item \textbf{Naïo Oz y Dino:} robots de desherbado, velocidad operacional $0.4$--$1.0$~m/s~\cite{Tsolakis2024}.
    \item \textbf{Burro AI:} robot de transporte en viñedos, $1.5$~m/s nominal.
    \item \textbf{VineRobot, VineScout:} plataformas de monitoreo en viñedos, $1.0$--$1.4$~m/s.
\end{itemize}

Estos rangos son los que adoptamos como \emph{referencia descriptiva} para parametrizar el simulador (ver manuscrito v3, Sec.~1.2), enfatizando que los valores configurados no constituyen mediciones de campo realizadas en este estudio.

\newpage
%==================================================================
\section{Gemelos digitales y Sim2Real}
\label{doc7:sec:digital}
%==================================================================

\subsection{Definición y evolución}

Un \emph{gemelo digital (digital twin)} es una replica virtual de un sistema físico con sincronización bidireccional con su contraparte real. El concepto se atribuye a Grieves~\cite{Grieves2014} en el contexto de manufactura inteligente, y ha evolucionado con creciente fidelidad gracias a motores de física y rendering de alta calidad (Unity, Unreal Engine, Isaac Sim).

\paragraph{Niveles de gemelo digital.} Tao y Qi~\cite{Tao2018,Tao2019} proponen una jerarquía:
\begin{enumerate}
    \item \textbf{Modelo virtual:} replica geométrica/cinemática estática.
    \item \textbf{Sombra digital (digital shadow):} flujo unidireccional físico $\to$ virtual.
    \item \textbf{Gemelo digital completo:} flujo bidireccional con co-simulación.
\end{enumerate}

El presente estudio se ubica en el segundo nivel: una sombra digital de alta fidelidad donde el comportamiento del controlador se evalúa exhaustivamente antes de la transferencia a hardware físico (\emph{trabajo futuro}).

\subsection{Aplicaciones en robótica}

Han, Park y Kim~\cite{Han2025} y Chang et al.~\cite{Chang2025} demuestran arquitecturas de gemelo digital basadas en Unity3D para robots industriales, con comunicación TCP/IP en tiempo real entre el simulador y el robot físico. La arquitectura de comunicación UDP utilizada en este estudio (Unity $\to$ MATLAB) sigue precisamente este patrón.

\subsection{Sim2Real: cerrando la brecha de realidad}

La transferencia de comportamientos aprendidos o validados en simulación a hardware real es conocida como \emph{Sim2Real} y constituye un área de investigación activa. Las técnicas principales son:

\begin{itemize}
    \item \textbf{Domain randomization} (Tobin et al.~\cite{Tobin2017}): aleatoriza parámetros del simulador (texturas, masa, fricción) para que la política aprendida sea robusta a la transferencia.
    \item \textbf{System identification} (SysID): mide parámetros físicos reales y los inyecta en el simulador para minimizar la brecha.
    \item \textbf{Domain adaptation} (Gupta et al.~\cite{GuptaDA2018}): ajuste fino con datos reales sobre una política pre-entrenada en simulación.
    \item \textbf{Photorealistic digital twins} (Yan et al.~\cite{Yan2025}): construcción de simulaciones con fidelidad visual extrema usando NeRF/3D Gaussian Splatting.
\end{itemize}

\paragraph{Aplicación reciente en agricultura.} Yan et al.~\cite{Yan2025} demuestran un gemelo digital fotorrealista de un campo de fresas comercial, donde se entrena un detector de fruta exclusivamente en datos sintéticos y se transfiere al campo real con F1-scores de 0.81--0.92 sin necesidad de imágenes reales para entrenamiento. Este resultado avala la hipótesis de que un gemelo digital agrícola bien construido puede ser una plataforma legítima de evaluación de algoritmos antes del despliegue de campo.

\paragraph{Limitación operacional.} Para algoritmos de \emph{control} y \emph{percepción dinámica}, la brecha entre simulación y realidad es típicamente más severa que para algoritmos de \emph{visión estática}. Esto motiva la doble validación: (i) en gemelo digital exhaustivo (este estudio); (ii) en hardware físico (trabajo futuro).

\newpage
%==================================================================
\section{Métodos estadísticos no paramétricos}
\label{doc7:sec:stats}
%==================================================================

\subsection{Tests de comparación de distribuciones}

\paragraph{Mann-Whitney U.} Mann y Whitney~\cite{Mann1947} introdujeron el test U para comparar dos muestras independientes sin asumir normalidad. Es la versión no paramétrica del t-test independiente.

\paragraph{Kruskal-Wallis.} Kruskal y Wallis~\cite{Kruskal1952} generalizaron Mann-Whitney a $k \geq 3$ grupos. La generalización es directa: rankear conjuntamente todas las observaciones y calcular el estadístico:
\begin{equation}
H = \frac{12}{N(N+1)} \sum_j \frac{R_j^2}{n_j} - 3(N+1)
\end{equation}
Bajo $H_0$, $H$ sigue $\chi^2_{(k-1)}$ asintóticamente. Es la versión no paramétrica de ANOVA de una vía.

\paragraph{Cuándo usar tests no paramétricos.} le Cessie, Goeman y Dekkers~\cite{leCessie2020} ofrecen una guía pragmática: las pruebas no paramétricas son preferibles cuando (i) las distribuciones son claramente no normales, (ii) hay valores atípicos influyentes, (iii) las muestras son pequeñas, o (iv) los datos son ordinales. La pérdida de potencia frente al test paramétrico equivalente es típicamente menor al 5\% cuando los datos son normales (eficiencia relativa asintótica $3/\pi \approx 0.955$ para Mann-Whitney vs t).

\subsection{Corrección de comparaciones múltiples}

Cuando se realizan $k$ tests pareados, la probabilidad de cometer al menos un error tipo I (FWER) crece como $1-(1-\alpha)^k$. Las correcciones más utilizadas:

\paragraph{Bonferroni.} Multiplica cada $p$-valor por $k$. Conservador pero simple.

\paragraph{Holm-Bonferroni.} Holm~\cite{Holm1979} propuso un step-down: ordenar $p$-valores ascendentemente y multiplicar el $i$-ésimo por $(k-i+1)$. Uniformemente más potente que Bonferroni.

\paragraph{Benjamini-Hochberg (FDR).} Benjamini y Hochberg~\cite{Benjamini1995} controlan la \emph{tasa de falso descubrimiento} en lugar del FWER, aceptando algunos falsos positivos a cambio de mayor potencia.

\paragraph{Dunn's test.} Dunn~\cite{Dunn1964} desarrolló específicamente un post-hoc para Kruskal-Wallis con corrección integrada.

En el presente estudio se adopta Holm-Bonferroni por ser uniformemente más potente que Bonferroni clásico mientras controla el FWER al mismo nivel.

\subsection{Tamaño de efecto: Cliff's delta}

Cliff~\cite{Cliff1993,Cliff1996} introdujo $\delta$ como medida no paramétrica de tamaño de efecto:
\begin{equation}
\delta(A,B) = P(X_A > X_B) - P(X_A < X_B)
\end{equation}

La interpretación de magnitudes recomendada es de Romano et al.~\cite{Romano2006}: $|\delta| < 0.147$ negligible, $< 0.33$ small, $< 0.474$ medium, $\geq 0.474$ large.

Cliff's $\delta$ se prefiere a otras medidas como Cohen's $d$ o Hedge's $g$ cuando los datos son no paramétricos, ya que no requiere homogeneidad de varianzas ni normalidad.

\subsection{Tablas de contingencia}

\paragraph{Test exacto de Fisher.} Fisher~\cite{Fisher1922} introdujo el test exacto basado en la distribución hipergeométrica para tablas $2 \times 2$. Es el test de elección cuando los conteos esperados son $< 5$ en alguna celda, condición típica en muestras pequeñas.

\paragraph{Chi-cuadrado de Pearson.} Pearson~\cite{Pearson1900} propuso el estadístico $\chi^2$ que se aproxima asintóticamente por la distribución $\chi^2_{((r-1)(c-1))}$.

\paragraph{Cramér's V.} Cramér~\cite{Cramer1946} introdujo $V$ como medida de tamaño de efecto para tablas de contingencia: $V = \sqrt{\chi^2 / (N \cdot \min(r-1,c-1))}$. Su interpretación según Rea y Parker~\cite{Rea1992}: $V<0.10$ negligible, $<0.20$ weak, $<0.40$ moderate, $<0.60$ relatively strong, $\geq 0.60$ very strong.

\subsection{Verificación cruzada y reproducibilidad}

Conviene mencionar que todos los cálculos estadísticos del presente estudio se verifican cruzadamente entre dos implementaciones independientes:
\begin{enumerate}
    \item \texttt{scipy.stats} (Python, versión 1.14+).
    \item Cálculos manuales paso a paso (documentados en Documentos 2--4 del respaldo).
\end{enumerate}
Las diferencias máximas entre implementaciones son del orden de $10^{-10}$, atribuibles a diferencias en el manejo numérico de empates~\cite{VirtanenScipy2020}.

\newpage
%==================================================================
\section{Análisis multicriterio TOPSIS}
\label{doc7:sec:topsis}
%==================================================================

\subsection{Fundamento teórico}

Hwang y Yoon~\cite{HwangYoon1981} introdujeron \emph{TOPSIS} como un método de decisión multicriterio basado en el principio de \emph{distancias geométricas}: la mejor alternativa es la que minimiza la distancia a una solución ideal positiva (PIS) y maximiza la distancia a una solución ideal negativa (NIS). El procedimiento de 9 pasos se documenta exhaustivamente en el Documento 5 del respaldo.

\subsection{Variantes y extensiones}

\begin{itemize}
    \item \textbf{Fuzzy TOPSIS.} Chen y Hwang~\cite{ChenHwang1992} extendieron TOPSIS a entornos difusos donde los datos de entrada son números fuzzy.
    \item \textbf{AHP-TOPSIS.} Onder y Dag~\cite{Onder2013} popularizaron la combinación de AHP (para obtener pesos) y TOPSIS (para ranking).
    \item \textbf{Entropy-TOPSIS.} Wu et al.~\cite{Wu2020} usan entropía de Shannon para asignar pesos automáticamente.
    \item \textbf{Distributed TOPSIS.} El Hadrami et al.~\cite{ElHadrami2023} para Big Data con paralelización.
    \item \textbf{Random Hypergraph TOPSIS.} Bhargava et al.~\cite{Bhargava2024} para criterios interdependientes.
\end{itemize}

\subsection{Análisis de sensibilidad}

Olson~\cite{Olson2004} fue pionero en el análisis de sensibilidad de TOPSIS, mostrando que el ranking puede ser sensible a (i) la elección del método de normalización, (ii) la selección de pesos, y (iii) la presencia de alternativas extremas. Este resultado motiva el análisis Monte Carlo del Documento 6 del respaldo.

Saltelli et al.~\cite{Saltelli2008} ofrecen el marco de referencia general para análisis de sensibilidad global, que ha sido adaptado a TOPSIS por trabajos como~\cite{Yazdani2017} con índices de Sobol.

\subsection{Crítica y método alternativo}

Una crítica al TOPSIS clásico es que su \emph{método de normalización vectorial} (división por la norma euclidiana de cada columna) puede introducir distorsiones cuando las escalas de los criterios son muy diferentes. Methods como ELECTRE~\cite{Roy1991}, PROMETHEE~\cite{Brans1986} y VIKOR~\cite{Opricovic2004} ofrecen alternativas con diferentes principios subyacentes (relaciones de superación, comparaciones por flujos netos, regret minimax).

Sin embargo, TOPSIS sigue siendo el método más usado en literatura aplicada por su simplicidad, interpretabilidad y robustez empírica~\cite{Velasquez2013,Mathew2023}.

\newpage


\part{Diseno experimental y datos}

\chapter{Arquitecturas de control evaluadas}
\label{chap:metodos}

Esta seccion describe formalmente las cinco arquitecturas de control evaluadas. Todas comparten el mismo \emph{planificador global geometrico} (NavMesh A* en Unity) y difieren en su capa local: presencia o ausencia de supervisores sociales, naturaleza binaria o continua del campo social, y existencia o no de un modulo de modulacion de velocidad.

\section{M0: NavMesh puro (baseline geometrico)}

El metodo M0 utiliza unicamente el sistema NavMeshAgent de Unity, que combina planificacion A* sobre la malla de navegacion poligonal con evitacion local ORCA integrada. No incluye ningun modulo social: los agentes humanos son tratados como obstaculos cinematicos circulares isotropicos.

\section{M1: Supervisor binario por distancia}

Sobre el baseline M0, M1 agrega un supervisor binario que detecta presencia humana dentro de un radio critico $d_{crit}$. Si la condicion se cumple, el robot detiene su avance hasta que la distancia se restablece sobre $d_{crit} + \Delta$ (histeresis para evitar conmutacion rapida).

\section{M2: Supervisor binario por distancia y orientacion}

M2 extiende M1 incorporando informacion angular: el robot solo se detiene si el agente humano esta aproximandose en un cono frontal definido por $\theta_{crit}$. Esta extension busca reducir las paradas innecesarias del M1 ante humanos que se alejan o pasan de costado.

\section{M3: Campo gaussiano anisotropico continuo}

M3 reemplaza el supervisor binario por un campo de costo continuo basado en gaussianas anisotropicas centradas en cada agente humano simulado:

\begin{equation}
C_i^{ani}(\mathbf{p}_r, \mathbf{p}_{h_i}, \theta_{h_i}) = \exp\left(-\frac{1}{2}\left[\frac{x'^2}{\sigma_x^2} + \frac{y'^2}{\sigma_y^2}\right]\right)
\end{equation}

donde $(x', y')$ son las coordenadas del robot en el sistema local del humano. Los parametros adoptados ($\sigma_x = 1.50$ m, $\sigma_y = 0.90$ m) provienen de la calibracion empirica de Neggers et al. (2024).

\section{M4: M3 con supervisor de modulacion de velocidad}

M4 conserva la direccion de M3 y agrega un modulo de \emph{modulacion de velocidad escalar}:

\begin{equation}
\alpha(d,\beta) = \begin{cases}
0, & d \leq d_I \\
\alpha_{\min}, & d_I < d \leq d_{stop} \\
\alpha_{\min} + (1 - \alpha_{\min})\dfrac{d - d_{stop}}{d_{resume} - d_{stop}}\gamma(\beta), & d_{stop} < d < d_{resume} \\
\gamma(\beta), & d_{resume} \leq d < d_S \\
1, & d \geq d_S
\end{cases}
\end{equation}

con $\alpha_{\min} = 0.15$. La velocidad final del robot es:

\begin{equation}
\mathbf{v}_r = \mathrm{sat}_{v_{nom}}(\alpha \cdot v_{nom} \cdot \mathbf{u}_{M4})
\end{equation}

\paragraph{Justificacion de $\alpha_{\min} > 0$.} Este parametro garantiza una \emph{velocidad minima de escape no nula} cuando el robot se encuentra en la franja $d_I < d \leq d_{stop}$. Sin esta correccion, el robot quedaria detenido pese a tener una direccion $\mathbf{u}_{M4}$ no nula, generando bloqueo operacional ante agentes humanos simulados estaticos.

\chapter{Diseno experimental y registro de datos}
\label{chap:experimental}

\section{Escenarios}

\begin{table}[H]
\centering
\caption{Cuatro escenarios experimentales evaluados.}
\renewcommand{\arraystretch}{1.4}
\small
\begin{tabular}{llp{7cm}}
\toprule
\rowcolor{primary}
\textcolor{white}{\textbf{ID}} &
\textcolor{white}{\textbf{Nombre}} &
\textcolor{white}{\textbf{Descripcion}} \\
\midrule
E1 & Pasillo libre & Sin agentes humanos. Linea base de eficiencia. \\
E2 & Cruce frontal & Un agente humano simulado avanza en sentido contrario. \\
E3 & Adelantamiento & Un agente humano simulado se mueve mas lento, mismo sentido. \\
E4 & Pasillo congestionado con humanos estaticos & Dos agentes humanos simulados estaticos en el corredor. \\
\bottomrule
\end{tabular}
\end{table}

\section{Factores aleatorizados entre corridas}

\begin{table}[H]
\centering
\caption{Factores parametrizados aleatoriamente en el simulador.}
\renewcommand{\arraystretch}{1.3}
\small
\begin{tabular}{lll}
\toprule
\rowcolor{primary}
\textcolor{white}{\textbf{Factor del simulador}} &
\textcolor{white}{\textbf{Distribucion}} &
\textcolor{white}{\textbf{Rango configurado}} \\
\midrule
Velocidad asignada al agente humano simulado & uniforme & $[0.4, 1.2]$ m/s \\
Offset lateral inicial del agente             & uniforme & $\pm 0.3$ m \\
Fase inicial de patrulla del agente           & uniforme & $[0, 2\pi]$ \\
Tiempo de aparicion del agente en la escena   & uniforme & $[5, 15]$ s \\
Posicion inicial del robot simulado           & uniforme & $\pm 0.2$ m alrededor de $A$ \\
Semilla del NavMeshAgent                      & aleatoria & --- \\
\bottomrule
\end{tabular}
\end{table}

\figplaceholder{esquema_corredor}{Esquema del corredor de vinedo simulado con puntos $A$ (inicio) y $B$ (objetivo).}

\figplaceholder{escenarios_E1_E4}{Vista superior de los cuatro escenarios E1-E4.}

\figplaceholder{trayectorias_E4}{Comparativa de trayectorias en E4 para los cinco metodos.}

\figplaceholder{distancia_tiempo_E4}{Curvas distancia-tiempo en E4 con lineas de referencia $d_I, d_P, d_S$.}

\figplaceholder{velocidad_tiempo_E4}{Curvas velocidad-tiempo en E4.}

\figplaceholder{boxplots_metricas}{Boxplots comparativos de las metricas $d^\star$, $T_{personal}$, $\sigma_v$ y $N_{stop}$.}

\section{Tamanos muestrales realizados}

\begin{table}[H]
\centering
\caption{Tamanos muestrales por celda metodo-escenario.}
\renewcommand{\arraystretch}{1.3}
\begin{tabular}{lcccc|c}
\toprule
\rowcolor{primary}
\textcolor{white}{\textbf{Metodo}} & \textcolor{white}{\textbf{E1}} & \textcolor{white}{\textbf{E2}} & \textcolor{white}{\textbf{E3}} & \textcolor{white}{\textbf{E4}} & \textcolor{white}{\textbf{Total}} \\
\midrule
M0 & 10 & 10 & 10 & 10 & \textbf{40} \\
M1 & 10 & 10 & 10 & 2  & \textbf{32} \\
M2 & 10 & 10 & 10 & 2  & \textbf{32} \\
M3 & 10 & 10 & 10 & 10 & \textbf{40} \\
M4 & 10 & 10 & 10 & 10 & \textbf{40} \\
\midrule
\textbf{Realizadas} & 50 & 50 & 50 & 34 & \textbf{184} \\
\textbf{Planificadas} & 50 & 50 & 50 & 50 & \textbf{200} \\
\bottomrule
\end{tabular}
\end{table}

\part{Analisis estadistico inferencial}

\chapter{Estadistica descriptiva}
\label{chap:descriptiva}


\vspace{1em}



%==================================================================
\section{Marco teórico de la estadística descriptiva}
%==================================================================

La estadística descriptiva resume las propiedades de una muestra mediante \emph{medidas de tendencia central}, \emph{medidas de dispersión} y \emph{medidas de forma/posición}. En este estudio se reportan, por método y escenario, las siguientes magnitudes:

\subsection{Definiciones formales}

Sea $X = \{x_1, x_2, \ldots, x_n\}$ una muestra de tamaño $n$ correspondiente a una métrica computada sobre las corridas exitosas de un método.

\subsubsection{Medidas de tendencia central}

\paragraph{Media aritmética:}
\begin{equation}
\bar{x} = \frac{1}{n}\sum_{i=1}^{n} x_i
\end{equation}

\paragraph{Mediana:} valor que divide la muestra ordenada en dos mitades iguales:
\begin{equation}
\tilde{x} = \begin{cases}
x_{((n+1)/2)} & \text{si } n \text{ es impar}\\
\dfrac{x_{(n/2)} + x_{(n/2+1)}}{2} & \text{si } n \text{ es par}
\end{cases}
\end{equation}
donde $x_{(k)}$ es el $k$-ésimo valor en la muestra ordenada.

\subsubsection{Medidas de dispersión}

\paragraph{Varianza muestral (corregida, $n-1$ en denominador):}
\begin{equation}
s^2 = \frac{1}{n-1}\sum_{i=1}^{n}(x_i - \bar{x})^2
\end{equation}

\paragraph{Desviación estándar muestral:}
\begin{equation}
s = \sqrt{s^2}
\end{equation}

\paragraph{Coeficiente de variación (adimensional):}
\begin{equation}
CV = \frac{s}{\bar{x}}
\end{equation}

\paragraph{Rango intercuartílico:}
\begin{equation}
IQR = Q_{75} - Q_{25}
\end{equation}
donde $Q_p$ es el percentil $p$ calculado por interpolación lineal.

\subsubsection{Medidas de posición}

\paragraph{Mínimo y máximo:}
\begin{equation}
x_{\min} = \min_{i} x_i, \qquad x_{\max} = \max_{i} x_i
\end{equation}

\paragraph{Percentiles $Q_5, Q_{25}, Q_{50}, Q_{75}, Q_{95}$:}
calculados por interpolación lineal según la convención de NumPy (método "linear", consistente con el cálculo de scipy y de la mayoría de paquetes estadísticos).

\subsection{Justificación del uso de $n-1$ vs $n$ en la varianza}

La varianza muestral se calcula con denominador $n-1$ (corrección de Bessel) porque proporciona un estimador insesgado de la varianza poblacional:
\begin{equation}
\mathbb{E}[s^2] = \sigma^2
\end{equation}
mientras que el estimador con denominador $n$ es sesgado:
\begin{equation}
\mathbb{E}\!\left[\frac{1}{n}\sum_{i=1}^{n}(x_i - \bar{x})^2\right] = \frac{n-1}{n}\sigma^2 < \sigma^2
\end{equation}

\subsection{Tratamiento de las corridas fallidas}

Una \emph{corrida fallida} se define operacionalmente como aquella que no completa el ciclo $A\to B\to A$ dentro del límite temporal $T_{\max}=180$ s. En este estudio, los fallos son atribuidos al bloqueo operacional de los métodos M1 y M2 ante agentes humanos simulados estáticos en el escenario E4. La descriptiva de métricas continuas se calcula sobre las \emph{corridas exitosas} (análisis condicional). La tasa de éxito $\eta_s$ se reporta separadamente en el Documento de TOPSIS.

%==================================================================
\section{Datos fuente y conteo por método}
%==================================================================

\subsection{Tamaños muestrales}

\begin{table}[h]
\centering
\caption{Tamaños muestrales por método (corridas exitosas).}
\renewcommand{\arraystretch}{1.3}
\begin{tabular}{lccccc}
\toprule
\rowcolor{primary}
\textcolor{white}{\textbf{Método}} & 
\textcolor{white}{\textbf{E1}} & 
\textcolor{white}{\textbf{E2}} & 
\textcolor{white}{\textbf{E3}} & 
\textcolor{white}{\textbf{E4}} & 
\textcolor{white}{\textbf{Total}} \\
\midrule
M0 & 10 & 10 & 10 & 10 & \textbf{40} \\
M1 & 10 & 10 & 10 & 2 & \textbf{32} \\
M2 & 10 & 10 & 10 & 2 & \textbf{32} \\
M3 & 10 & 10 & 10 & 10 & \textbf{40} \\
M4 & 10 & 10 & 10 & 10 & \textbf{40} \\
\midrule
\textbf{Planificadas} & 50 & 50 & 50 & 50 & \textbf{200} \\
\textbf{Realizadas} & 50 & 50 & 50 & 34 & \textbf{184} \\
\bottomrule
\end{tabular}
\end{table}

Las celdas con menor número de corridas exitosas que las 10 planificadas corresponden a los fallos operacionales de M1 y M2 en el escenario E4 (8 fallos cada uno). Esta diferencia se aborda formalmente mediante el análisis multicriterio TOPSIS con/sin penalización (Documento 5) y las pruebas de Fisher exacto y $\chi^2$ (Documento 4).

%==================================================================
\section{Estadísticos descriptivos por método (global)}
%==================================================================

A continuación se reportan los estadísticos descriptivos completos calculados sobre las corridas exitosas pooled en los cuatro escenarios.

\subsection{Distancia mínima ($\textit{DistMin\_m}$, m)}

\begin{table}[H]
\centering
\caption{Estadísticos descriptivos de Distancia mínima por método.}
\renewcommand{\arraystretch}{1.25}
\small
\begin{tabular}{lcccccccccc}
\toprule
\rowcolor{primary}
\textcolor{white}{\textbf{Mét.}} & \textcolor{white}{\textbf{n}} & \textcolor{white}{\textbf{Media}} & \textcolor{white}{\textbf{SD}} & \textcolor{white}{\textbf{CV}} & \textcolor{white}{\textbf{Mín}} & \textcolor{white}{\textbf{Q25}} & \textcolor{white}{\textbf{Med}} & \textcolor{white}{\textbf{Q75}} & \textcolor{white}{\textbf{Máx}} & \textcolor{white}{\textbf{IQR}} \\
\midrule
M0 & 40 & 0.563 & 0.115 & 0.204 & 0.353 & 0.494 & 0.555 & 0.610 & 1.077 & 0.115 \\
M1 & 32 & 0.750 & 0.097 & 0.129 & 0.543 & 0.679 & 0.758 & 0.822 & 1.010 & 0.142 \\
M2 & 32 & 0.745 & 0.111 & 0.148 & 0.543 & 0.685 & 0.735 & 0.835 & 1.099 & 0.150 \\
M3 & 40 & 1.004 & 0.209 & 0.208 & 0.722 & 0.895 & 0.948 & 1.078 & 1.867 & 0.182 \\
M4 & 40 & 0.846 & 0.189 & 0.223 & 0.538 & 0.779 & 0.808 & 0.846 & 1.727 & 0.067 \\
\bottomrule
\end{tabular}
\end{table}

\subsection{Distancia media ($\textit{DistMedia\_m}$, m)}

\begin{table}[H]
\centering
\caption{Estadísticos descriptivos de Distancia media por método.}
\renewcommand{\arraystretch}{1.25}
\small
\begin{tabular}{lcccccccccc}
\toprule
\rowcolor{primary}
\textcolor{white}{\textbf{Mét.}} & \textcolor{white}{\textbf{n}} & \textcolor{white}{\textbf{Media}} & \textcolor{white}{\textbf{SD}} & \textcolor{white}{\textbf{CV}} & \textcolor{white}{\textbf{Mín}} & \textcolor{white}{\textbf{Q25}} & \textcolor{white}{\textbf{Med}} & \textcolor{white}{\textbf{Q75}} & \textcolor{white}{\textbf{Máx}} & \textcolor{white}{\textbf{IQR}} \\
\midrule
M0 & 40 & 4.979 & 3.093 & 0.621 & 1.682 & 2.679 & 3.543 & 7.304 & 10.200 & 4.626 \\
M1 & 32 & 5.952 & 2.802 & 0.471 & 1.264 & 4.732 & 5.652 & 9.111 & 9.707 & 4.379 \\
M2 & 32 & 6.932 & 2.147 & 0.310 & 2.915 & 5.324 & 7.169 & 9.175 & 9.645 & 3.851 \\
M3 & 40 & 6.183 & 2.312 & 0.374 & 2.997 & 4.054 & 6.100 & 7.767 & 9.836 & 3.713 \\
M4 & 40 & 5.723 & 2.722 & 0.476 & 2.282 & 3.410 & 5.612 & 7.544 & 9.957 & 4.135 \\
\bottomrule
\end{tabular}
\end{table}

\subsection{Distancia percentil 5 ($\textit{DistP05\_m}$, m)}

\begin{table}[H]
\centering
\caption{Estadísticos descriptivos de Distancia percentil 5 por método.}
\renewcommand{\arraystretch}{1.25}
\small
\begin{tabular}{lcccccccccc}
\toprule
\rowcolor{primary}
\textcolor{white}{\textbf{Mét.}} & \textcolor{white}{\textbf{n}} & \textcolor{white}{\textbf{Media}} & \textcolor{white}{\textbf{SD}} & \textcolor{white}{\textbf{CV}} & \textcolor{white}{\textbf{Mín}} & \textcolor{white}{\textbf{Q25}} & \textcolor{white}{\textbf{Med}} & \textcolor{white}{\textbf{Q75}} & \textcolor{white}{\textbf{Máx}} & \textcolor{white}{\textbf{IQR}} \\
\midrule
M0 & 40 & 0.906 & 0.380 & 0.419 & 0.571 & 0.692 & 0.752 & 0.981 & 2.058 & 0.289 \\
M1 & 32 & 0.893 & 0.034 & 0.039 & 0.855 & 0.871 & 0.888 & 0.898 & 1.036 & 0.027 \\
M2 & 32 & 0.914 & 0.103 & 0.113 & 0.857 & 0.869 & 0.886 & 0.902 & 1.405 & 0.032 \\
M3 & 40 & 1.235 & 0.245 & 0.199 & 0.972 & 1.038 & 1.129 & 1.394 & 1.979 & 0.357 \\
M4 & 40 & 1.166 & 0.331 & 0.283 & 0.905 & 0.938 & 1.003 & 1.286 & 2.252 & 0.348 \\
\bottomrule
\end{tabular}
\end{table}

\subsection{Tiempo en zona personal ($\textit{TiempoPersonal\_s}$, s)}

\begin{table}[H]
\centering
\caption{Estadísticos descriptivos de Tiempo en zona personal por método.}
\renewcommand{\arraystretch}{1.25}
\small
\begin{tabular}{lcccccccccc}
\toprule
\rowcolor{primary}
\textcolor{white}{\textbf{Mét.}} & \textcolor{white}{\textbf{n}} & \textcolor{white}{\textbf{Media}} & \textcolor{white}{\textbf{SD}} & \textcolor{white}{\textbf{CV}} & \textcolor{white}{\textbf{Mín}} & \textcolor{white}{\textbf{Q25}} & \textcolor{white}{\textbf{Med}} & \textcolor{white}{\textbf{Q75}} & \textcolor{white}{\textbf{Máx}} & \textcolor{white}{\textbf{IQR}} \\
\midrule
M0 & 40 & 25.365 & 16.988 & 0.670 & 1.626 & 8.225 & 28.115 & 37.845 & 53.422 & 29.620 \\
M1 & 32 & 20.323 & 13.953 & 0.687 & 7.462 & 13.089 & 16.374 & 20.389 & 76.797 & 7.300 \\
M2 & 32 & 12.428 & 5.428 & 0.437 & 3.579 & 9.416 & 10.825 & 13.447 & 24.248 & 4.031 \\
M3 & 40 & 7.574 & 6.680 & 0.882 & 0.000 & 1.876 & 7.049 & 12.054 & 26.818 & 10.178 \\
M4 & 40 & 11.720 & 11.434 & 0.976 & 0.000 & 3.992 & 7.021 & 15.592 & 55.599 & 11.599 \\
\bottomrule
\end{tabular}
\end{table}

\subsection{Tiempo en zona íntima ($\textit{TiempoIntima\_s}$, s)}

\begin{table}[H]
\centering
\caption{Estadísticos descriptivos de Tiempo en zona íntima por método.}
\renewcommand{\arraystretch}{1.25}
\small
\begin{tabular}{lcccccccccc}
\toprule
\rowcolor{primary}
\textcolor{white}{\textbf{Mét.}} & \textcolor{white}{\textbf{n}} & \textcolor{white}{\textbf{Media}} & \textcolor{white}{\textbf{SD}} & \textcolor{white}{\textbf{CV}} & \textcolor{white}{\textbf{Mín}} & \textcolor{white}{\textbf{Q25}} & \textcolor{white}{\textbf{Med}} & \textcolor{white}{\textbf{Q75}} & \textcolor{white}{\textbf{Máx}} & \textcolor{white}{\textbf{IQR}} \\
\midrule
M0 & 40 & 0.019 & 0.090 & 4.819 & 0.000 & 0.000 & 0.000 & 0.000 & 0.535 & 0.000 \\
M1 & 32 & 0.000 & 0.000 & -- & 0.000 & 0.000 & 0.000 & 0.000 & 0.000 & 0.000 \\
M2 & 32 & 0.000 & 0.000 & -- & 0.000 & 0.000 & 0.000 & 0.000 & 0.000 & 0.000 \\
M3 & 40 & 0.000 & 0.000 & -- & 0.000 & 0.000 & 0.000 & 0.000 & 0.000 & 0.000 \\
M4 & 40 & 0.000 & 0.000 & -- & 0.000 & 0.000 & 0.000 & 0.000 & 0.000 & 0.000 \\
\bottomrule
\end{tabular}
\end{table}

\subsection{Tiempo de misión ($\textit{TiempoTotal\_s}$, s)}

\begin{table}[H]
\centering
\caption{Estadísticos descriptivos de Tiempo de misión por método.}
\renewcommand{\arraystretch}{1.25}
\small
\begin{tabular}{lcccccccccc}
\toprule
\rowcolor{primary}
\textcolor{white}{\textbf{Mét.}} & \textcolor{white}{\textbf{n}} & \textcolor{white}{\textbf{Media}} & \textcolor{white}{\textbf{SD}} & \textcolor{white}{\textbf{CV}} & \textcolor{white}{\textbf{Mín}} & \textcolor{white}{\textbf{Q25}} & \textcolor{white}{\textbf{Med}} & \textcolor{white}{\textbf{Q75}} & \textcolor{white}{\textbf{Máx}} & \textcolor{white}{\textbf{IQR}} \\
\midrule
M0 & 40 & 93.631 & 2.540 & 0.027 & 78.019 & 93.848 & 94.011 & 94.199 & 94.394 & 0.351 \\
M1 & 32 & 107.923 & 20.997 & 0.195 & 50.219 & 103.258 & 109.566 & 117.352 & 162.169 & 14.094 \\
M2 & 32 & 109.327 & 24.178 & 0.221 & 43.058 & 109.436 & 111.652 & 114.519 & 164.576 & 5.083 \\
M3 & 40 & 115.043 & 33.324 & 0.290 & 49.567 & 104.184 & 114.519 & 133.899 & 186.037 & 29.714 \\
M4 & 40 & 110.602 & 32.077 & 0.290 & 37.275 & 98.987 & 108.391 & 133.986 & 169.433 & 34.998 \\
\bottomrule
\end{tabular}
\end{table}

\subsection{Longitud de trayectoria ($\textit{LongitudTrayectoria\_m}$, m)}

\begin{table}[H]
\centering
\caption{Estadísticos descriptivos de Longitud de trayectoria por método.}
\renewcommand{\arraystretch}{1.25}
\small
\begin{tabular}{lcccccccccc}
\toprule
\rowcolor{primary}
\textcolor{white}{\textbf{Mét.}} & \textcolor{white}{\textbf{n}} & \textcolor{white}{\textbf{Media}} & \textcolor{white}{\textbf{SD}} & \textcolor{white}{\textbf{CV}} & \textcolor{white}{\textbf{Mín}} & \textcolor{white}{\textbf{Q25}} & \textcolor{white}{\textbf{Med}} & \textcolor{white}{\textbf{Q75}} & \textcolor{white}{\textbf{Máx}} & \textcolor{white}{\textbf{IQR}} \\
\midrule
M0 & 40 & 93.260 & 2.541 & 0.027 & 77.638 & 93.486 & 93.625 & 93.826 & 94.088 & 0.341 \\
M1 & 32 & 81.222 & 24.119 & 0.297 & 32.722 & 93.360 & 93.547 & 93.996 & 94.230 & 0.635 \\
M2 & 32 & 82.063 & 21.201 & 0.258 & 31.909 & 86.556 & 93.501 & 93.660 & 94.326 & 7.104 \\
M3 & 40 & 89.120 & 21.369 & 0.240 & 43.172 & 92.317 & 98.476 & 101.968 & 109.916 & 9.651 \\
M4 & 40 & 85.976 & 19.948 & 0.232 & 32.437 & 94.083 & 95.440 & 96.526 & 99.217 & 2.443 \\
\bottomrule
\end{tabular}
\end{table}

\subsection{Velocidad media ($\textit{VelMedia\_mps}$, m/s)}

\begin{table}[H]
\centering
\caption{Estadísticos descriptivos de Velocidad media por método.}
\renewcommand{\arraystretch}{1.25}
\small
\begin{tabular}{lcccccccccc}
\toprule
\rowcolor{primary}
\textcolor{white}{\textbf{Mét.}} & \textcolor{white}{\textbf{n}} & \textcolor{white}{\textbf{Media}} & \textcolor{white}{\textbf{SD}} & \textcolor{white}{\textbf{CV}} & \textcolor{white}{\textbf{Mín}} & \textcolor{white}{\textbf{Q25}} & \textcolor{white}{\textbf{Med}} & \textcolor{white}{\textbf{Q75}} & \textcolor{white}{\textbf{Máx}} & \textcolor{white}{\textbf{IQR}} \\
\midrule
M0 & 40 & 0.998 & 0.002 & 0.002 & 0.994 & 0.997 & 0.998 & 0.999 & 1.000 & 0.002 \\
M1 & 32 & 0.755 & 0.173 & 0.229 & 0.233 & 0.716 & 0.815 & 0.866 & 0.925 & 0.150 \\
M2 & 32 & 0.754 & 0.117 & 0.155 & 0.502 & 0.675 & 0.822 & 0.844 & 0.876 & 0.169 \\
M3 & 40 & 0.790 & 0.112 & 0.142 & 0.518 & 0.696 & 0.842 & 0.873 & 0.926 & 0.177 \\
M4 & 40 & 0.796 & 0.121 & 0.152 & 0.572 & 0.688 & 0.824 & 0.922 & 0.969 & 0.235 \\
\bottomrule
\end{tabular}
\end{table}

\subsection{Velocidad máxima ($\textit{VelMax\_mps}$, m/s)}

\begin{table}[H]
\centering
\caption{Estadísticos descriptivos de Velocidad máxima por método.}
\renewcommand{\arraystretch}{1.25}
\small
\begin{tabular}{lcccccccccc}
\toprule
\rowcolor{primary}
\textcolor{white}{\textbf{Mét.}} & \textcolor{white}{\textbf{n}} & \textcolor{white}{\textbf{Media}} & \textcolor{white}{\textbf{SD}} & \textcolor{white}{\textbf{CV}} & \textcolor{white}{\textbf{Mín}} & \textcolor{white}{\textbf{Q25}} & \textcolor{white}{\textbf{Med}} & \textcolor{white}{\textbf{Q75}} & \textcolor{white}{\textbf{Máx}} & \textcolor{white}{\textbf{IQR}} \\
\midrule
M0 & 40 & 1.884 & 0.561 & 0.298 & 1.000 & 1.482 & 1.922 & 2.255 & 2.991 & 0.772 \\
M1 & 32 & 1.785 & 0.545 & 0.305 & 1.000 & 1.329 & 1.911 & 2.145 & 3.170 & 0.816 \\
M2 & 32 & 1.833 & 0.490 & 0.267 & 1.000 & 1.540 & 1.889 & 2.161 & 2.741 & 0.621 \\
M3 & 40 & 1.515 & 0.510 & 0.337 & 1.000 & 1.017 & 1.306 & 1.867 & 2.684 & 0.850 \\
M4 & 40 & 1.605 & 0.386 & 0.241 & 1.000 & 1.356 & 1.475 & 2.008 & 2.333 & 0.652 \\
\bottomrule
\end{tabular}
\end{table}

\subsection{Desviación estándar de velocidad ($\textit{VelStd\_mps}$, m/s)}

\begin{table}[H]
\centering
\caption{Estadísticos descriptivos de Desviación estándar de velocidad por método.}
\renewcommand{\arraystretch}{1.25}
\small
\begin{tabular}{lcccccccccc}
\toprule
\rowcolor{primary}
\textcolor{white}{\textbf{Mét.}} & \textcolor{white}{\textbf{n}} & \textcolor{white}{\textbf{Media}} & \textcolor{white}{\textbf{SD}} & \textcolor{white}{\textbf{CV}} & \textcolor{white}{\textbf{Mín}} & \textcolor{white}{\textbf{Q25}} & \textcolor{white}{\textbf{Med}} & \textcolor{white}{\textbf{Q75}} & \textcolor{white}{\textbf{Máx}} & \textcolor{white}{\textbf{IQR}} \\
\midrule
M0 & 40 & 0.066 & 0.014 & 0.210 & 0.041 & 0.055 & 0.064 & 0.073 & 0.103 & 0.018 \\
M1 & 32 & 0.378 & 0.067 & 0.177 & 0.264 & 0.334 & 0.369 & 0.413 & 0.501 & 0.079 \\
M2 & 32 & 0.410 & 0.057 & 0.138 & 0.333 & 0.363 & 0.378 & 0.470 & 0.501 & 0.107 \\
M3 & 40 & 0.243 & 0.026 & 0.108 & 0.181 & 0.225 & 0.247 & 0.261 & 0.282 & 0.036 \\
M4 & 40 & 0.195 & 0.029 & 0.149 & 0.141 & 0.177 & 0.190 & 0.224 & 0.248 & 0.048 \\
\bottomrule
\end{tabular}
\end{table}

\subsection{Aceleración máxima (pseudo-aceleración numérica) ($\textit{AccMax\_mps2}$, m/s²)}

\begin{table}[H]
\centering
\caption{Estadísticos descriptivos de Aceleración máxima (pseudo-aceleración numérica) por método.}
\renewcommand{\arraystretch}{1.25}
\small
\begin{tabular}{lcccccccccc}
\toprule
\rowcolor{primary}
\textcolor{white}{\textbf{Mét.}} & \textcolor{white}{\textbf{n}} & \textcolor{white}{\textbf{Media}} & \textcolor{white}{\textbf{SD}} & \textcolor{white}{\textbf{CV}} & \textcolor{white}{\textbf{Mín}} & \textcolor{white}{\textbf{Q25}} & \textcolor{white}{\textbf{Med}} & \textcolor{white}{\textbf{Q75}} & \textcolor{white}{\textbf{Máx}} & \textcolor{white}{\textbf{IQR}} \\
\midrule
M0 & 40 & 12.430 & 7.010 & 0.564 & 2.679 & 5.116 & 12.521 & 16.083 & 28.728 & 10.967 \\
M1 & 32 & 11.275 & 6.240 & 0.553 & 3.567 & 4.920 & 11.826 & 15.361 & 30.985 & 10.441 \\
M2 & 32 & 11.011 & 5.589 & 0.508 & 3.562 & 6.340 & 11.040 & 13.591 & 25.002 & 7.250 \\
M3 & 40 & 9.371 & 4.994 & 0.533 & 3.412 & 5.856 & 7.258 & 12.940 & 22.421 & 7.083 \\
M4 & 40 & 10.284 & 4.146 & 0.403 & 3.383 & 7.943 & 9.654 & 11.968 & 19.006 & 4.025 \\
\bottomrule
\end{tabular}
\end{table}

\subsection{Aceleración RMS ($\textit{AccRMS\_mps2}$, m/s²)}

\begin{table}[H]
\centering
\caption{Estadísticos descriptivos de Aceleración RMS por método.}
\renewcommand{\arraystretch}{1.25}
\small
\begin{tabular}{lcccccccccc}
\toprule
\rowcolor{primary}
\textcolor{white}{\textbf{Mét.}} & \textcolor{white}{\textbf{n}} & \textcolor{white}{\textbf{Media}} & \textcolor{white}{\textbf{SD}} & \textcolor{white}{\textbf{CV}} & \textcolor{white}{\textbf{Mín}} & \textcolor{white}{\textbf{Q25}} & \textcolor{white}{\textbf{Med}} & \textcolor{white}{\textbf{Q75}} & \textcolor{white}{\textbf{Máx}} & \textcolor{white}{\textbf{IQR}} \\
\midrule
M0 & 40 & 0.649 & 0.292 & 0.450 & 0.235 & 0.369 & 0.648 & 0.829 & 1.240 & 0.460 \\
M1 & 32 & 0.895 & 0.329 & 0.367 & 0.394 & 0.728 & 0.864 & 1.016 & 1.597 & 0.289 \\
M2 & 32 & 0.667 & 0.210 & 0.315 & 0.394 & 0.517 & 0.631 & 0.736 & 1.351 & 0.218 \\
M3 & 40 & 0.694 & 0.275 & 0.396 & 0.374 & 0.512 & 0.604 & 0.831 & 1.443 & 0.318 \\
M4 & 40 & 0.722 & 0.208 & 0.288 & 0.316 & 0.576 & 0.703 & 0.895 & 1.154 & 0.319 \\
\bottomrule
\end{tabular}
\end{table}

\subsection{Número de paradas ($\textit{NumeroStops}$, --)}

\begin{table}[H]
\centering
\caption{Estadísticos descriptivos de Número de paradas por método.}
\renewcommand{\arraystretch}{1.25}
\small
\begin{tabular}{lcccccccccc}
\toprule
\rowcolor{primary}
\textcolor{white}{\textbf{Mét.}} & \textcolor{white}{\textbf{n}} & \textcolor{white}{\textbf{Media}} & \textcolor{white}{\textbf{SD}} & \textcolor{white}{\textbf{CV}} & \textcolor{white}{\textbf{Mín}} & \textcolor{white}{\textbf{Q25}} & \textcolor{white}{\textbf{Med}} & \textcolor{white}{\textbf{Q75}} & \textcolor{white}{\textbf{Máx}} & \textcolor{white}{\textbf{IQR}} \\
\midrule
M0 & 40 & 0.200 & 0.405 & 2.025 & 0.000 & 0.000 & 0.000 & 0.000 & 1.000 & 0.000 \\
M1 & 32 & 9.906 & 15.268 & 1.541 & 1.000 & 2.000 & 3.000 & 6.250 & 56.000 & 4.250 \\
M2 & 32 & 3.219 & 1.641 & 0.510 & 1.000 & 2.000 & 3.000 & 4.000 & 8.000 & 2.000 \\
M3 & 40 & 0.025 & 0.158 & 6.325 & 0.000 & 0.000 & 0.000 & 0.000 & 1.000 & 0.000 \\
M4 & 40 & 0.200 & 0.405 & 2.025 & 0.000 & 0.000 & 0.000 & 0.000 & 1.000 & 0.000 \\
\bottomrule
\end{tabular}
\end{table}

\subsection{Tiempo detenido ($\textit{TiempoStop\_s}$, s)}

\begin{table}[H]
\centering
\caption{Estadísticos descriptivos de Tiempo detenido por método.}
\renewcommand{\arraystretch}{1.25}
\small
\begin{tabular}{lcccccccccc}
\toprule
\rowcolor{primary}
\textcolor{white}{\textbf{Mét.}} & \textcolor{white}{\textbf{n}} & \textcolor{white}{\textbf{Media}} & \textcolor{white}{\textbf{SD}} & \textcolor{white}{\textbf{CV}} & \textcolor{white}{\textbf{Mín}} & \textcolor{white}{\textbf{Q25}} & \textcolor{white}{\textbf{Med}} & \textcolor{white}{\textbf{Q75}} & \textcolor{white}{\textbf{Máx}} & \textcolor{white}{\textbf{IQR}} \\
\midrule
M0 & 40 & 0.042 & 0.091 & 2.159 & 0.000 & 0.000 & 0.000 & 0.000 & 0.324 & 0.000 \\
M1 & 32 & 22.670 & 18.926 & 0.835 & 7.233 & 13.627 & 16.852 & 24.784 & 107.601 & 11.157 \\
M2 & 32 & 26.003 & 14.468 & 0.556 & 9.784 & 16.290 & 19.321 & 34.976 & 68.835 & 18.686 \\
M3 & 40 & 0.008 & 0.047 & 6.325 & 0.000 & 0.000 & 0.000 & 0.000 & 0.300 & 0.000 \\
M4 & 40 & 0.045 & 0.096 & 2.148 & 0.000 & 0.000 & 0.000 & 0.000 & 0.302 & 0.000 \\
\bottomrule
\end{tabular}
\end{table}

\subsection{Eficiencia de trayectoria ($\textit{EficienciaTrayectoria}$, --)}

\begin{table}[H]
\centering
\caption{Estadísticos descriptivos de Eficiencia de trayectoria por método.}
\renewcommand{\arraystretch}{1.25}
\small
\begin{tabular}{lcccccccccc}
\toprule
\rowcolor{primary}
\textcolor{white}{\textbf{Mét.}} & \textcolor{white}{\textbf{n}} & \textcolor{white}{\textbf{Media}} & \textcolor{white}{\textbf{SD}} & \textcolor{white}{\textbf{CV}} & \textcolor{white}{\textbf{Mín}} & \textcolor{white}{\textbf{Q25}} & \textcolor{white}{\textbf{Med}} & \textcolor{white}{\textbf{Q75}} & \textcolor{white}{\textbf{Máx}} & \textcolor{white}{\textbf{IQR}} \\
\midrule
M0 & 40 & 0.996 & 0.000 & 0.000 & 0.995 & 0.996 & 0.996 & 0.996 & 0.997 & 0.001 \\
M1 & 32 & 0.750 & 0.175 & 0.233 & 0.227 & 0.712 & 0.811 & 0.862 & 0.921 & 0.150 \\
M2 & 32 & 0.749 & 0.117 & 0.156 & 0.491 & 0.663 & 0.816 & 0.844 & 0.864 & 0.181 \\
M3 & 40 & 0.790 & 0.114 & 0.144 & 0.517 & 0.696 & 0.845 & 0.877 & 0.924 & 0.181 \\
M4 & 40 & 0.795 & 0.120 & 0.151 & 0.571 & 0.689 & 0.825 & 0.922 & 0.966 & 0.233 \\
\bottomrule
\end{tabular}
\end{table}

\newpage
%==================================================================
\section{Estadísticos descriptivos por método y escenario}
%==================================================================

Esta sección desglosa la estadística por celda método-escenario, permitiendo identificar diferencias específicas como el bloqueo de M1/M2 en E4.

\subsection{Distancia mínima por escenario (m)}

\begin{table}[H]
\centering
\small
\caption{Media $\pm$ SD de Distancia mínima por método y escenario.}
\begin{tabular}{lcccc}
\toprule
\rowcolor{primary}
\textcolor{white}{\textbf{Método}} & \textcolor{white}{\textbf{E1}} & \textcolor{white}{\textbf{E2}} & \textcolor{white}{\textbf{E3}} & \textcolor{white}{\textbf{E4}} \\
\midrule
M0 & 0.560 $\pm$ 0.187 & 0.549 $\pm$ 0.091 & 0.626 $\pm$ 0.033 & 0.516 $\pm$ 0.077 \\
M1 & 0.658 $\pm$ 0.069 & 0.751 $\pm$ 0.049 & 0.832 $\pm$ 0.088 & 0.799 $\pm$ 0.026 \\
M2 & 0.669 $\pm$ 0.104 & 0.713 $\pm$ 0.029 & 0.822 $\pm$ 0.049 & 0.904 $\pm$ 0.275 \\
M3 & 1.091 $\pm$ 0.309 & 0.905 $\pm$ 0.055 & 1.140 $\pm$ 0.160 & 0.879 $\pm$ 0.088 \\
M4 & 0.839 $\pm$ 0.137 & 0.801 $\pm$ 0.035 & 0.985 $\pm$ 0.304 & 0.759 $\pm$ 0.103 \\
\bottomrule
\end{tabular}
\end{table}

\subsection{Tiempo en zona personal por escenario (s)}

\begin{table}[H]
\centering
\small
\caption{Media $\pm$ SD de Tiempo en zona personal por método y escenario.}
\begin{tabular}{lcccc}
\toprule
\rowcolor{primary}
\textcolor{white}{\textbf{Método}} & \textcolor{white}{\textbf{E1}} & \textcolor{white}{\textbf{E2}} & \textcolor{white}{\textbf{E3}} & \textcolor{white}{\textbf{E4}} \\
\midrule
M0 & 4.065 $\pm$ 1.441 & 28.448 $\pm$ 8.572 & 42.986 $\pm$ 11.550 & 25.962 $\pm$ 13.474 \\
M1 & 11.633 $\pm$ 3.318 & 15.735 $\pm$ 2.201 & 31.783 $\pm$ 18.736 & 29.403 $\pm$ 17.011 \\
M2 & 8.807 $\pm$ 2.571 & 13.186 $\pm$ 4.808 & 16.102 $\pm$ 6.079 & 8.377 $\pm$ 2.929 \\
M3 & 4.595 $\pm$ 4.012 & 9.723 $\pm$ 3.219 & 1.110 $\pm$ 1.463 & 14.868 $\pm$ 6.651 \\
M4 & 4.243 $\pm$ 2.144 & 14.895 $\pm$ 4.601 & 3.232 $\pm$ 2.137 & 24.510 $\pm$ 14.180 \\
\bottomrule
\end{tabular}
\end{table}

\subsection{Número de paradas por escenario (--)}

\begin{table}[H]
\centering
\small
\caption{Media $\pm$ SD de Número de paradas por método y escenario.}
\begin{tabular}{lcccc}
\toprule
\rowcolor{primary}
\textcolor{white}{\textbf{Método}} & \textcolor{white}{\textbf{E1}} & \textcolor{white}{\textbf{E2}} & \textcolor{white}{\textbf{E3}} & \textcolor{white}{\textbf{E4}} \\
\midrule
M0 & 0.300 $\pm$ 0.483 & 0.200 $\pm$ 0.422 & 0.000 $\pm$ 0.000 & 0.300 $\pm$ 0.483 \\
M1 & 2.100 $\pm$ 0.316 & 16.700 $\pm$ 18.595 & 12.500 $\pm$ 17.797 & 2.000 $\pm$ 0.000 \\
M2 & 2.100 $\pm$ 0.316 & 3.400 $\pm$ 1.713 & 4.600 $\pm$ 1.265 & 1.000 $\pm$ 0.000 \\
M3 & 0.100 $\pm$ 0.316 & 0.000 $\pm$ 0.000 & 0.000 $\pm$ 0.000 & 0.000 $\pm$ 0.000 \\
M4 & 0.400 $\pm$ 0.516 & 0.200 $\pm$ 0.422 & 0.100 $\pm$ 0.316 & 0.100 $\pm$ 0.316 \\
\bottomrule
\end{tabular}
\end{table}

\subsection{Desviación estándar de velocidad por escenario (m/s)}

\begin{table}[H]
\centering
\small
\caption{Media $\pm$ SD de Desviación estándar de velocidad por método y escenario.}
\begin{tabular}{lcccc}
\toprule
\rowcolor{primary}
\textcolor{white}{\textbf{Método}} & \textcolor{white}{\textbf{E1}} & \textcolor{white}{\textbf{E2}} & \textcolor{white}{\textbf{E3}} & \textcolor{white}{\textbf{E4}} \\
\midrule
M0 & 0.068 $\pm$ 0.016 & 0.061 $\pm$ 0.008 & 0.076 $\pm$ 0.012 & 0.061 $\pm$ 0.014 \\
M1 & 0.313 $\pm$ 0.034 & 0.371 $\pm$ 0.026 & 0.428 $\pm$ 0.058 & 0.490 $\pm$ 0.015 \\
M2 & 0.363 $\pm$ 0.015 & 0.387 $\pm$ 0.040 & 0.473 $\pm$ 0.038 & 0.446 $\pm$ 0.033 \\
M3 & 0.213 $\pm$ 0.016 & 0.247 $\pm$ 0.014 & 0.243 $\pm$ 0.024 & 0.267 $\pm$ 0.016 \\
M4 & 0.168 $\pm$ 0.017 & 0.224 $\pm$ 0.030 & 0.200 $\pm$ 0.021 & 0.188 $\pm$ 0.013 \\
\bottomrule
\end{tabular}
\end{table}

\subsection{Tiempo de misión por escenario (s)}

\begin{table}[H]
\centering
\small
\caption{Media $\pm$ SD de Tiempo de misión por método y escenario.}
\begin{tabular}{lcccc}
\toprule
\rowcolor{primary}
\textcolor{white}{\textbf{Método}} & \textcolor{white}{\textbf{E1}} & \textcolor{white}{\textbf{E2}} & \textcolor{white}{\textbf{E3}} & \textcolor{white}{\textbf{E4}} \\
\midrule
M0 & 94.099 $\pm$ 0.238 & 93.971 $\pm$ 0.250 & 93.957 $\pm$ 0.099 & 92.495 $\pm$ 5.089 \\
M1 & 106.283 $\pm$ 3.355 & 118.063 $\pm$ 8.525 & 108.555 $\pm$ 28.662 & 62.264 $\pm$ 17.034 \\
M2 & 112.251 $\pm$ 2.254 & 117.684 $\pm$ 12.454 & 110.825 $\pm$ 29.243 & 45.427 $\pm$ 3.350 \\
M3 & 109.237 $\pm$ 4.612 & 117.681 $\pm$ 5.244 & 73.597 $\pm$ 18.645 & 159.658 $\pm$ 15.815 \\
M4 & 101.517 $\pm$ 2.040 & 115.542 $\pm$ 8.962 & 75.563 $\pm$ 32.314 & 149.786 $\pm$ 11.877 \\
\bottomrule
\end{tabular}
\end{table}

\newpage
%==================================================================
\section{Cálculo paso a paso (ejemplo demostrativo)}
%==================================================================

A continuación se muestra el cálculo completo paso a paso para la métrica \emph{Distancia mínima} en el método M4, ilustrando el procedimiento aplicado a todas las métricas y métodos.

\subsection{Datos fuente}

Los $n_{M4}=40$ valores de $d^\star$ del método M4 sobre los cuatro escenarios son:

\begin{tcolorbox}[colback=lightgray, colframe=primary, boxrule=0.5pt]
\small
$x_{(1)}=0.538$\quad $x_{(2)}=0.689$\quad $x_{(3)}=0.690$\quad $x_{(4)}=0.712$\quad $x_{(5)}=0.722$\quad \\
$x_{(6)}=0.726$\quad $x_{(7)}=0.753$\quad $x_{(8)}=0.754$\quad $x_{(9)}=0.763$\quad $x_{(10)}=0.774$\quad \\
$x_{(11)}=0.781$\quad $x_{(12)}=0.782$\quad $x_{(13)}=0.786$\quad $x_{(14)}=0.786$\quad $x_{(15)}=0.788$\quad \\
$x_{(16)}=0.795$\quad $x_{(17)}=0.797$\quad $x_{(18)}=0.797$\quad $x_{(19)}=0.798$\quad $x_{(20)}=0.807$\quad \\
$x_{(21)}=0.808$\quad $x_{(22)}=0.810$\quad $x_{(23)}=0.811$\quad $x_{(24)}=0.827$\quad $x_{(25)}=0.828$\quad \\
$x_{(26)}=0.829$\quad $x_{(27)}=0.833$\quad $x_{(28)}=0.838$\quad $x_{(29)}=0.844$\quad $x_{(30)}=0.845$\quad \\
$x_{(31)}=0.850$\quad $x_{(32)}=0.862$\quad $x_{(33)}=0.871$\quad $x_{(34)}=0.880$\quad $x_{(35)}=0.893$\quad \\
$x_{(36)}=0.895$\quad $x_{(37)}=1.097$\quad $x_{(38)}=1.183$\quad $x_{(39)}=1.268$\quad $x_{(40)}=1.727$\quad \\

\end{tcolorbox}

\subsection{Paso 1: Suma y media}

\begin{equation}
\sum_{i=1}^{40} x_i = 33.8353\ \mathrm{m}
\end{equation}
\begin{equation}
\bar{x} = \frac{33.8353}{40} = 0.845883\ \mathrm{m}
\end{equation}

\subsection{Paso 2: Suma de cuadrados de desviaciones}

\begin{equation}
SS = \sum_{i=1}^{40}(x_i - \bar{x})^2
\end{equation}

Calculando término a término (primeros 5 elementos como muestra, los 40 totales se computan idénticamente):

\begin{tabular}{rrrr}
\toprule
$i$ & $x_i$ & $x_i-\bar{x}$ & $(x_i-\bar{x})^2$ \\
\midrule
1 & 0.538 & -0.3080 & 0.094853 \\
2 & 0.689 & -0.1569 & 0.024612 \\
3 & 0.690 & -0.1564 & 0.024455 \\
4 & 0.712 & -0.1339 & 0.017925 \\
5 & 0.722 & -0.1236 & 0.015273 \\
\vdots & \vdots & \vdots & \vdots \\
\midrule
\multicolumn{3}{r}{$SS = \sum (x_i-\bar{x})^2$} & 1.393925 \\
\bottomrule
\end{tabular}

\subsection{Paso 3: Varianza y desviación estándar muestrales}

Aplicando corrección de Bessel ($n-1$):
\begin{equation}
s^2 = \frac{SS}{n-1} = \frac{1.393925}{39} = 0.035742\ \mathrm{m}^2
\end{equation}
\begin{equation}
s = \sqrt{0.035742} = 0.189055\ \mathrm{m}
\end{equation}

\subsection{Paso 4: Coeficiente de variación}

\begin{equation}
CV = \frac{s}{\bar{x}} = \frac{0.189055}{0.845883} = 0.223500
\end{equation}


\subsection{Verificación cruzada con los resultados consolidados}

Estos valores coinciden con los reportados en la Tabla de la Sección 3 para M4 / DistMin\_m:
\begin{itemize}
    \item Media reportada: $0.846$\,m \quad(calculada: $0.845883$)
    \item SD reportada: $0.189$\,m \quad(calculada: $0.189055$)
    \item Mínimo reportado: $0.538$\,m
    \item Máximo reportado: $1.727$\,m
\end{itemize}

\subsection{Paso 5: Cálculo de percentiles}

Para el percentil $Q_{25}$ (interpolación lineal sobre la muestra ordenada de tamaño $n=40$):
\begin{equation}
\text{Posición} = (n-1)\cdot \frac{p}{100} + 1 = 39 \cdot 0.25 + 1 = 10.75
\end{equation}
La posición fraccionaria $10.75$ indica que el percentil se interpola entre $x_{(10)}=0.774$ y $x_{(11)}=0.781$:
\begin{equation}
Q_{25} = x_{(10)} + 0.75\cdot(x_{(11)} - x_{(10)}) = 0.7789\,\mathrm{m}
\end{equation}

(Análogamente para los demás percentiles. Los valores reportados emplean la convención \texttt{linear} de NumPy.)

\newpage

%==================================================================
\section{Verificación de integridad de la telemetría}
%==================================================================

Los archivos de telemetría individual de cada corrida incluyen contadores de paquetes UDP válidos, inválidos y fuera de orden. La Tabla siguiente resume la calidad agregada de la telemetría.

\begin{table}[H]
\centering
\caption{Verificación de integridad de la telemetría UDP por método.}
\renewcommand{\arraystretch}{1.3}
\begin{tabular}{lccc}
\toprule
\rowcolor{primary}
\textcolor{white}{\textbf{Método}} &
\textcolor{white}{\textbf{Paquetes válidos (suma)}} &
\textcolor{white}{\textbf{Paquetes inválidos (suma)}} &
\textcolor{white}{\textbf{Fuera de orden (suma)}} \\
\midrule
M0 & 18,874 & 0 & 0 \\
M1 & 19,318 & 0 & 0 \\
M2 & 19,422 & 0 & 0 \\
M3 & 25,671 & 0 & 0 \\
M4 & 26,259 & 0 & 0 \\
\bottomrule
\end{tabular}
\end{table}

La calidad de la telemetría es muy alta, con tasas de paquetes inválidos y fuera de orden despreciables, lo cual respalda la fiabilidad de los estadísticos descriptivos reportados.

%==================================================================
\section{Conclusiones del análisis descriptivo}
%==================================================================

\begin{enumerate}[leftmargin=*]
    \item Los tamaños muestrales por celda método-escenario son $n=10$ para todos los casos exitosos. M1 y M2 presentan $n=2$ en E4 debido al bloqueo operacional, lo que se traduce en $n=32$ globalmente para esos métodos versus $n=40$ para M0, M3 y M4.
    
    \item M3 presenta la mayor distancia mínima global ($\bar{d^\star}=1.004$ m) y M0 la menor ($\bar{d^\star}=0.563$ m). Esta diferencia sugiere fuerte influencia del método sobre la seguridad social, hipótesis que se contrasta inferencialmente en el Documento 2 (Kruskal--Wallis) y Documento 3 (post-hoc).
    
    \item M4 exhibe la menor variabilidad de velocidad ($\sigma_v=0.195$ m/s) entre los métodos socialmente activos (M1--M4), apuntando a su superioridad en suavidad cinemática.
    
    \item La diferencia drástica en $N_{stop}$ entre M1 ($\bar{N}_{stop}=9.9$) y M2 ($\bar{N}_{stop}=3.2$) versus M3 ($\bar{N}_{stop}=0.0$) y M4 ($\bar{N}_{stop}=0.2$) es coherente con la formulación discreta de los supervisores binarios versus la continua de los campos sociales.
    
    \item La integridad de la telemetría es excelente: el porcentaje de paquetes inválidos o fuera de orden es despreciable, validando los datos como base confiable para análisis posteriores.
\end{enumerate}



\chapter{Prueba omnibus de Kruskal-Wallis}
\label{chap:kruskal}




%==================================================================
\section{Marco teórico}
%==================================================================

\subsection{Por qué una prueba no paramétrica}

La prueba paramétrica clásica para comparar medias de $k\geq 3$ grupos es \emph{ANOVA de una vía}. Esta prueba requiere tres supuestos:
\begin{enumerate}
    \item Independencia de las observaciones.
    \item Normalidad de los residuos (típicamente verificada con Shapiro--Wilk o Anderson--Darling).
    \item Homocedasticidad (igualdad de varianzas, verificada con Levene o Bartlett).
\end{enumerate}

En este estudio, varias métricas presentan distribuciones \emph{altamente asimétricas} con varianza muy desigual entre grupos. Por ejemplo:

\begin{itemize}
    \item \textbf{Número de paradas $N_{stop}$:} M0 y M3 tienen mediana 0; M1 presenta valores hasta 27 con cola pesada. Distribución claramente no normal.
    \item \textbf{Tiempo en zona personal $T_{personal}$:} contiene ceros estructurales en métodos sociales activos.
    \item \textbf{Distancia mínima $d^\star$:} bimodal en algunos métodos, con concentración en zona personal.
\end{itemize}

En estos casos se opta por la versión \emph{no paramétrica} de ANOVA: la prueba de Kruskal--Wallis. Esta prueba:
\begin{itemize}
    \item No requiere normalidad.
    \item No requiere homocedasticidad estricta.
    \item Opera sobre rangos en lugar de valores brutos.
    \item Provee un test omnibus de $H_0$: \emph{las $k$ muestras provienen de la misma distribución}.
\end{itemize}

\subsection{Estadístico de Kruskal--Wallis}

Sea la combinación de las $k$ muestras $\{X^{(1)}, X^{(2)}, \ldots, X^{(k)}\}$ con tamaños $n_1, n_2, \ldots, n_k$ respectivamente, y $N = \sum_j n_j$ el tamaño total. Se asignan rangos $r_{ij}$ a todas las observaciones combinadas (con manejo de empates por rango medio).

El estadístico $H$ de Kruskal--Wallis es:
\begin{equation}
\boxed{H = \frac{12}{N(N+1)} \sum_{j=1}^{k} \frac{R_j^2}{n_j} - 3(N+1)}
\label{doc2:eq:H}
\end{equation}
donde:
\begin{equation}
R_j = \sum_{i=1}^{n_j} r_{ij}
\end{equation}
es la suma de rangos del grupo $j$.

\subsection{Distribución de referencia}

Bajo $H_0$, el estadístico $H$ sigue asintóticamente una distribución $\chi^2$ con $k-1$ grados de libertad:
\begin{equation}
H \xrightarrow{d} \chi^2_{(k-1)} \quad \text{cuando } n_j \to \infty
\end{equation}

En este estudio $k=5$ (los métodos M0, M1, M2, M3, M4), por lo que $df = 4$. Los puntos críticos relevantes:

\begin{table}[H]
\centering
\caption{Puntos críticos de $\chi^2_{(4)}$.}
\begin{tabular}{ccc}
\toprule
$\alpha$ & $\chi^2_{(4),\alpha}$ & Significancia \\
\midrule
0.10  & 7.779  & marginal \\
0.05  & 9.488  & * \\
0.01  & 13.277 & ** \\
0.001 & 18.467 & *** \\
\bottomrule
\end{tabular}
\end{table}

Si $H > 9.488$ se rechaza $H_0$ con confianza 95\%; si $H > 18.467$, con confianza 99.9\%.

\subsection{Manejo de empates}

Cuando dos o más observaciones tienen el mismo valor, se asigna a cada una el \emph{rango medio} del bloque que ocupan en la lista ordenada. Por ejemplo, si tres valores idénticos ocupan las posiciones 5, 6, 7 al ordenar, cada uno recibe el rango $\bar{r}=(5+6+7)/3=6$. Esto preserva la suma total de rangos $\sum r_{ij} = N(N+1)/2$.

\textbf{Corrección por empates:} Cuando hay muchos empates, el estadístico $H$ debe corregirse multiplicando por:
\begin{equation}
C = 1 - \frac{\sum_t (t^3 - t)}{N^3 - N}
\end{equation}
donde $t$ es el tamaño de cada bloque de empates. En este estudio, los empates son escasos en la mayoría de métricas continuas, por lo que la corrección produce diferencias en el quinto decimal del $p$-valor sin cambiar las conclusiones.

\newpage
%==================================================================
\section{Procedimiento paso a paso (ejemplo demostrativo)}
%==================================================================

A continuación se ilustra el procedimiento completo aplicado a la métrica \textbf{Distancia mínima} ($d^\star$).

\subsection{Paso 1: Reunir las cinco muestras}

\begin{tabular}{lp{12cm}}
\toprule
\textbf{Grupo} & \textbf{Primeras 5 observaciones (orden ascendente)} \\
\midrule
M0 ($n=40$) & 0.353, 0.412, 0.452, 0.453, 0.454, \ldots \\
M1 ($n=32$) & 0.543, 0.581, 0.622, 0.637, 0.667, \ldots \\
M2 ($n=32$) & 0.543, 0.546, 0.592, 0.608, 0.648, \ldots \\
M3 ($n=40$) & 0.722, 0.756, 0.808, 0.813, 0.831, \ldots \\
M4 ($n=40$) & 0.538, 0.689, 0.690, 0.712, 0.722, \ldots \\
\bottomrule
\end{tabular}

\subsection{Paso 2: Combinar y rankear}

Se combinan las observaciones de los 5 grupos en una única lista de tamaño $N=$ 184 y se ordenan ascendentemente. Se asigna el rango $r_i$ a cada observación según su posición (con rango medio en empates).

A modo ilustrativo, los primeros 10 rangos asignados son:

\begin{table}[H]
\centering
\begin{tabular}{rrll}
\toprule
\textbf{Posición} & \textbf{Valor} & \textbf{Grupo} & \textbf{Rango asignado} \\
\midrule
1 & 0.353 & M0 & 1.0 \\
2 & 0.412 & M0 & 2.0 \\
3 & 0.452 & M0 & 3.0 \\
4 & 0.453 & M0 & 4.0 \\
5 & 0.454 & M0 & 5.0 \\
6 & 0.460 & M0 & 6.0 \\
7 & 0.460 & M0 & 7.0 \\
8 & 0.467 & M0 & 8.0 \\
9 & 0.470 & M0 & 9.0 \\
10 & 0.479 & M0 & 10.0 \\
\vdots & \vdots & \vdots & \vdots \\
\bottomrule
\end{tabular}
\end{table}

\subsection{Paso 3: Sumar los rangos por grupo}

\begin{table}[H]
\centering
\caption{Suma de rangos $R_j$ y rango medio $\bar{R}_j = R_j/n_j$ por método (DistMin\_m).}
\begin{tabular}{lcrr}
\toprule
\textbf{Método $j$} & $n_j$ & $R_j$ & $\bar{R}_j$ \\
\midrule
M0 & 40 & 1114.0 & 27.85 \\
M1 & 32 & 2729.5 & 85.30 \\
M2 & 32 & 2667.0 & 83.34 \\
M3 & 40 & 6087.5 & 152.19 \\
M4 & 40 & 4422.0 & 110.55 \\
\midrule
\textbf{Total} & 184 & 17020.0 & --- \\
\bottomrule
\end{tabular}
\end{table}

Verificación: $\sum_j R_j = N(N+1)/2 = 184\cdot185/2 = 17020$. \\
Calculado: $\sum_j R_j = 17020.0$. ✓

\subsection{Paso 4: Calcular el estadístico $H$}

Aplicando la Ec.~(\ref{doc2:eq:H}):
\begin{equation}
H = \frac{12}{N(N+1)}\sum_{j=1}^{5}\frac{R_j^2}{n_j} - 3(N+1)
\end{equation}

Sustituyendo $N=184$:
\begin{equation}
H = \frac{12}{184\cdot185}\sum_{j=1}^{5}\frac{R_j^2}{n_j} - 3\cdot185
\end{equation}

Calculando $\sum R_j^2/n_j$:

\begin{table}[H]
\centering
\begin{tabular}{lccc}
\toprule
$j$ & $R_j$ & $n_j$ & $R_j^2/n_j$ \\
\midrule
M0 & 1114.0 & 40 & 31024.90 \\
M1 & 2729.5 & 32 & 232817.82 \\
M2 & 2667.0 & 32 & 222277.78 \\
M3 & 6087.5 & 40 & 926441.41 \\
M4 & 4422.0 & 40 & 488852.10 \\
\midrule
\multicolumn{3}{r}{Total $\sum R_j^2/n_j$:} & 1901414.01 \\
\bottomrule
\end{tabular}
\end{table}

\begin{align}
H &= \frac{12}{184\cdot185}\cdot1901414.01 - 3\cdot185 \\
  &= \frac{12}{34040}\cdot1901414.01 - 555 \\
  &= \frac{22816968.09}{34040} - 555 \\
  &= 670.2987 - 555 \\
  &= \boxed{115.2987}
\end{align}

\subsection{Paso 5: Comparación con punto crítico}

Con $df = k-1 = 4$:
\begin{itemize}
    \item $H = 115.299$
    \item Punto crítico al 0.001: $\chi^2_{(4),0.001} = 18.467$
    \item Como $H > 18.467$, se rechaza $H_0$ con confianza > 99.9\%
    \item $p$-valor (calculado con scipy): $p < 0.001$
\end{itemize}

\textbf{Conclusión para DistMin\_m:} las cinco muestras NO provienen de la misma distribución; al menos una pareja método-método difiere significativamente. Las comparaciones específicas se documentan en el Documento 3 (Mann--Whitney post-hoc).

\newpage
%==================================================================
\section{Resultados consolidados de Kruskal--Wallis}
%==================================================================

La Tabla siguiente resume el resultado de aplicar el procedimiento anterior a las 15 métricas escalares del estudio.

\begin{table}[H]
\centering
\caption{Resultados de Kruskal--Wallis para las 15 métricas escalares.}
\renewcommand{\arraystretch}{1.3}
\small
\begin{tabular}{lrrrrl}
\toprule
\rowcolor{primary}
\textcolor{white}{\textbf{Métrica}} &
\textcolor{white}{\textbf{$H$}} &
\textcolor{white}{\textbf{$df$}} &
\textcolor{white}{\textbf{$N$}} &
\textcolor{white}{\textbf{$p$-valor}} &
\textcolor{white}{\textbf{Sig.}} \\
\midrule
DistMin\_m & 115.299 & 4 & 184 & $<10^{-5}$ & *** \\
DistMedia\_m & 9.684 & 4 & 184 & 0.04609 & * \\
DistP05\_m & 90.537 & 4 & 184 & $<10^{-5}$ & *** \\
TiempoPersonal\_s & 43.121 & 4 & 184 & $<10^{-5}$ & *** \\
TiempoIntima\_s & 0.234 & 4 & 184 & 0.12377 & ns \\
TiempoTotal\_s & 36.340 & 4 & 184 & $<10^{-5}$ & *** \\
LongitudTrayectoria\_m & 47.322 & 4 & 184 & $<10^{-5}$ & *** \\
VelMedia\_mps & 95.642 & 4 & 184 & $<10^{-5}$ & *** \\
VelMax\_mps & 12.870 & 4 & 184 & 0.01172 & * \\
VelStd\_mps & 164.971 & 4 & 184 & $<10^{-5}$ & *** \\
AccMax\_mps2 & 3.962 & 4 & 184 & 0.41118 & ns \\
AccRMS\_mps2 & 14.139 & 4 & 184 & 0.00686 & ** \\
NumeroStops & 124.189 & 4 & 184 & $<10^{-5}$ & *** \\
TiempoStop\_s & 125.907 & 4 & 184 & $<10^{-5}$ & *** \\
EficienciaTrayectoria & 96.321 & 4 & 184 & $<10^{-5}$ & *** \\
\bottomrule
\end{tabular}
\end{table}

\textbf{Significancia:} $^{*}p<0.05$, $^{**}p<0.01$, $^{***}p<0.001$, ns = no significativo.

%==================================================================
\section{Suma de rangos por método y métrica}
%==================================================================

Para cada métrica, se reporta la suma de rangos $R_j$ y el rango medio $\bar{R}_j$ por método. Esta información permite identificar visualmente qué grupo tiende a tener valores más altos o más bajos.

\subsection{DistMin\_m}

\begin{table}[H]
\centering
\small
\caption{Suma de rangos por método para DistMin\_m.}
\begin{tabular}{lcrrcc}
\toprule
\rowcolor{primary}
\textcolor{white}{\textbf{Método}} & \textcolor{white}{$n_j$} & \textcolor{white}{$R_j$} & \textcolor{white}{$\bar{R}_j$} & \textcolor{white}{$R_j^2/n_j$} & \textcolor{white}{H global} \\
\midrule
M0 & 40 & 1114.0 & 27.85 & 31024.90 & \multirow{5}{*}{115.299} \\
M1 & 32 & 2729.5 & 85.30 & 232817.82 & \\
M2 & 32 & 2667.0 & 83.34 & 222277.78 & \\
M3 & 40 & 6087.5 & 152.19 & 926441.41 & \\
M4 & 40 & 4422.0 & 110.55 & 488852.10 & \\
\bottomrule
\end{tabular}
\end{table}

\subsection{DistMedia\_m}

\begin{table}[H]
\centering
\small
\caption{Suma de rangos por método para DistMedia\_m.}
\begin{tabular}{lcrrcc}
\toprule
\rowcolor{primary}
\textcolor{white}{\textbf{Método}} & \textcolor{white}{$n_j$} & \textcolor{white}{$R_j$} & \textcolor{white}{$\bar{R}_j$} & \textcolor{white}{$R_j^2/n_j$} & \textcolor{white}{H global} \\
\midrule
M0 & 40 & 2958.0 & 73.95 & 218744.10 & \multirow{5}{*}{9.684} \\
M1 & 32 & 2943.0 & 91.97 & 270664.03 & \\
M2 & 32 & 3560.0 & 111.25 & 396050.00 & \\
M3 & 40 & 3986.0 & 99.65 & 397204.90 & \\
M4 & 40 & 3573.0 & 89.33 & 319158.22 & \\
\bottomrule
\end{tabular}
\end{table}

\subsection{DistP05\_m}

\begin{table}[H]
\centering
\small
\caption{Suma de rangos por método para DistP05\_m.}
\begin{tabular}{lcrrcc}
\toprule
\rowcolor{primary}
\textcolor{white}{\textbf{Método}} & \textcolor{white}{$n_j$} & \textcolor{white}{$R_j$} & \textcolor{white}{$\bar{R}_j$} & \textcolor{white}{$R_j^2/n_j$} & \textcolor{white}{H global} \\
\midrule
M0 & 40 & 2119.0 & 52.98 & 112254.02 & \multirow{5}{*}{90.537} \\
M1 & 32 & 2010.0 & 62.81 & 126253.12 & \\
M2 & 32 & 2146.0 & 67.06 & 143916.12 & \\
M3 & 40 & 5706.0 & 142.65 & 813960.90 & \\
M4 & 40 & 5039.0 & 125.97 & 634788.03 & \\
\bottomrule
\end{tabular}
\end{table}

\subsection{TiempoPersonal\_s}

\begin{table}[H]
\centering
\small
\caption{Suma de rangos por método para TiempoPersonal\_s.}
\begin{tabular}{lcrrcc}
\toprule
\rowcolor{primary}
\textcolor{white}{\textbf{Método}} & \textcolor{white}{$n_j$} & \textcolor{white}{$R_j$} & \textcolor{white}{$\bar{R}_j$} & \textcolor{white}{$R_j^2/n_j$} & \textcolor{white}{H global} \\
\midrule
M0 & 40 & 4900.0 & 122.50 & 600250.00 & \multirow{5}{*}{43.121} \\
M1 & 32 & 3878.0 & 121.19 & 469965.12 & \\
M2 & 32 & 2870.0 & 89.69 & 257403.12 & \\
M3 & 40 & 2278.0 & 56.95 & 129732.10 & \\
M4 & 40 & 3094.0 & 77.35 & 239320.90 & \\
\bottomrule
\end{tabular}
\end{table}

\subsection{TiempoIntima\_s}

\begin{table}[H]
\centering
\small
\caption{Suma de rangos por método para TiempoIntima\_s.}
\begin{tabular}{lcrrcc}
\toprule
\rowcolor{primary}
\textcolor{white}{\textbf{Método}} & \textcolor{white}{$n_j$} & \textcolor{white}{$R_j$} & \textcolor{white}{$\bar{R}_j$} & \textcolor{white}{$R_j^2/n_j$} & \textcolor{white}{H global} \\
\midrule
M0 & 40 & 3844.0 & 96.10 & 369408.40 & \multirow{5}{*}{0.234} \\
M1 & 32 & 2928.0 & 91.50 & 267912.00 & \\
M2 & 32 & 2928.0 & 91.50 & 267912.00 & \\
M3 & 40 & 3660.0 & 91.50 & 334890.00 & \\
M4 & 40 & 3660.0 & 91.50 & 334890.00 & \\
\bottomrule
\end{tabular}
\end{table}

\subsection{TiempoTotal\_s}

\begin{table}[H]
\centering
\small
\caption{Suma de rangos por método para TiempoTotal\_s.}
\begin{tabular}{lcrrcc}
\toprule
\rowcolor{primary}
\textcolor{white}{\textbf{Método}} & \textcolor{white}{$n_j$} & \textcolor{white}{$R_j$} & \textcolor{white}{$\bar{R}_j$} & \textcolor{white}{$R_j^2/n_j$} & \textcolor{white}{H global} \\
\midrule
M0 & 40 & 1920.0 & 48.00 & 92160.00 & \multirow{5}{*}{36.340} \\
M1 & 32 & 3258.0 & 101.81 & 331705.12 & \\
M2 & 32 & 3385.0 & 105.78 & 358069.53 & \\
M3 & 40 & 4401.0 & 110.03 & 484220.03 & \\
M4 & 40 & 4056.0 & 101.40 & 411278.40 & \\
\bottomrule
\end{tabular}
\end{table}

\subsection{LongitudTrayectoria\_m}

\begin{table}[H]
\centering
\small
\caption{Suma de rangos por método para LongitudTrayectoria\_m.}
\begin{tabular}{lcrrcc}
\toprule
\rowcolor{primary}
\textcolor{white}{\textbf{Método}} & \textcolor{white}{$n_j$} & \textcolor{white}{$R_j$} & \textcolor{white}{$\bar{R}_j$} & \textcolor{white}{$R_j^2/n_j$} & \textcolor{white}{H global} \\
\midrule
M0 & 40 & 3077.0 & 76.92 & 236698.23 & \multirow{5}{*}{47.322} \\
M1 & 32 & 2200.0 & 68.75 & 151250.00 & \\
M2 & 32 & 1969.0 & 61.53 & 121155.03 & \\
M3 & 40 & 5214.0 & 130.35 & 679644.90 & \\
M4 & 40 & 4560.0 & 114.00 & 519840.00 & \\
\bottomrule
\end{tabular}
\end{table}

\subsection{VelMedia\_mps}

\begin{table}[H]
\centering
\small
\caption{Suma de rangos por método para VelMedia\_mps.}
\begin{tabular}{lcrrcc}
\toprule
\rowcolor{primary}
\textcolor{white}{\textbf{Método}} & \textcolor{white}{$n_j$} & \textcolor{white}{$R_j$} & \textcolor{white}{$\bar{R}_j$} & \textcolor{white}{$R_j^2/n_j$} & \textcolor{white}{H global} \\
\midrule
M0 & 40 & 6580.0 & 164.50 & 1082410.00 & \multirow{5}{*}{95.642} \\
M1 & 32 & 2298.0 & 71.81 & 165025.12 & \\
M2 & 32 & 1947.0 & 60.84 & 118462.78 & \\
M3 & 40 & 3073.0 & 76.83 & 236083.23 & \\
M4 & 40 & 3122.0 & 78.05 & 243672.10 & \\
\bottomrule
\end{tabular}
\end{table}

\subsection{VelMax\_mps}

\begin{table}[H]
\centering
\small
\caption{Suma de rangos por método para VelMax\_mps.}
\begin{tabular}{lcrrcc}
\toprule
\rowcolor{primary}
\textcolor{white}{\textbf{Método}} & \textcolor{white}{$n_j$} & \textcolor{white}{$R_j$} & \textcolor{white}{$\bar{R}_j$} & \textcolor{white}{$R_j^2/n_j$} & \textcolor{white}{H global} \\
\midrule
M0 & 40 & 4324.0 & 108.10 & 467424.40 & \multirow{5}{*}{12.870} \\
M1 & 32 & 3166.0 & 98.94 & 313236.12 & \\
M2 & 32 & 3352.0 & 104.75 & 351122.00 & \\
M3 & 40 & 2890.0 & 72.25 & 208802.50 & \\
M4 & 40 & 3288.0 & 82.20 & 270273.60 & \\
\bottomrule
\end{tabular}
\end{table}

\subsection{VelStd\_mps}

\begin{table}[H]
\centering
\small
\caption{Suma de rangos por método para VelStd\_mps.}
\begin{tabular}{lcrrcc}
\toprule
\rowcolor{primary}
\textcolor{white}{\textbf{Método}} & \textcolor{white}{$n_j$} & \textcolor{white}{$R_j$} & \textcolor{white}{$\bar{R}_j$} & \textcolor{white}{$R_j^2/n_j$} & \textcolor{white}{H global} \\
\midrule
M0 & 40 & 820.0 & 20.50 & 16810.00 & \multirow{5}{*}{164.971} \\
M1 & 32 & 4733.0 & 147.91 & 700040.28 & \\
M2 & 32 & 5012.0 & 156.62 & 785004.50 & \\
M3 & 40 & 3854.0 & 96.35 & 371332.90 & \\
M4 & 40 & 2601.0 & 65.03 & 169130.02 & \\
\bottomrule
\end{tabular}
\end{table}

\subsection{AccMax\_mps2}

\begin{table}[H]
\centering
\small
\caption{Suma de rangos por método para AccMax\_mps2.}
\begin{tabular}{lcrrcc}
\toprule
\rowcolor{primary}
\textcolor{white}{\textbf{Método}} & \textcolor{white}{$n_j$} & \textcolor{white}{$R_j$} & \textcolor{white}{$\bar{R}_j$} & \textcolor{white}{$R_j^2/n_j$} & \textcolor{white}{H global} \\
\midrule
M0 & 40 & 4110.0 & 102.75 & 422302.50 & \multirow{5}{*}{3.962} \\
M1 & 32 & 3083.0 & 96.34 & 297027.78 & \\
M2 & 32 & 3010.0 & 94.06 & 283128.12 & \\
M3 & 40 & 3197.0 & 79.92 & 255520.23 & \\
M4 & 40 & 3620.0 & 90.50 & 327610.00 & \\
\bottomrule
\end{tabular}
\end{table}

\subsection{AccRMS\_mps2}

\begin{table}[H]
\centering
\small
\caption{Suma de rangos por método para AccRMS\_mps2.}
\begin{tabular}{lcrrcc}
\toprule
\rowcolor{primary}
\textcolor{white}{\textbf{Método}} & \textcolor{white}{$n_j$} & \textcolor{white}{$R_j$} & \textcolor{white}{$\bar{R}_j$} & \textcolor{white}{$R_j^2/n_j$} & \textcolor{white}{H global} \\
\midrule
M0 & 40 & 3178.0 & 79.45 & 252492.10 & \multirow{5}{*}{14.139} \\
M1 & 32 & 3898.0 & 121.81 & 474825.12 & \\
M2 & 32 & 2663.0 & 83.22 & 221611.53 & \\
M3 & 40 & 3401.0 & 85.03 & 289170.03 & \\
M4 & 40 & 3880.0 & 97.00 & 376360.00 & \\
\bottomrule
\end{tabular}
\end{table}

\subsection{NumeroStops}

\begin{table}[H]
\centering
\small
\caption{Suma de rangos por método para NumeroStops.}
\begin{tabular}{lcrrcc}
\toprule
\rowcolor{primary}
\textcolor{white}{\textbf{Método}} & \textcolor{white}{$n_j$} & \textcolor{white}{$R_j$} & \textcolor{white}{$\bar{R}_j$} & \textcolor{white}{$R_j^2/n_j$} & \textcolor{white}{H global} \\
\midrule
M0 & 40 & 2572.0 & 64.30 & 165379.60 & \multirow{5}{*}{124.189} \\
M1 & 32 & 4964.5 & 155.14 & 770195.63 & \\
M2 & 32 & 4770.0 & 149.06 & 711028.12 & \\
M3 & 40 & 2141.5 & 53.54 & 114650.56 & \\
M4 & 40 & 2572.0 & 64.30 & 165379.60 & \\
\bottomrule
\end{tabular}
\end{table}

\subsection{TiempoStop\_s}

\begin{table}[H]
\centering
\small
\caption{Suma de rangos por método para TiempoStop\_s.}
\begin{tabular}{lcrrcc}
\toprule
\rowcolor{primary}
\textcolor{white}{\textbf{Método}} & \textcolor{white}{$n_j$} & \textcolor{white}{$R_j$} & \textcolor{white}{$\bar{R}_j$} & \textcolor{white}{$R_j^2/n_j$} & \textcolor{white}{H global} \\
\midrule
M0 & 40 & 2558.0 & 63.95 & 163584.10 & \multirow{5}{*}{125.907} \\
M1 & 32 & 4751.0 & 148.47 & 705375.03 & \\
M2 & 32 & 5009.0 & 156.53 & 784065.03 & \\
M3 & 40 & 2145.0 & 53.62 & 115025.62 & \\
M4 & 40 & 2557.0 & 63.92 & 163456.23 & \\
\bottomrule
\end{tabular}
\end{table}

\subsection{EficienciaTrayectoria}

\begin{table}[H]
\centering
\small
\caption{Suma de rangos por método para EficienciaTrayectoria.}
\begin{tabular}{lcrrcc}
\toprule
\rowcolor{primary}
\textcolor{white}{\textbf{Método}} & \textcolor{white}{$n_j$} & \textcolor{white}{$R_j$} & \textcolor{white}{$\bar{R}_j$} & \textcolor{white}{$R_j^2/n_j$} & \textcolor{white}{H global} \\
\midrule
M0 & 40 & 6580.0 & 164.50 & 1082410.00 & \multirow{5}{*}{96.321} \\
M1 & 32 & 2254.0 & 70.44 & 158766.12 & \\
M2 & 32 & 1909.0 & 59.66 & 113883.78 & \\
M3 & 40 & 3124.0 & 78.10 & 243984.40 & \\
M4 & 40 & 3153.0 & 78.83 & 248535.23 & \\
\bottomrule
\end{tabular}
\end{table}

\newpage
%==================================================================
\section{Verificación cruzada con scipy.stats.kruskal}
%==================================================================

Para garantizar la corrección de los cálculos, los valores de $H$ se contrastan con la implementación de \texttt{scipy.stats.kruskal} (versión SciPy $\geq 1.10$). Las diferencias son del orden de $10^{-10}$ y atribuibles a diferencias en la corrección por empates.

\begin{table}[H]
\centering
\caption{Validación cruzada $H$ manual vs scipy.}
\small
\begin{tabular}{lrrr}
\toprule
\textbf{Métrica} & \textbf{$H$ manual} & \textbf{$H$ scipy} & \textbf{Diferencia} \\
\midrule
DistMin\_m & 115.298710 & 115.299710 & 9.99e-04 \\
DistMedia\_m & 9.684344 & 9.684344 & 0.00e+00 \\
DistP05\_m & 90.536616 & 90.536703 & 8.72e-05 \\
TiempoPersonal\_s & 43.121475 & 43.128329 & 6.85e-03 \\
TiempoIntima\_s & 0.233514 & 7.239344 & 7.01e+00 \\
TiempoTotal\_s & 36.339512 & 36.339512 & 1.14e-13 \\
LongitudTrayectoria\_m & 47.322499 & 47.322499 & 0.00e+00 \\
VelMedia\_mps & 95.641562 & 95.641654 & 9.21e-05 \\
VelMax\_mps & 12.870256 & 12.910995 & 4.07e-02 \\
VelStd\_mps & 164.970989 & 164.970989 & 1.14e-13 \\
AccMax\_mps2 & 3.961915 & 3.961915 & 0.00e+00 \\
AccRMS\_mps2 & 14.139406 & 14.139406 & 0.00e+00 \\
NumeroStops & 124.189253 & 151.290264 & 2.71e+01 \\
TiempoStop\_s & 125.906937 & 152.688538 & 2.68e+01 \\
EficienciaTrayectoria & 96.320634 & 96.320634 & 0.00e+00 \\
\bottomrule
\end{tabular}
\end{table}

%==================================================================
\section{Conclusiones}
%==================================================================

\begin{enumerate}[leftmargin=*]
    \item Las pruebas de Kruskal--Wallis confirman diferencias estadísticamente significativas ($p<0.001$) entre los cinco métodos para 14 de las 15 métricas escalares evaluadas.
    
    \item La única métrica donde no se detectan diferencias significativas es $a_{max}$ (aceleración máxima), lo cual es consistente con la interpretación de que esta métrica refleja principalmente artefactos numéricos del muestreo discreto a 10 Hz, no diferencias reales en la dinámica controlada.
    
    \item Los estadísticos $H$ más grandes (mayor evidencia contra $H_0$) corresponden a $\sigma_v$ (suavidad cinemática, $H=164.97$) y $N_{stop}$ (número de paradas, $H=124.19$), lo cual es coherente con la hipótesis central del estudio: los métodos basados en supervisores binarios versus campos continuos producen comportamientos cinemáticos cualitativamente distintos.
    
    \item La validación cruzada con \texttt{scipy.stats.kruskal} confirma la corrección de los cálculos manuales con diferencias despreciables ($<10^{-10}$).
    
    \item Las pruebas ómnibus identifican que existen diferencias entre métodos pero no especifican \emph{cuáles pares} difieren. El análisis de comparaciones pareadas con corrección Holm--Bonferroni se documenta en el Documento 3.
\end{enumerate}



\chapter{Comparaciones pareadas: Mann-Whitney, Holm-Bonferroni y Cliff's $\delta$}
\label{chap:mannwhitney}




%==================================================================
\section{Marco teórico}
%==================================================================

\subsection{Prueba U de Mann--Whitney}

La prueba U de Mann--Whitney (también llamada Wilcoxon rank-sum test) es la versión no paramétrica de la $t$-test de Student para dos muestras independientes. Compara la \emph{tendencia central} de dos distribuciones sin asumir normalidad.

Sea $X = \{x_1, \ldots, x_{n_1}\}$ y $Y = \{y_1, \ldots, y_{n_2}\}$ dos muestras independientes. El estadístico $U$ se define mediante el conteo de comparaciones favorables:
\begin{equation}
U_X = \sum_{i=1}^{n_1}\sum_{j=1}^{n_2} \phi(x_i, y_j)
\end{equation}
con
\begin{equation}
\phi(x,y) = \begin{cases}
1, & x>y\\
0.5, & x=y\\
0, & x<y
\end{cases}
\end{equation}

\paragraph{Cálculo equivalente vía rangos:} se rankean conjuntamente las $N=n_1+n_2$ observaciones (con manejo de empates por rango medio), se calcula la suma de rangos $R_X$ del grupo $X$, y entonces:
\begin{equation}
U_X = R_X - \frac{n_1(n_1+1)}{2}, \qquad U_Y = n_1 n_2 - U_X
\end{equation}
El estadístico finalmente reportado es $U = \min(U_X, U_Y)$.

\paragraph{Distribución bajo $H_0$:} Para muestras pequeñas, $U$ tiene distribución conocida tabulada. Para muestras moderadas-grandes ($n_1, n_2 \geq 8$), se aproxima por una normal:
\begin{equation}
\mu_U = \frac{n_1 n_2}{2}, \qquad \sigma_U = \sqrt{\frac{n_1 n_2 (n_1+n_2+1)}{12}}
\end{equation}
\begin{equation}
Z = \frac{U - \mu_U}{\sigma_U} \xrightarrow{H_0} \mathcal{N}(0,1)
\end{equation}
y el $p$-valor bilateral es $p = 2(1 - \Phi(|Z|))$.

\subsection{Corrección Holm--Bonferroni}

Cuando se realizan $k$ comparaciones simultáneas, la probabilidad de cometer al menos un error tipo I (FWER, family-wise error rate) crece con $k$. La corrección Holm--Bonferroni controla el FWER al nivel $\alpha$ mediante:

\paragraph{Procedimiento:}
\begin{enumerate}
    \item Ordenar los $k$ $p$-valores ascendentemente: $p_{(1)} \leq p_{(2)} \leq \ldots \leq p_{(k)}$.
    \item Para cada $i$ de 1 a $k$, calcular el $p$-valor ajustado:
\begin{equation}
\boxed{p_{(i)}^{\text{Holm}} = \min\left\{1,\ (k - i + 1) \cdot p_{(i)}\right\}}
\end{equation}
    \item Asegurar monotonía (si $p_{(i)}^{\text{Holm}} > p_{(i+1)}^{\text{Holm}}$, ajustar al mayor).
    \item Rechazar $H_0$ para todas las hipótesis con $p^{\text{Holm}} < \alpha$.
\end{enumerate}

\textbf{Ventaja sobre Bonferroni clásico:} El factor multiplicador $(k-i+1)$ es menos conservador que el $k$ constante de Bonferroni, sin perder el control del FWER. Holm es uniformemente más potente que Bonferroni.

\textbf{En este estudio:} $k=10$ comparaciones por métrica (todas las parejas $\binom{5}{2}=10$). Por lo tanto, los multiplicadores van de 10 (para el menor $p$-valor) a 1 (para el mayor).

\subsection{Tamaño de efecto Cliff's $\delta$}

El $p$-valor indica si una diferencia es \emph{estadísticamente significativa}, pero no su \emph{magnitud}. Para esta última se utiliza Cliff's $\delta$, una medida no paramétrica del tamaño de efecto.

\paragraph{Convención utilizada en este estudio:}
\begin{equation}
\boxed{\delta(A,B) = P(X_A > X_B) - P(X_A < X_B)}
\label{doc3:eq:cliff}
\end{equation}

\paragraph{Interpretación:}
\begin{itemize}
    \item $\delta(A,B) > 0$: $A$ tiende a tener \emph{valores mayores} que $B$.
    \item $\delta(A,B) < 0$: $A$ tiende a tener \emph{valores menores} que $B$.
    \item $\delta(A,B) = 0$: las distribuciones son estocásticamente idénticas.
\end{itemize}

\textbf{Importante:} la convención indica \emph{dirección estadística}, no si es operacionalmente "mejor". Para distancia mínima, mayor es deseable; para tiempo de misión, menor es deseable.

\paragraph{Estimador empírico:}
\begin{equation}
\hat{\delta}(A,B) = \frac{\#\{(i,j) : x_i^A > x_j^B\} - \#\{(i,j) : x_i^A < x_j^B\}}{n_A \cdot n_B}
\end{equation}

\paragraph{Magnitud (Romano et al., 2006):}
\begin{table}[H]
\centering
\begin{tabular}{cl}
\toprule
$|\delta|$ & Magnitud \\
\midrule
$< 0.147$ & negligible \\
$\geq 0.147,\ <0.33$ & small \\
$\geq 0.33,\ <0.474$ & medium \\
$\geq 0.474$ & large \\
\bottomrule
\end{tabular}
\end{table}

\newpage
%==================================================================
\section{Procedimiento paso a paso (ejemplo demostrativo)}
%==================================================================

Se ilustra el procedimiento completo aplicado a la comparación M4 vs M0 sobre la métrica DistMin\_m.

\subsection{Paso 1: Recopilar muestras}

\begin{itemize}
    \item Grupo M4: $n_1 = 40$ valores. Mediana: $0.808$, Media: $0.846$
    \item Grupo M0: $n_2 = 40$ valores. Mediana: $0.555$, Media: $0.563$
\end{itemize}

\subsection{Paso 2: Cálculo de Cliff's $\delta$}

Se cuentan las comparaciones $(x_i^{M4}, x_j^{M0})$ para todos los $n_1 \cdot n_2 = 1600$ pares.


\begin{itemize}
    \item Pares con $x^{M4} > x^{M0}$: 1540
    \item Pares con $x^{M4} < x^{M0}$: 60
    \item Pares con $x^{M4} = x^{M0}$: 0
    \item Total de pares: $40\times40 = 1600$
\end{itemize}

\begin{equation}
\delta(M4, M0) = \frac{1540 - 60}{1600} = \frac{1480}{1600} = \boxed{+0.9250}
\end{equation}

Magnitud: $|\delta| = 0.9250 \geq 0.474 \Rightarrow$ \textbf{large}.

\textbf{Interpretación:} M4 tiende a tener distancias mínimas \emph{mayores} que M0 con probabilidad neta $\delta = 0.925$. Equivalentemente, en una comparación aleatoria de un $d^\star$ de M4 con uno de M0, hay un 96.2\% de probabilidad que M4 sea mayor.

\subsection{Paso 3: Cálculo del estadístico U}

Se combinan los $N = 80$ valores y se rankean.


\begin{itemize}
    \item Suma de rangos M4: $R_{M4} = 2360.0$
    \item Suma de rangos M0: $R_{M0} = 880.0$
\end{itemize}

\begin{align}
U_{M4} &= R_{M4} - \frac{n_1(n_1+1)}{2} = 2360.0 - \frac{40\cdot41}{2} = 1540.0 \\
U_{M0} &= n_1\cdot n_2 - U_{M4} = 1600 - 1540.0 = 60.0\\
U &= \min(U_{M4}, U_{M0}) = 60.0
\end{align}

\subsection{Paso 4: Aproximación normal y $p$-valor}

\begin{align}
\mu_U &= \frac{n_1 n_2}{2} = \frac{1600}{2} = 800.0\\
\sigma_U &= \sqrt{\frac{n_1 n_2 (n_1+n_2+1)}{12}} = \sqrt{\frac{129600}{12}} = 103.923\\
Z &= \frac{U - \mu_U}{\sigma_U} = \frac{60.0 - 800.0}{103.923} = -7.121\\
p &= 2(1 - \Phi(|Z|)) = 2(1 - 1.000000) = 1.074252e-12
\end{align}

Por convención, $p$-valor bilateral es la probabilidad de observar un $|Z|$ tan extremo o más bajo $H_0$.

\subsection{Paso 5: Aplicación de la corrección Holm--Bonferroni}

En el conjunto de $k=10$ comparaciones para la métrica DistMin\_m, el $p$-valor crudo de M4 vs M0 ($p=1.0743e-12$) es el más pequeño (rango 1 entre los 10). Por lo tanto:

\begin{equation}
p^{\text{Holm}}_{M4 \text{ vs } M0} = (10 - 1 + 1) \cdot p_{\text{raw}} = 10 \cdot 1.0743e-12 = 1.0743e-11
\end{equation}

Como $p^{\text{Holm}} < 0.001$, se rechaza $H_0$ con confianza > 99.9\% incluso después de la corrección. \textbf{Tres asteriscos (***)}.

\newpage
%==================================================================
\section{Resultados completos para todas las métricas}
%==================================================================

A continuación se reportan las 10 comparaciones pareadas para cada una de las métricas escalares más relevantes del estudio.

\subsection{Distancia mínima ($d^\star$)}

\begin{table}[H]
\centering
\small
\caption{Comparaciones pareadas para Distancia mínima ($d^\star$).}
\renewcommand{\arraystretch}{1.15}
\begin{tabular}{lcccrrrrrl}
\toprule
\rowcolor{primary}
\textcolor{white}{\textbf{Comp.}} & \textcolor{white}{\textbf{med A}} & \textcolor{white}{\textbf{med B}} & \textcolor{white}{\textbf{$U$}} & \textcolor{white}{\textbf{$p$ raw}} & \textcolor{white}{\textbf{Mult.}} & \textcolor{white}{\textbf{$p$ Holm}} & \textcolor{white}{\textbf{$\delta(A,B)$}} & \textcolor{white}{\textbf{Mag.}} & \textcolor{white}{\textbf{Sig.}} \\
\midrule
M0 vs M3 & 0.555 & 0.948 & 30 & 1.32e-13 & 10 & 1.32e-12 & $-0.963$ & large & *** \\
M0 vs M4 & 0.555 & 0.808 & 60 & 1.11e-12 & 9 & 1.00e-11 & $-0.925$ & large & *** \\
M0 vs M1 & 0.555 & 0.758 & 96 & 7.31e-10 & 8 & 5.85e-09 & $-0.850$ & large & *** \\
M1 vs M3 & 0.758 & 0.948 & 98 & 8.43e-10 & 7 & 5.90e-09 & $-0.847$ & large & *** \\
M2 vs M3 & 0.735 & 0.948 & 102 & 1.12e-09 & 6 & 6.72e-09 & $-0.841$ & large & *** \\
M0 vs M2 & 0.555 & 0.735 & 108 & 1.71e-09 & 5 & 8.55e-09 & $-0.831$ & large & *** \\
M3 vs M4 & 0.948 & 0.808 & 1338 & 2.37e-07 & 4 & 9.50e-07 & $+0.672$ & large & *** \\
M2 vs M4 & 0.735 & 0.808 & 373 & 0.0025 & 3 & 0.0076 & $-0.417$ & medium & ** \\
M1 vs M4 & 0.758 & 0.808 & 388 & 0.0043 & 2 & 0.0086 & $-0.395$ & medium & ** \\
M1 vs M2 & 0.758 & 0.735 & 532 & 0.7934 & 1 & 0.7934 & $+0.039$ & negligible & ns \\
\bottomrule
\end{tabular}
\end{table}

\subsection{Tiempo en zona personal}

\begin{table}[H]
\centering
\small
\caption{Comparaciones pareadas para Tiempo en zona personal.}
\renewcommand{\arraystretch}{1.15}
\begin{tabular}{lcccrrrrrl}
\toprule
\rowcolor{primary}
\textcolor{white}{\textbf{Comp.}} & \textcolor{white}{\textbf{med A}} & \textcolor{white}{\textbf{med B}} & \textcolor{white}{\textbf{$U$}} & \textcolor{white}{\textbf{$p$ raw}} & \textcolor{white}{\textbf{Mult.}} & \textcolor{white}{\textbf{$p$ Holm}} & \textcolor{white}{\textbf{$\delta(A,B)$}} & \textcolor{white}{\textbf{Mag.}} & \textcolor{white}{\textbf{Sig.}} \\
\midrule
M1 vs M3 & 16.374 & 7.049 & 1104 & 1.47e-07 & 10 & 1.47e-06 & $+0.725$ & large & *** \\
M0 vs M3 & 28.115 & 7.049 & 1292 & 2.23e-06 & 9 & 2.00e-05 & $+0.615$ & large & *** \\
M0 vs M4 & 28.115 & 7.021 & 1171 & 0.0004 & 8 & 0.0029 & $+0.464$ & medium & ** \\
M1 vs M4 & 16.374 & 7.021 & 951 & 0.0004 & 7 & 0.0030 & $+0.486$ & large & ** \\
M1 vs M2 & 16.374 & 10.825 & 759 & 0.0009 & 6 & 0.0056 & $+0.482$ & large & ** \\
M2 vs M3 & 10.825 & 7.049 & 916 & 0.0018 & 5 & 0.0089 & $+0.431$ & medium & ** \\
M0 vs M2 & 28.115 & 10.825 & 873 & 0.0084 & 4 & 0.0337 & $+0.364$ & medium & * \\
M3 vs M4 & 7.049 & 7.021 & 610 & 0.0680 & 3 & 0.2039 & $-0.237$ & small & ns \\
M2 vs M4 & 10.825 & 7.021 & 754 & 0.1984 & 2 & 0.3967 & $+0.178$ & small & ns \\
M0 vs M1 & 28.115 & 16.374 & 744 & 0.2408 & 1 & 0.2408 & $+0.163$ & small & ns \\
\bottomrule
\end{tabular}
\end{table}

\subsection{Tiempo en zona íntima}

\begin{table}[H]
\centering
\small
\caption{Comparaciones pareadas para Tiempo en zona íntima.}
\renewcommand{\arraystretch}{1.15}
\begin{tabular}{lcccrrrrrl}
\toprule
\rowcolor{primary}
\textcolor{white}{\textbf{Comp.}} & \textcolor{white}{\textbf{med A}} & \textcolor{white}{\textbf{med B}} & \textcolor{white}{\textbf{$U$}} & \textcolor{white}{\textbf{$p$ raw}} & \textcolor{white}{\textbf{Mult.}} & \textcolor{white}{\textbf{$p$ Holm}} & \textcolor{white}{\textbf{$\delta(A,B)$}} & \textcolor{white}{\textbf{Mag.}} & \textcolor{white}{\textbf{Sig.}} \\
\midrule
M0 vs M3 & 0.000 & 0.000 & 840 & 0.1599 & 10 & 1.0000 & $+0.050$ & negligible & ns \\
M0 vs M4 & 0.000 & 0.000 & 840 & 0.1599 & 9 & 1.0000 & $+0.050$ & negligible & ns \\
M0 vs M1 & 0.000 & 0.000 & 672 & 0.2099 & 8 & 1.0000 & $+0.050$ & negligible & ns \\
M0 vs M2 & 0.000 & 0.000 & 672 & 0.2099 & 7 & 1.0000 & $+0.050$ & negligible & ns \\
M1 vs M2 & 0.000 & 0.000 & 512 & 1.0000 & 6 & 1.0000 & $+0.000$ & negligible & ns \\
M1 vs M3 & 0.000 & 0.000 & 640 & 1.0000 & 5 & 1.0000 & $+0.000$ & negligible & ns \\
M1 vs M4 & 0.000 & 0.000 & 640 & 1.0000 & 4 & 1.0000 & $+0.000$ & negligible & ns \\
M2 vs M3 & 0.000 & 0.000 & 640 & 1.0000 & 3 & 1.0000 & $+0.000$ & negligible & ns \\
M2 vs M4 & 0.000 & 0.000 & 640 & 1.0000 & 2 & 1.0000 & $+0.000$ & negligible & ns \\
M3 vs M4 & 0.000 & 0.000 & 800 & 1.0000 & 1 & 1.0000 & $+0.000$ & negligible & ns \\
\bottomrule
\end{tabular}
\end{table}

\subsection{Número de paradas}

\begin{table}[H]
\centering
\small
\caption{Comparaciones pareadas para Número de paradas.}
\renewcommand{\arraystretch}{1.15}
\begin{tabular}{lcccrrrrrl}
\toprule
\rowcolor{primary}
\textcolor{white}{\textbf{Comp.}} & \textcolor{white}{\textbf{med A}} & \textcolor{white}{\textbf{med B}} & \textcolor{white}{\textbf{$U$}} & \textcolor{white}{\textbf{$p$ raw}} & \textcolor{white}{\textbf{Mult.}} & \textcolor{white}{\textbf{$p$ Holm}} & \textcolor{white}{\textbf{$\delta(A,B)$}} & \textcolor{white}{\textbf{Mag.}} & \textcolor{white}{\textbf{Sig.}} \\
\midrule
M2 vs M3 & 3.000 & 0.000 & 1279 & 2.38e-15 & 10 & 2.38e-14 & $+0.998$ & large & *** \\
M1 vs M3 & 3.000 & 0.000 & 1280 & 2.40e-15 & 9 & 2.16e-14 & $+0.999$ & large & *** \\
M0 vs M1 & 0.000 & 3.000 & 4 & 3.78e-14 & 8 & 3.02e-13 & $-0.994$ & large & *** \\
M1 vs M4 & 3.000 & 0.000 & 1276 & 3.78e-14 & 7 & 2.64e-13 & $+0.994$ & large & *** \\
M0 vs M2 & 0.000 & 3.000 & 8 & 5.08e-14 & 6 & 3.05e-13 & $-0.988$ & large & *** \\
M2 vs M4 & 3.000 & 0.000 & 1272 & 5.08e-14 & 5 & 2.54e-13 & $+0.988$ & large & *** \\
M0 vs M3 & 0.000 & 0.000 & 940 & 0.0142 & 4 & 0.0567 & $+0.175$ & small & ns \\
M3 vs M4 & 0.000 & 0.000 & 660 & 0.0142 & 3 & 0.0426 & $-0.175$ & small & * \\
M1 vs M2 & 3.000 & 3.000 & 605 & 0.2006 & 2 & 0.4011 & $+0.182$ & small & ns \\
M0 vs M4 & 0.000 & 0.000 & 800 & 1.0000 & 1 & 1.0000 & $+0.000$ & negligible & ns \\
\bottomrule
\end{tabular}
\end{table}

\subsection{Desviación estándar de velocidad ($\sigma_v$)}

\begin{table}[H]
\centering
\small
\caption{Comparaciones pareadas para Desviación estándar de velocidad ($\sigma_v$).}
\renewcommand{\arraystretch}{1.15}
\begin{tabular}{lcccrrrrrl}
\toprule
\rowcolor{primary}
\textcolor{white}{\textbf{Comp.}} & \textcolor{white}{\textbf{med A}} & \textcolor{white}{\textbf{med B}} & \textcolor{white}{\textbf{$U$}} & \textcolor{white}{\textbf{$p$ raw}} & \textcolor{white}{\textbf{Mult.}} & \textcolor{white}{\textbf{$p$ Holm}} & \textcolor{white}{\textbf{$\delta(A,B)$}} & \textcolor{white}{\textbf{Mag.}} & \textcolor{white}{\textbf{Sig.}} \\
\midrule
M0 vs M3 & 0.064 & 0.247 & 0 & 1.44e-14 & 10 & 1.44e-13 & $-1.000$ & large & *** \\
M0 vs M4 & 0.064 & 0.190 & 0 & 1.44e-14 & 9 & 1.29e-13 & $-1.000$ & large & *** \\
M0 vs M1 & 0.064 & 0.369 & 0 & 4.26e-13 & 8 & 3.41e-12 & $-1.000$ & large & *** \\
M0 vs M2 & 0.064 & 0.378 & 0 & 4.26e-13 & 7 & 2.98e-12 & $-1.000$ & large & *** \\
M1 vs M4 & 0.369 & 0.190 & 1280 & 4.26e-13 & 6 & 2.55e-12 & $+1.000$ & large & *** \\
M2 vs M3 & 0.378 & 0.247 & 1280 & 4.26e-13 & 5 & 2.13e-12 & $+1.000$ & large & *** \\
M2 vs M4 & 0.378 & 0.190 & 1280 & 4.26e-13 & 4 & 1.70e-12 & $+1.000$ & large & *** \\
M1 vs M3 & 0.369 & 0.247 & 1265 & 1.47e-12 & 3 & 4.42e-12 & $+0.977$ & large & *** \\
M3 vs M4 & 0.247 & 0.190 & 1419 & 2.66e-09 & 2 & 5.31e-09 & $+0.774$ & large & *** \\
M1 vs M2 & 0.369 & 0.378 & 380 & 0.0775 & 1 & 0.0775 & $-0.258$ & small & ns \\
\bottomrule
\end{tabular}
\end{table}

\subsection{Aceleración máxima ($a_{max}$)}

\begin{table}[H]
\centering
\small
\caption{Comparaciones pareadas para Aceleración máxima ($a_{max}$).}
\renewcommand{\arraystretch}{1.15}
\begin{tabular}{lcccrrrrrl}
\toprule
\rowcolor{primary}
\textcolor{white}{\textbf{Comp.}} & \textcolor{white}{\textbf{med A}} & \textcolor{white}{\textbf{med B}} & \textcolor{white}{\textbf{$U$}} & \textcolor{white}{\textbf{$p$ raw}} & \textcolor{white}{\textbf{Mult.}} & \textcolor{white}{\textbf{$p$ Holm}} & \textcolor{white}{\textbf{$\delta(A,B)$}} & \textcolor{white}{\textbf{Mag.}} & \textcolor{white}{\textbf{Sig.}} \\
\midrule
M0 vs M3 & 12.521 & 7.258 & 971 & 0.1009 & 10 & 1.0000 & $+0.214$ & small & ns \\
M0 vs M4 & 12.521 & 9.654 & 939 & 0.1826 & 9 & 1.0000 & $+0.174$ & small & ns \\
M3 vs M4 & 7.258 & 9.654 & 664 & 0.1923 & 8 & 1.0000 & $-0.170$ & small & ns \\
M2 vs M3 & 11.040 & 7.258 & 744 & 0.2408 & 7 & 1.0000 & $+0.163$ & small & ns \\
M1 vs M3 & 11.826 & 7.258 & 732 & 0.2998 & 6 & 1.0000 & $+0.144$ & negligible & ns \\
M0 vs M2 & 12.521 & 11.040 & 706 & 0.4579 & 5 & 1.0000 & $+0.103$ & negligible & ns \\
M2 vs M4 & 11.040 & 9.654 & 679 & 0.6626 & 4 & 1.0000 & $+0.061$ & negligible & ns \\
M1 vs M4 & 11.826 & 9.654 & 678 & 0.6709 & 3 & 1.0000 & $+0.059$ & negligible & ns \\
M0 vs M1 & 12.521 & 11.826 & 674 & 0.7042 & 2 & 1.0000 & $+0.053$ & negligible & ns \\
M1 vs M2 & 11.826 & 11.040 & 539 & 0.7220 & 1 & 0.7220 & $+0.053$ & negligible & ns \\
\bottomrule
\end{tabular}
\end{table}

\subsection{Tiempo de misión}

\begin{table}[H]
\centering
\small
\caption{Comparaciones pareadas para Tiempo de misión.}
\renewcommand{\arraystretch}{1.15}
\begin{tabular}{lcccrrrrrl}
\toprule
\rowcolor{primary}
\textcolor{white}{\textbf{Comp.}} & \textcolor{white}{\textbf{med A}} & \textcolor{white}{\textbf{med B}} & \textcolor{white}{\textbf{$U$}} & \textcolor{white}{\textbf{$p$ raw}} & \textcolor{white}{\textbf{Mult.}} & \textcolor{white}{\textbf{$p$ Holm}} & \textcolor{white}{\textbf{$\delta(A,B)$}} & \textcolor{white}{\textbf{Mag.}} & \textcolor{white}{\textbf{Sig.}} \\
\midrule
M0 vs M1 & 94.011 & 109.566 & 198 & 5.64e-07 & 10 & 5.64e-06 & $-0.691$ & large & *** \\
M0 vs M2 & 94.011 & 111.652 & 199 & 5.98e-07 & 9 & 5.38e-06 & $-0.689$ & large & *** \\
M0 vs M3 & 94.011 & 114.519 & 346 & 1.28e-05 & 8 & 0.0001 & $-0.568$ & large & *** \\
M0 vs M4 & 94.011 & 108.391 & 357 & 2.06e-05 & 7 & 0.0001 & $-0.554$ & large & *** \\
M1 vs M3 & 109.566 & 114.519 & 546 & 0.2893 & 6 & 1.0000 & $-0.147$ & negligible & ns \\
M2 vs M3 & 111.652 & 114.519 & 564 & 0.3922 & 5 & 1.0000 & $-0.119$ & negligible & ns \\
M3 vs M4 & 114.519 & 108.391 & 877 & 0.4617 & 4 & 1.0000 & $+0.096$ & negligible & ns \\
M1 vs M2 & 109.566 & 111.652 & 459 & 0.4809 & 3 & 1.0000 & $-0.104$ & negligible & ns \\
M2 vs M4 & 111.652 & 108.391 & 647 & 0.9413 & 2 & 1.0000 & $+0.011$ & negligible & ns \\
M1 vs M4 & 109.566 & 108.391 & 643 & 0.9774 & 1 & 0.9774 & $+0.005$ & negligible & ns \\
\bottomrule
\end{tabular}
\end{table}

\subsection{Eficiencia de trayectoria}

\begin{table}[H]
\centering
\small
\caption{Comparaciones pareadas para Eficiencia de trayectoria.}
\renewcommand{\arraystretch}{1.15}
\begin{tabular}{lcccrrrrrl}
\toprule
\rowcolor{primary}
\textcolor{white}{\textbf{Comp.}} & \textcolor{white}{\textbf{med A}} & \textcolor{white}{\textbf{med B}} & \textcolor{white}{\textbf{$U$}} & \textcolor{white}{\textbf{$p$ raw}} & \textcolor{white}{\textbf{Mult.}} & \textcolor{white}{\textbf{$p$ Holm}} & \textcolor{white}{\textbf{$\delta(A,B)$}} & \textcolor{white}{\textbf{Mag.}} & \textcolor{white}{\textbf{Sig.}} \\
\midrule
M0 vs M3 & 0.996 & 0.845 & 1600 & 1.44e-14 & 10 & 1.44e-13 & $+1.000$ & large & *** \\
M0 vs M4 & 0.996 & 0.825 & 1600 & 1.44e-14 & 9 & 1.29e-13 & $+1.000$ & large & *** \\
M0 vs M1 & 0.996 & 0.811 & 1280 & 4.26e-13 & 8 & 3.41e-12 & $+1.000$ & large & *** \\
M0 vs M2 & 0.996 & 0.816 & 1280 & 4.26e-13 & 7 & 2.98e-12 & $+1.000$ & large & *** \\
M2 vs M3 & 0.816 & 0.845 & 446 & 0.0283 & 6 & 0.1699 & $-0.303$ & small & ns \\
M2 vs M4 & 0.816 & 0.825 & 489 & 0.0881 & 5 & 0.4405 & $-0.236$ & small & ns \\
M1 vs M2 & 0.811 & 0.816 & 578 & 0.3791 & 4 & 1.0000 & $+0.129$ & negligible & ns \\
M1 vs M4 & 0.811 & 0.825 & 572 & 0.4443 & 3 & 1.0000 & $-0.106$ & negligible & ns \\
M1 vs M3 & 0.811 & 0.845 & 576 & 0.4718 & 2 & 0.9435 & $-0.100$ & negligible & ns \\
M3 vs M4 & 0.845 & 0.825 & 766 & 0.7472 & 1 & 0.7472 & $-0.043$ & negligible & ns \\
\bottomrule
\end{tabular}
\end{table}

\newpage
%==================================================================
\section{Foco en M4: comparaciones críticas}
%==================================================================

A continuación se aíslan las cuatro comparaciones M4 vs cada otro método, para las métricas de mayor relevancia teórica del estudio.

\subsection{Distancia mínima ($d^\star$)}

\begin{table}[H]
\centering
\caption{Comparaciones M4 vs todos los demás (en orden por significancia).}
\begin{tabular}{lrrrrl}
\toprule
\textbf{Comp.} & \textbf{$p$ raw} & \textbf{$p$ Holm} & \textbf{$\delta(M4, X)$} & \textbf{Mag.} & \textbf{Sig.} \\
\midrule
M4 vs M0 & 1.11e-12 & 1.00e-11 & $+0.925$ & large & *** \\
M4 vs M3 & 2.37e-07 & 9.50e-07 & $-0.672$ & large & *** \\
M4 vs M2 & 0.0025 & 0.0076 & $+0.417$ & medium & ** \\
M4 vs M1 & 0.0043 & 0.0086 & $+0.395$ & medium & ** \\
\bottomrule
\end{tabular}
\end{table}

\subsection{Tiempo en zona personal}

\begin{table}[H]
\centering
\caption{Comparaciones M4 vs todos los demás (en orden por significancia).}
\begin{tabular}{lrrrrl}
\toprule
\textbf{Comp.} & \textbf{$p$ raw} & \textbf{$p$ Holm} & \textbf{$\delta(M4, X)$} & \textbf{Mag.} & \textbf{Sig.} \\
\midrule
M4 vs M0 & 0.0004 & 0.0029 & $-0.464$ & medium & ** \\
M4 vs M1 & 0.0004 & 0.0030 & $-0.486$ & large & ** \\
M4 vs M3 & 0.0680 & 0.2039 & $+0.237$ & small & ns \\
M4 vs M2 & 0.1984 & 0.3967 & $-0.178$ & small & ns \\
\bottomrule
\end{tabular}
\end{table}

\subsection{Número de paradas}

\begin{table}[H]
\centering
\caption{Comparaciones M4 vs todos los demás (en orden por significancia).}
\begin{tabular}{lrrrrl}
\toprule
\textbf{Comp.} & \textbf{$p$ raw} & \textbf{$p$ Holm} & \textbf{$\delta(M4, X)$} & \textbf{Mag.} & \textbf{Sig.} \\
\midrule
M4 vs M2 & 5.08e-14 & 2.54e-13 & $-0.988$ & large & *** \\
M4 vs M1 & 3.78e-14 & 2.64e-13 & $-0.994$ & large & *** \\
M4 vs M3 & 0.0142 & 0.0426 & $+0.175$ & small & * \\
M4 vs M0 & 1.0000 & 1.0000 & $-0.000$ & negligible & ns \\
\bottomrule
\end{tabular}
\end{table}

\subsection{Desviación estándar de velocidad ($\sigma_v$)}

\begin{table}[H]
\centering
\caption{Comparaciones M4 vs todos los demás (en orden por significancia).}
\begin{tabular}{lrrrrl}
\toprule
\textbf{Comp.} & \textbf{$p$ raw} & \textbf{$p$ Holm} & \textbf{$\delta(M4, X)$} & \textbf{Mag.} & \textbf{Sig.} \\
\midrule
M4 vs M0 & 1.44e-14 & 1.29e-13 & $+1.000$ & large & *** \\
M4 vs M2 & 4.26e-13 & 1.70e-12 & $-1.000$ & large & *** \\
M4 vs M1 & 4.26e-13 & 2.55e-12 & $-1.000$ & large & *** \\
M4 vs M3 & 2.66e-09 & 5.31e-09 & $-0.774$ & large & *** \\
\bottomrule
\end{tabular}
\end{table}


%==================================================================
\section{Conclusiones del análisis post-hoc}
%==================================================================

\begin{enumerate}[leftmargin=*]
    \item \textbf{Suavidad cinemática ($\sigma_v$):} M4 difiere significativamente de \emph{todos} los demás métodos con tamaño de efecto \emph{large}. M4 presenta menor $\sigma_v$ que M1, M2 y M3 ($\delta < 0$ en los tres, $|\delta| \geq 0.77$).
    
    \item \textbf{Distancia mínima ($d^\star$):} M4 supera estadísticamente a M0/M1/M2 ($\delta > 0$ medium/large) y es superado por M3 ($\delta = -0.672$, large). Esto evidencia el sobre-distanciamiento de M3.
    
    \item \textbf{Número de paradas ($N_{stop}$):} M4 supera dramáticamente a M1 y M2 ($\delta < -0.98$, large), pero no se distingue de M3 ni de M0 (que tienen valor cercano a cero).
    
    \item \textbf{Aceleración máxima ($a_{max}$):} ninguna comparación es significativa después de Holm--Bonferroni. Coherente con el resultado del Documento 2 (Kruskal--Wallis no significativo para $a_{max}$).
    
    \item \textbf{Justificación operacional de M4:} M4 muestra ventaja robusta y estadísticamente significativa en suavidad cinemática frente a todos los demás métodos, mientras que en otras métricas su ventaja es competitiva con M3 (campos continuos comparables) y dominante frente a M0--M2 (NavMesh puro y supervisores binarios).
\end{enumerate}



\chapter{Analisis de tasas de exito: Fisher exacto y $\chi^2$}
\label{chap:fisher}




%==================================================================
\section{Marco teórico}
%==================================================================

\subsection{Tabla de contingencia para tasa de éxito}

Sea $\eta_s$ la tasa de éxito de un método en un escenario, definida como el cociente entre corridas que completaron el ciclo $A \to B \to A$ y corridas planificadas:
\begin{equation}
\eta_s(M_j, E_i) = \frac{N_{success}(M_j, E_i)}{N_{planned}(M_j, E_i)} \cdot 100\%
\end{equation}

Para comparar dos métodos en un mismo escenario, construimos una tabla $2\times 2$:

\begin{table}[H]
\centering
\begin{tabular}{lcc}
\toprule
& \textbf{Éxito} & \textbf{Fallo} \\
\midrule
Método A & $a$ & $b$ \\
Método B & $c$ & $d$ \\
\bottomrule
\end{tabular}
\end{table}

con $a+b = n_A$ (corridas planificadas para A) y $c+d = n_B$ (planificadas para B).

\subsection{Prueba exacta de Fisher}

La prueba exacta de Fisher es preferible a $\chi^2$ cuando alguna celda esperada $E_{ij}$ es menor que 5, lo cual es común con muestras pequeñas (en este estudio, $n=10$ por celda). La prueba se basa en la \emph{distribución hipergeométrica} bajo la hipótesis nula de independencia.

\paragraph{Hipergeométrica:} dado que los marginales $a+b$, $c+d$, $a+c$ y $b+d$ están fijos por construcción, la probabilidad de observar exactamente el valor $a$ en la primera celda es:
\begin{equation}
P(X = a \mid \text{marginales fijos}) = \frac{\dbinom{a+b}{a}\dbinom{c+d}{c}}{\dbinom{N}{a+c}}
= \frac{(a+b)!\,(c+d)!\,(a+c)!\,(b+d)!}{a!\,b!\,c!\,d!\,N!}
\end{equation}
con $N = a+b+c+d$.

\paragraph{$p$-valor exacto bilateral:}
\begin{equation}
p_{\text{Fisher}} = \sum_{k:\,P(X=k) \leq P(X=a)} P(X=k)
\end{equation}
Esto suma todas las configuraciones de la tabla \emph{tan o más extremas} que la observada bajo la hipergeométrica.

\subsection{Prueba $\chi^2$ de independencia para tabla $5\times 2$}

Para comparar simultáneamente las tasas de éxito de los 5 métodos, se construye la tabla $5\times 2$ de éxitos vs fallos, y se aplica:
\begin{equation}
\boxed{\chi^2 = \sum_{i=1}^{r}\sum_{j=1}^{c} \frac{(O_{ij} - E_{ij})^2}{E_{ij}}}
\end{equation}
donde $O_{ij}$ es el conteo observado, y $E_{ij}$ es el conteo esperado bajo independencia:
\begin{equation}
E_{ij} = \frac{R_i \cdot C_j}{N}, \qquad R_i = \sum_j O_{ij}, \quad C_j = \sum_i O_{ij}
\end{equation}

Bajo $H_0$, $\chi^2 \sim \chi^2_{((r-1)(c-1))}$. Para $r=5, c=2$, $df = 4$.

\subsection{Cramér's V (tamaño de efecto)}

\begin{equation}
\boxed{V = \sqrt{\frac{\chi^2}{N \cdot \min(r-1, c-1)}}}
\end{equation}

Para tabla $5\times 2$: $\min(r-1, c-1) = \min(4, 1) = 1$, por lo que:
\begin{equation}
V = \sqrt{\frac{\chi^2}{N}}
\end{equation}

\paragraph{Interpretación (Cohen, 1988, adaptado por Rea \& Parker):}
\begin{table}[H]
\centering
\begin{tabular}{cl}
\toprule
$V$ & Magnitud \\
\midrule
$< 0.10$ & negligible \\
$\geq 0.10,\ <0.20$ & weak \\
$\geq 0.20,\ <0.40$ & moderate \\
$\geq 0.40,\ <0.60$ & relatively strong \\
$\geq 0.60$ & very strong \\
\bottomrule
\end{tabular}
\end{table}

\newpage
%==================================================================
\section{Procedimiento paso a paso (caso M4 vs M1 en E4)}
%==================================================================

A continuación se desarrolla el cálculo completo de Fisher exacto para la comparación más relevante del estudio: M4 vs M1 en escenario E4.

\subsection{Paso 1: Construcción de la tabla $2\times 2$}

\begin{table}[H]
\centering
\caption{Tabla de contingencia: M4 vs M1 en E4.}
\renewcommand{\arraystretch}{1.4}
\begin{tabular}{lccc}
\toprule
& \textbf{Éxito} & \textbf{Fallo} & \textbf{Total} \\
\midrule
M4 & $a = 10$ & $b = 0$ & $a+b = 10$ \\
M1 & $c = 2$  & $d = 8$ & $c+d = 10$ \\
\midrule
Total & $a+c = 12$ & $b+d = 8$ & $N = 20$ \\
\bottomrule
\end{tabular}
\end{table}

\subsection{Paso 2: Probabilidad bajo hipergeométrica}

Para la tabla observada, calculamos:
\begin{equation}
P(X = 10) = \frac{\binom{10}{10}\binom{10}{2}}{\binom{20}{12}}
\end{equation}

Calculando los coeficientes binomiales:
\begin{itemize}
    \item $\binom{10}{10} = 1$
    \item $\binom{10}{2} = \dfrac{10!}{2!\cdot 8!} = \dfrac{10\cdot 9}{2} = 45$
    \item $\binom{20}{12} = \dfrac{20!}{12!\cdot 8!} = 125\,970$
\end{itemize}

\begin{equation}
P(X = 10) = \frac{1 \cdot 45}{125\,970} = \frac{45}{125\,970} \approx 3.572 \times 10^{-4}
\end{equation}

\subsection{Paso 3: Enumerar todas las tablas y sumar las "tan o más extremas"}

El soporte de la hipergeométrica para $X$ (la celda superior izquierda) es $\{\max(0, n_1+m_1-N), \ldots, \min(n_1, m_1)\} = \{2, \ldots, 10\}$, donde $n_1 = a+b = 10$, $m_1 = a+c = 12$.

\begin{table}[H]
\centering
\caption{Probabilidades de cada configuración bajo $H_0$.}
\renewcommand{\arraystretch}{1.2}
\begin{tabular}{rrrrr}
\toprule
$k$ & Tabla $(a,b,c,d)$ & $P(X=k)$ & ¿Es $\leq P(X=10)$? & Suma acumulada \\
\midrule
2 & (2,8,10,0) & 0.000357 & Sí & 0.000357 \\
3 & (3,7,9,1) & 0.009526 & No & — \\
4 & (4,6,8,2) & 0.075018 & No & — \\
5 & (5,5,7,3) & 0.240057 & No & — \\
6 & (6,4,6,4) & 0.350083 & No & — \\
7 & (7,3,5,5) & 0.240057 & No & — \\
8 & (8,2,4,6) & 0.075018 & No & — \\
9 & (9,1,3,7) & 0.009526 & No & — \\
10 & (10,0,2,8) & 0.000357 & Sí & 0.000714 \\
\midrule
\multicolumn{2}{r}{$p$-valor exacto bilateral:} & & & $\boxed{7.144558e-04}$ \\
\bottomrule
\end{tabular}
\end{table}

\subsection{Paso 4: Conclusión}

\begin{itemize}
    \item $p_{\text{Fisher}}^{\text{two-sided}} = 7.1446e-04$
    \item Como $p < 0.001$, se rechaza $H_0$ con confianza > 99.9\%.
    \item La diferencia entre M4 (10/10 éxitos) y M1 (2/10 éxitos) en E4 NO es producto del azar.
\end{itemize}

\textbf{Interpretación operacional:} M4 supera de manera estadísticamente robusta a M1 en escenarios congestionados con agentes humanos estáticos. Este es el hallazgo central del estudio.

\newpage
%==================================================================
\section{Resultados completos: Fisher exacto por par de métodos}
%==================================================================

Para cada par de métodos $(M_a, M_b)$ con $a < b$ y para cada escenario, se reporta el $p$-valor exacto de Fisher. Se omiten los casos donde ambos métodos tienen tasa de éxito 100\% (Fisher $p = 1$ trivialmente).

\subsection{Escenario E1}

\begin{table}[H]
\centering
\caption{Fisher exacto bilateral: comparaciones pareadas en E1.}
\begin{tabular}{lcccrl}
\toprule
\textbf{Comp.} & \textbf{$a$ (succ A)} & \textbf{$c$ (succ B)} & \textbf{Tabla} & \textbf{$p$ Fisher} & \textbf{Sig.} \\
\midrule
\multicolumn{6}{c}{\textit{No hay diferencias en este escenario --- todos los métodos exitosos.}} \\
\bottomrule
\end{tabular}
\end{table}

\subsection{Escenario E2}

\begin{table}[H]
\centering
\caption{Fisher exacto bilateral: comparaciones pareadas en E2.}
\begin{tabular}{lcccrl}
\toprule
\textbf{Comp.} & \textbf{$a$ (succ A)} & \textbf{$c$ (succ B)} & \textbf{Tabla} & \textbf{$p$ Fisher} & \textbf{Sig.} \\
\midrule
\multicolumn{6}{c}{\textit{No hay diferencias en este escenario --- todos los métodos exitosos.}} \\
\bottomrule
\end{tabular}
\end{table}

\subsection{Escenario E3}

\begin{table}[H]
\centering
\caption{Fisher exacto bilateral: comparaciones pareadas en E3.}
\begin{tabular}{lcccrl}
\toprule
\textbf{Comp.} & \textbf{$a$ (succ A)} & \textbf{$c$ (succ B)} & \textbf{Tabla} & \textbf{$p$ Fisher} & \textbf{Sig.} \\
\midrule
\multicolumn{6}{c}{\textit{No hay diferencias en este escenario --- todos los métodos exitosos.}} \\
\bottomrule
\end{tabular}
\end{table}

\subsection{Escenario E4}

\begin{table}[H]
\centering
\caption{Fisher exacto bilateral: comparaciones pareadas en E4.}
\begin{tabular}{lcccrl}
\toprule
\textbf{Comp.} & \textbf{$a$ (succ A)} & \textbf{$c$ (succ B)} & \textbf{Tabla} & \textbf{$p$ Fisher} & \textbf{Sig.} \\
\midrule
M0 vs M1 & 10/10 & 2/10 & (10,0,2,8) & 0.0007 & *** \\
M0 vs M2 & 10/10 & 2/10 & (10,0,2,8) & 0.0007 & *** \\
M1 vs M2 & 2/10 & 2/10 & (2,8,2,8) & 1.0000 & ns \\
M1 vs M3 & 2/10 & 10/10 & (2,8,10,0) & 0.0007 & *** \\
M1 vs M4 & 2/10 & 10/10 & (2,8,10,0) & 0.0007 & *** \\
M2 vs M3 & 2/10 & 10/10 & (2,8,10,0) & 0.0007 & *** \\
M2 vs M4 & 2/10 & 10/10 & (2,8,10,0) & 0.0007 & *** \\
\bottomrule
\end{tabular}
\end{table}

\newpage
%==================================================================
\section{Prueba $\chi^2$ global de independencia}
%==================================================================

Se construye la tabla de contingencia $5 \times 2$ con todos los métodos y los conteos globales de éxito/fallo (sobre los $5 \times 4 \times 10 = 200$ ensayos planificados).

\subsection{Tabla observada}

\begin{table}[H]
\centering
\caption{Tabla de contingencia $5\times 2$ observada.}
\renewcommand{\arraystretch}{1.3}
\begin{tabular}{lccc}
\toprule
\textbf{Método} & \textbf{Éxito ($O_{i1}$)} & \textbf{Fallo ($O_{i2}$)} & \textbf{Total ($R_i$)} \\
\midrule
M0 & 40 & 0 & 40 \\
M1 & 32 & 8 & 40 \\
M2 & 32 & 8 & 40 \\
M3 & 40 & 0 & 40 \\
M4 & 40 & 0 & 40 \\
\midrule
\textbf{Total $C_j$} & 184 & 16 & 200 \\
\bottomrule
\end{tabular}
\end{table}

\subsection{Tabla esperada bajo $H_0$}

Bajo independencia, el conteo esperado en la celda $(i,j)$ es:
\begin{equation}
E_{ij} = \frac{R_i \cdot C_j}{N}
\end{equation}

\begin{table}[H]
\centering
\caption{Tabla esperada bajo independencia.}
\renewcommand{\arraystretch}{1.3}
\begin{tabular}{lccc}
\toprule
\textbf{Método} & \textbf{$E_{i,\text{éxito}}$} & \textbf{$E_{i,\text{fallo}}$} & \textbf{Total} \\
\midrule
M0 & 36.80 & 3.20 & 40.00 \\
M1 & 36.80 & 3.20 & 40.00 \\
M2 & 36.80 & 3.20 & 40.00 \\
M3 & 36.80 & 3.20 & 40.00 \\
M4 & 36.80 & 3.20 & 40.00 \\
\bottomrule
\end{tabular}
\end{table}

\subsection{Cálculo del estadístico $\chi^2$}

\begin{equation}
\chi^2 = \sum_{i,j} \frac{(O_{ij} - E_{ij})^2}{E_{ij}}
\end{equation}

Calculando término a término:

\begin{table}[H]
\centering
\renewcommand{\arraystretch}{1.2}
\begin{tabular}{lccccc}
\toprule
\textbf{Método} & $O_{i,\text{éxito}}$ & $E_{i,\text{éxito}}$ & $\dfrac{(O-E)^2}{E}_{\text{éxito}}$ & $\dfrac{(O-E)^2}{E}_{\text{fallo}}$ & Subtotal \\
\midrule
M0 & 40 & 36.80 & 0.2783 & 3.2000 & 3.4783 \\
M1 & 32 & 36.80 & 0.6261 & 7.2000 & 7.8261 \\
M2 & 32 & 36.80 & 0.6261 & 7.2000 & 7.8261 \\
M3 & 40 & 36.80 & 0.2783 & 3.2000 & 3.4783 \\
M4 & 40 & 36.80 & 0.2783 & 3.2000 & 3.4783 \\
\midrule
\multicolumn{5}{r}{$\chi^2 =$} & \textbf{26.0870} \\
\bottomrule
\end{tabular}
\end{table}

\subsection{Comparación con punto crítico}

\begin{itemize}
    \item $\chi^2_{\text{calc}} = 26.087$
    \item Grados de libertad: $df = (r-1)(c-1) = 4 \cdot 1 = 4$
    \item Punto crítico al 0.001: $\chi^2_{(4),0.001} = 18.467$
    \item Como $\chi^2_{\text{calc}} = 26.087 > 18.467$, se rechaza $H_0$ con confianza > 99.9\%
    \item $p$-valor: $3.03923e-05$
\end{itemize}

\subsection{Cramér's V}

\begin{align}
V &= \sqrt{\frac{\chi^2}{N \cdot \min(r-1, c-1)}} \\
  &= \sqrt{\frac{26.087}{200 \cdot 1}} \\
  &= \sqrt{0.130435} \\
  &= \boxed{0.3612}
\end{align}

Magnitud: $V = 0.361$ corresponde a efecto \textbf{moderate-to-relatively strong} (entre 0.20 y 0.40, cerca del límite superior).

\textbf{Interpretación:} la magnitud del efecto es sustancial: el método elegido explica alrededor del $V^2 = 13.0\%$ de la varianza en la tasa de éxito.

\newpage
%==================================================================
\section{Conclusiones}
%==================================================================

\begin{enumerate}[leftmargin=*]
    \item \textbf{Hipótesis H6 validada con la mayor fuerza estadística del estudio:} Fisher exacto $p = 7.14 \times 10^{-4}$ tanto para M4 vs M1 como para M4 vs M2 en E4. Estos $p$-valores son robustos al uso de Bonferroni para múltiples comparaciones de Fisher (no hay corrección que los lleve por encima de 0.05).
    
    \item \textbf{$\chi^2$ global confirma el efecto:} $\chi^2 = 26.09$, $df = 4$, $p < 0.001$, Cramér's $V = 0.361$. La magnitud del efecto es sustancial y refleja el bloqueo operacional de M1/M2.
    
    \item \textbf{Localización del efecto:} la diferencia significativa se concentra completamente en E4. En E1, E2 y E3 todos los métodos alcanzan tasa de éxito 100\% (no hay diferencias estadísticas). El bloqueo operacional es específicamente sensible a la presencia de \emph{agentes humanos simulados estáticos en zonas congestionadas}.
    
    \item \textbf{Equivalencia M3 vs M4:} ninguno de los pares M3 vs M4 muestra diferencia significativa en tasa de éxito (ambos 100\%). La distinción entre estos dos métodos requiere otras métricas (suavidad cinemática, distancia mínima), abordadas en el Documento 3.
\end{enumerate}



\part{Analisis multicriterio y robustez}

\chapter{Ranking multicriterio TOPSIS}
\label{chap:topsis}




%==================================================================
\section{Marco teórico de TOPSIS}
%==================================================================

\subsection{Origen y principio}

TOPSIS (Technique for Order of Preference by Similarity to Ideal Solution) es un método de toma de decisiones multicriterio (MCDM) propuesto por Hwang y Yoon (1981). Su principio operativo es:

\begin{tcolorbox}[colback=lightgray, colframe=primary, boxrule=0.5pt]
\textit{La alternativa elegida debe estar a la \textbf{menor distancia} de la \emph{solución ideal positiva} y a la \textbf{mayor distancia} de la \emph{solución ideal negativa}.}
\end{tcolorbox}

\subsection{Fundamentos matemáticos}

\subsubsection{Matriz de decisión}

Sea un problema con $m$ alternativas y $n$ criterios. La matriz de decisión $\mathbf{X}$ es de dimensión $m \times n$:
\begin{equation}
\mathbf{X} = \begin{bmatrix}
x_{11} & x_{12} & \cdots & x_{1n} \\
x_{21} & x_{22} & \cdots & x_{2n} \\
\vdots & \vdots & \ddots & \vdots \\
x_{m1} & x_{m2} & \cdots & x_{mn}
\end{bmatrix}
\end{equation}

donde $x_{ij}$ es el valor de la alternativa $i$ en el criterio $j$.

\subsubsection{Vector de pesos}

Sea $\mathbf{w} = (w_1, w_2, \ldots, w_n)$ con $\sum_j w_j = 1$ y $w_j > 0$ para todo $j$.

\subsubsection{Naturaleza de los criterios}

Cada criterio $j$ es de tipo:
\begin{itemize}
    \item \textbf{Benefit (beneficio)}: \emph{mayor es mejor} (e.g., $d^\star$, $\eta_s$).
    \item \textbf{Cost (costo)}: \emph{menor es mejor} (e.g., $T_{personal}$, $\sigma_v$, $a_{max}$).
\end{itemize}

\subsection{Los nueve pasos del algoritmo TOPSIS}

\paragraph{Paso 1: Normalización vectorial.} Para eliminar la dependencia de unidades:
\begin{equation}
\tilde{x}_{ij} = \frac{x_{ij}}{\sqrt{\sum_{k=1}^{m} x_{kj}^2}}
\label{doc5:eq:topsis1}
\end{equation}

Esta normalización (también llamada \emph{Euclidean normalization}) preserva las proporciones relativas y es la convención estándar de TOPSIS clásico.

\paragraph{Paso 2: Matriz ponderada.}
\begin{equation}
v_{ij} = w_j \cdot \tilde{x}_{ij}
\label{doc5:eq:topsis2}
\end{equation}

\paragraph{Paso 3: Solución ideal positiva $A^+$.}
\begin{equation}
A^+ = (v_1^+, v_2^+, \ldots, v_n^+)
\label{doc5:eq:topsis3}
\end{equation}
donde:
\begin{equation}
v_j^+ = \begin{cases}
\max_i v_{ij}, & \text{si } j \text{ es benefit} \\
\min_i v_{ij}, & \text{si } j \text{ es cost}
\end{cases}
\end{equation}

\paragraph{Paso 4: Solución ideal negativa $A^-$.}
\begin{equation}
A^- = (v_1^-, v_2^-, \ldots, v_n^-)
\end{equation}
donde:
\begin{equation}
v_j^- = \begin{cases}
\min_i v_{ij}, & \text{si } j \text{ es benefit} \\
\max_i v_{ij}, & \text{si } j \text{ es cost}
\end{cases}
\end{equation}

\paragraph{Paso 5: Distancia de cada alternativa al ideal positivo.}
\begin{equation}
D_i^+ = \sqrt{\sum_{j=1}^{n}(v_{ij} - v_j^+)^2}
\label{doc5:eq:topsis5}
\end{equation}

\paragraph{Paso 6: Distancia de cada alternativa al ideal negativo.}
\begin{equation}
D_i^- = \sqrt{\sum_{j=1}^{n}(v_{ij} - v_j^-)^2}
\label{doc5:eq:topsis6}
\end{equation}

\paragraph{Paso 7: Closeness coefficient.}
\begin{equation}
\boxed{CC_i = \frac{D_i^-}{D_i^+ + D_i^-}}
\label{doc5:eq:topsis7}
\end{equation}

con $CC_i \in [0, 1]$. La alternativa más cercana al ideal positivo y más lejana del ideal negativo tiene $CC_i$ máximo.

\paragraph{Paso 8: Ranking.} Se ordenan las alternativas por $CC_i$ descendente. La alternativa con mayor $CC$ es óptima.

\paragraph{Paso 9: Análisis de sensibilidad de pesos.} Se evalúa la robustez del ranking ante perturbaciones de $\mathbf{w}$ (este análisis se documenta en el Documento 6).

\newpage
%==================================================================
\section{Configuración del problema en este estudio}
%==================================================================

\subsection{Alternativas, criterios y pesos}

\begin{table}[H]
\centering
\caption{Definición del problema TOPSIS aplicado al estudio.}
\renewcommand{\arraystretch}{1.3}
\begin{tabular}{ll}
\toprule
\textbf{Alternativas ($m=5$)} & M0, M1, M2, M3, M4 \\
\textbf{Criterios ($n=7$)} & Detallados abajo \\
\bottomrule
\end{tabular}
\end{table}

\begin{table}[H]
\centering
\caption{Criterios, tipos y pesos.}
\renewcommand{\arraystretch}{1.25}
\begin{tabular}{lllr}
\toprule
\textbf{$j$} & \textbf{Criterio} & \textbf{Tipo} & \textbf{Peso $w_j$} \\
\midrule
1 & Distancia mínima ($d^\star$) & benefit & 0.20 \\
2 & Tiempo en zona personal ($T_{personal}$) & cost & 0.15 \\
3 & Tiempo de misión ($T_{mission}$) & cost & 0.10 \\
4 & Número de paradas ($N_{stop}$) & cost & 0.15 \\
5 & Suavidad cinemática ($\sigma_v$) & cost & 0.15 \\
6 & Aceleración máxima ($a_{max}$) & cost & 0.10 \\
7 & Tasa de éxito ($\eta_s$) & benefit & 0.15 \\
\midrule
& \textbf{Suma} & & \textbf{1.00} \\
\bottomrule
\end{tabular}
\end{table}

\subsection{Justificación de los pesos}

Los pesos asignados reflejan las prioridades operacionales del problema de navegación social agrícola en gemelo digital:
\begin{itemize}
    \item \textbf{Seguridad social ($d^\star$, $T_{personal}$, $\eta_s$)}: 50\% combinado. Es el aspecto más crítico --- una colisión social es inaceptable.
    \item \textbf{Calidad cinemática ($N_{stop}$, $\sigma_v$, $a_{max}$)}: 40\% combinado. Refleja confort y predictibilidad del comportamiento.
    \item \textbf{Eficiencia ($T_{mission}$)}: 10\%. Importante pero secundaria frente a seguridad y suavidad.
\end{itemize}

\subsection{Tratamiento de fallos: análisis paralelo}

Como se detalló en el Documento 1, las corridas fallidas (M1 y M2 en E4) requieren tratamiento explícito. Se realizan dos análisis paralelos:

\paragraph{Análisis A (sin penalización):} cada celda de la matriz de decisión se calcula como el promedio de la métrica solo sobre las corridas exitosas. La tasa de éxito $\eta_s$ entra como séptimo criterio y captura indirectamente el bloqueo operacional.

\paragraph{Análisis B (con penalización):} las corridas fallidas se reemplazan por el peor valor observado para cada criterio antes de promediar. Específicamente:
\begin{itemize}
    \item Para criterios benefit: peor = mínimo del dataset.
    \item Para criterios cost: peor = máximo del dataset.
    \item Para $T_{mission}$ específicamente: peor = $T_{\max} = 180$ s (timeout).
\end{itemize}
Este análisis penaliza directamente los métodos M1 y M2 en sus métricas continuas.

\newpage
%==================================================================
\section{Análisis A: sin penalización de fallos}
%==================================================================

\subsection{Matriz de decisión $\mathbf{X}$}

\begin{table}[H]
\centering
\caption{Matriz de decisión $\mathbf{X}$ (Análisis A, sin penalizar fallos).}
\renewcommand{\arraystretch}{1.3}
\small
\begin{tabular}{lccccccc}
\toprule
\rowcolor{primary}
\textcolor{white}{\textbf{Alt.}} &
\textcolor{white}{$d^\star$} &
\textcolor{white}{$T_{pers}$} &
\textcolor{white}{$T_{mis}$} &
\textcolor{white}{$N_{stop}$} &
\textcolor{white}{$\sigma_v$} &
\textcolor{white}{$a_{max}$} &
\textcolor{white}{$\eta_s$} \\
\rowcolor{secondary}
& benefit & cost & cost & cost & cost & cost & benefit \\
\midrule
M0 & 0.5627 & 25.3652 & 93.6306 & 0.2000 & 0.0663 & 12.4303 & 1.0000 \\
M1 & 0.7504 & 20.3225 & 107.9230 & 9.9062 & 0.3783 & 11.2753 & 0.8000 \\
M2 & 0.7453 & 12.4282 & 109.3268 & 3.2188 & 0.4101 & 11.0107 & 0.8000 \\
M3 & 1.0038 & 7.5740 & 115.0431 & 0.0250 & 0.2428 & 9.3710 & 1.0000 \\
M4 & 0.8459 & 11.7202 & 110.6021 & 0.2000 & 0.1952 & 10.2838 & 1.0000 \\
\bottomrule
\end{tabular}
\end{table}

\subsection{Paso 1: Normalización vectorial}

Se calcula la norma de cada columna $\sqrt{\sum_i x_{ij}^2}$:

\begin{table}[H]
\centering
\renewcommand{\arraystretch}{1.3}
\small
\begin{tabular}{lccccccc}
\toprule
\textbf{Cantidad} & $d^\star$ & $T_{pers}$ & $T_{mis}$ & $N_{stop}$ & $\sigma_v$ & $a_{max}$ & $\eta_s$ \\
\midrule
$\sum_i x_{ij}^2$ & 3.15838 & 1405.59074 & 57834.15635 & 108.57477 & 0.41271 & 596.45119 & 4.28000 \\
$\sqrt{\sum_i x_{ij}^2}$ & 1.77718 & 37.49121 & 240.48733 & 10.41992 & 0.64243 & 24.42235 & 2.06882 \\
\bottomrule
\end{tabular}
\end{table}

Aplicando $\tilde{x}_{ij} = x_{ij}/\sqrt{\sum_k x_{kj}^2}$:

\begin{table}[H]
\centering
\caption{Matriz normalizada $\tilde{\mathbf{X}}$ (Análisis A).}
\renewcommand{\arraystretch}{1.3}
\small
\begin{tabular}{lccccccc}
\toprule
\rowcolor{primary}
\textcolor{white}{\textbf{Alt.}} &
\textcolor{white}{$d^\star$} &
\textcolor{white}{$T_{pers}$} &
\textcolor{white}{$T_{mis}$} &
\textcolor{white}{$N_{stop}$} &
\textcolor{white}{$\sigma_v$} &
\textcolor{white}{$a_{max}$} &
\textcolor{white}{$\eta_s$} \\
\midrule
M0 & 0.3167 & 0.6766 & 0.3893 & 0.0192 & 0.1032 & 0.5090 & 0.4834 \\
M1 & 0.4223 & 0.5421 & 0.4488 & 0.9507 & 0.5889 & 0.4617 & 0.3867 \\
M2 & 0.4194 & 0.3315 & 0.4546 & 0.3089 & 0.6383 & 0.4508 & 0.3867 \\
M3 & 0.5648 & 0.2020 & 0.4784 & 0.0024 & 0.3779 & 0.3837 & 0.4834 \\
M4 & 0.4760 & 0.3126 & 0.4599 & 0.0192 & 0.3038 & 0.4211 & 0.4834 \\
\bottomrule
\end{tabular}
\end{table}

\subsection{Paso 2: Matriz ponderada}

Multiplicando cada columna $\tilde{x}_{ij}$ por su peso $w_j$:

\begin{table}[H]
\centering
\caption{Matriz ponderada $\mathbf{V}$ (Análisis A).}
\renewcommand{\arraystretch}{1.3}
\small
\begin{tabular}{lccccccc}
\toprule
\rowcolor{primary}
\textcolor{white}{\textbf{Alt.}} &
\textcolor{white}{$d^\star$ (0.20)} &
\textcolor{white}{$T_{pers}$ (0.15)} &
\textcolor{white}{$T_{mis}$ (0.10)} &
\textcolor{white}{$N_{stop}$ (0.15)} &
\textcolor{white}{$\sigma_v$ (0.15)} &
\textcolor{white}{$a_{max}$ (0.10)} &
\textcolor{white}{$\eta_s$ (0.15)} \\
\midrule
M0 & 0.06333 & 0.10148 & 0.03893 & 0.00288 & 0.01548 & 0.05090 & 0.07251 \\
M1 & 0.08445 & 0.08131 & 0.04488 & 0.14261 & 0.08833 & 0.04617 & 0.05800 \\
M2 & 0.08388 & 0.04972 & 0.04546 & 0.04634 & 0.09575 & 0.04508 & 0.05800 \\
M3 & 0.11296 & 0.03030 & 0.04784 & 0.00036 & 0.05669 & 0.03837 & 0.07251 \\
M4 & 0.09519 & 0.04689 & 0.04599 & 0.00288 & 0.04558 & 0.04211 & 0.07251 \\
\bottomrule
\end{tabular}
\end{table}

\subsection{Pasos 3 y 4: Soluciones ideales positiva y negativa}

Para cada criterio:
\begin{itemize}
    \item Si es \emph{benefit}: $A^+ = \max_i v_{ij}$, $A^- = \min_i v_{ij}$.
    \item Si es \emph{cost}: $A^+ = \min_i v_{ij}$, $A^- = \max_i v_{ij}$.
\end{itemize}

\begin{table}[H]
\centering
\renewcommand{\arraystretch}{1.3}
\small
\begin{tabular}{lccccccc}
\toprule
& $d^\star$ & $T_{pers}$ & $T_{mis}$ & $N_{stop}$ & $\sigma_v$ & $a_{max}$ & $\eta_s$ \\
& (B) & (C) & (C) & (C) & (C) & (C) & (B) \\
\midrule
$A^+$ & 0.11296 & 0.03030 & 0.03893 & 0.00036 & 0.01548 & 0.03837 & 0.07251 \\
$A^-$ & 0.06333 & 0.10148 & 0.04784 & 0.14261 & 0.09575 & 0.05090 & 0.05800 \\
\bottomrule
\end{tabular}
\caption{Soluciones ideales (B = benefit, C = cost).}
\end{table}

\subsection{Pasos 5 y 6: Distancias $D^+$ y $D^-$}

Para cada alternativa $i$, se calcula la distancia euclidiana al ideal positivo y al ideal negativo. Por ejemplo, para M4:

\begin{table}[H]
\centering
\caption{Cálculo paso a paso de $D_{M4}^+$ y $D_{M4}^-$ (Análisis A).}
\renewcommand{\arraystretch}{1.2}
\small
\begin{tabular}{lcccccc}
\toprule
$j$ & $v_{M4,j}$ & $A^+_j$ & $(v_{M4,j}-A^+_j)^2$ & $A^-_j$ & $(v_{M4,j}-A^-_j)^2$ \\
\midrule
$d^\star$ & 0.09519 & 0.11296 & 0.0003157 & 0.06333 & 0.0010153 \\
$T_{pers}$ & 0.04689 & 0.03030 & 0.0002752 & 0.10148 & 0.0029804 \\
$T_{mis}$ & 0.04599 & 0.03893 & 0.0000498 & 0.04784 & 0.0000034 \\
$N_{stop}$ & 0.00288 & 0.00036 & 0.0000063 & 0.14261 & 0.0195234 \\
$\sigma_v$ & 0.04558 & 0.01548 & 0.0009059 & 0.09575 & 0.0025171 \\
$a_{max}$ & 0.04211 & 0.03837 & 0.0000140 & 0.05090 & 0.0000772 \\
$\eta_s$ & 0.07251 & 0.07251 & 0.0000000 & 0.05800 & 0.0002103 \\
\midrule
\multicolumn{3}{r}{Suma:} & 0.0015669 & & 0.0263271 \\
\multicolumn{3}{r}{Raíz:} & $D^+_{M4} = 0.03958$ & & $D^-_{M4} = 0.16226$ \\
\bottomrule
\end{tabular}
\end{table}

\subsection{Paso 7: Closeness coefficient para todas las alternativas}

\begin{equation}
CC_{M4} = \frac{D^-_{M4}}{D^+_{M4} + D^-_{M4}} = \frac{0.16226}{0.03958 + 0.16226} = \frac{0.16226}{0.20184} = \boxed{0.8039}
\end{equation}

Aplicando el mismo procedimiento a las cinco alternativas:

\begin{table}[H]
\centering
\caption{Resultados TOPSIS Análisis A (sin penalización).}
\renewcommand{\arraystretch}{1.3}
\begin{tabular}{lcccc}
\toprule
\rowcolor{primary}
\textcolor{white}{\textbf{Alt.}} &
\textcolor{white}{\textbf{$D^+$}} &
\textcolor{white}{\textbf{$D^-$}} &
\textcolor{white}{\textbf{$CC$}} &
\textcolor{white}{\textbf{Rank}} \\
\midrule
M0 & 0.08771 & 0.16204 & 0.6488 & 3 \\
M1 & 0.17106 & 0.03065 & 0.1519 & 5 \\
M2 & 0.10039 & 0.11139 & 0.5260 & 4 \\
M3 & 0.04216 & 0.17221 & 0.8033 & 2 \\
M4 & 0.03958 & 0.16226 & 0.8039 & 1 \\
\bottomrule
\end{tabular}
\end{table}

\textbf{Ranking final (Análisis A):} M4 > M3 > M0 > M2 > M1

\newpage
%==================================================================
\section{Análisis B: con penalización de fallos}
%==================================================================

Para penalizar el bloqueo operacional, las corridas fallidas se reemplazan por el peor valor observado en el dataset para cada criterio. La matriz resultante penaliza directamente a M1 y M2 en métricas continuas.

\subsection{Matriz de decisión penalizada}

\begin{table}[H]
\centering
\caption{Matriz $\mathbf{X}$ con penalización de fallos (Análisis B).}
\renewcommand{\arraystretch}{1.3}
\small
\begin{tabular}{lccccccc}
\toprule
\rowcolor{primary}
\textcolor{white}{\textbf{Alt.}} &
\textcolor{white}{$d^\star$} &
\textcolor{white}{$T_{pers}$} &
\textcolor{white}{$T_{mis}$} &
\textcolor{white}{$N_{stop}$} &
\textcolor{white}{$\sigma_v$} &
\textcolor{white}{$a_{max}$} &
\textcolor{white}{$\eta_s$} \\
\midrule
M0 & 0.5627 & 25.3652 & 93.6306 & 0.2000 & 0.0663 & 12.4303 & 1.0000 \\
M1 & 0.6709 & 31.6174 & 122.3384 & 19.1250 & 0.4029 & 15.2172 & 0.8000 \\
M2 & 0.6668 & 25.3020 & 123.4614 & 13.7750 & 0.4283 & 15.0055 & 0.8000 \\
M3 & 1.0038 & 7.5740 & 115.0431 & 0.0250 & 0.2428 & 9.3710 & 1.0000 \\
M4 & 0.8459 & 11.7202 & 110.6021 & 0.2000 & 0.1952 & 10.2838 & 1.0000 \\
\bottomrule
\end{tabular}
\end{table}

\subsection{Resultados Análisis B}

Aplicando los mismos 7 pasos del algoritmo TOPSIS sobre la matriz penalizada:

\begin{table}[H]
\centering
\caption{Resultados TOPSIS Análisis B (con penalización).}
\renewcommand{\arraystretch}{1.3}
\begin{tabular}{lcccc}
\toprule
\rowcolor{primary}
\textcolor{white}{\textbf{Alt.}} &
\textcolor{white}{\textbf{$D^+$}} &
\textcolor{white}{\textbf{$D^-$}} &
\textcolor{white}{\textbf{$CC$}} &
\textcolor{white}{\textbf{Rank}} \\
\midrule
M0 & 0.07512 & 0.14798 & 0.6633 & 3 \\
M1 & 0.16732 & 0.01386 & 0.0765 & 5 \\
M2 & 0.13926 & 0.04086 & 0.2268 & 4 \\
M3 & 0.04048 & 0.15828 & 0.7963 & 2 \\
M4 & 0.03726 & 0.14988 & 0.8009 & 1 \\
\bottomrule
\end{tabular}
\end{table}

\textbf{Ranking final (Análisis B):} M4 > M3 > M0 > M2 > M1

\subsection{Comparación entre análisis}

\begin{table}[H]
\centering
\caption{Comparación de $CC$ entre Análisis A (sin pen.) y B (con pen.).}
\renewcommand{\arraystretch}{1.3}
\begin{tabular}{lccccc}
\toprule
\rowcolor{primary}
\textcolor{white}{\textbf{Método}} &
\textcolor{white}{\textbf{$CC$ (A)}} &
\textcolor{white}{\textbf{$CC$ (B)}} &
\textcolor{white}{\textbf{$\Delta CC$}} &
\textcolor{white}{\textbf{Rank A}} &
\textcolor{white}{\textbf{Rank B}} \\
\midrule
M0 & 0.6488 & 0.6633 & +0.0145 & 3 & 3 \\
M1 & 0.1519 & 0.0765 & -0.0755 & 5 & 5 \\
M2 & 0.5260 & 0.2268 & -0.2991 & 4 & 4 \\
M3 & 0.8033 & 0.7963 & -0.0070 & 2 & 2 \\
M4 & 0.8039 & 0.8009 & -0.0030 & 1 & 1 \\
\bottomrule
\end{tabular}
\end{table}

\textbf{Observaciones clave:}
\begin{itemize}
    \item M4 mantiene Rank \#1 en ambos análisis ($CC = 0.804$ y $0.801$).
    \item M3 mantiene Rank \#2 muy cerca de M4 (diferencia $\Delta CC \approx 0.005$).
    \item M1 y M2 son penalizados drásticamente en el Análisis B: $CC_{M2}$ pasa de $0.526$ a $0.227$; $CC_{M1}$ de $0.152$ a $0.077$.
    \item El ranking ordinal se conserva entre A y B, pero la separación entre los métodos sociales continuos (M3, M4) y los binarios (M1, M2) se amplifica con la penalización.
\end{itemize}

%==================================================================
\section{Conclusiones}
%==================================================================

\begin{enumerate}[leftmargin=*]
    \item TOPSIS posiciona consistentemente a M4 como Rank \#1 en ambos análisis (con y sin penalización), seguido muy de cerca por M3.
    \item La diferencia entre $CC_{M4}$ y $CC_{M3}$ es marginal (alrededor de 0.0006-0.0046). El Documento 6 (Sensibilidad Monte Carlo) explora cuán robusto es este orden ante perturbaciones de los pesos.
    \item Los métodos basados en supervisores binarios (M1, M2) reciben los $CC$ más bajos. La penalización de fallos amplifica esta penalización.
    \item Para una afirmación robusta de superioridad de M4, debe combinarse el ranking TOPSIS con: (i) la dominancia significativa en $\sigma_v$ (ver Documento 3), y (ii) el análisis de sensibilidad de pesos (Documento 6).
\end{enumerate}



\chapter{Analisis de sensibilidad Monte Carlo}
\label{chap:montecarlo}




%==================================================================
\section{Marco teórico del análisis de sensibilidad}
%==================================================================

\subsection{Motivación}

El método TOPSIS produce un ranking que depende crucialmente del vector de pesos $\mathbf{w}$. Si el ranking es \emph{sensible} a pequeñas variaciones de $\mathbf{w}$, las conclusiones del análisis multicriterio son frágiles. Si es \emph{robusto}, las conclusiones son defendibles incluso si los pesos exactos están sujetos a controversia.

\subsection{Formalización del problema}

Sea $\mathbf{w} = (w_1, w_2, \ldots, w_n)$ el vector de pesos nominales con $\sum_j w_j = 1$. Se define una perturbación multiplicativa:
\begin{equation}
w'_j = w_j (1 + \epsilon_j), \qquad \epsilon_j \sim \mathcal{U}(-\epsilon_{\max}, +\epsilon_{\max})
\label{doc6:eq:pertubation}
\end{equation}

donde $\epsilon_{\max}$ controla la amplitud de la perturbación. En este estudio se evalúan dos niveles:
\begin{itemize}
    \item $\epsilon_{\max} = 0.10$ (perturbación moderada, $\pm 10\%$)
    \item $\epsilon_{\max} = 0.20$ (perturbación amplia, $\pm 20\%$)
\end{itemize}

\subsection{Renormalización}

Después de aplicar las perturbaciones independientes, $\sum_j w'_j \neq 1$ en general. Se renormaliza:
\begin{equation}
\boxed{\hat{w}_j = \frac{w'_j}{\sum_{k=1}^{n} w'_k}}
\label{doc6:eq:renorm}
\end{equation}

Esta renormalización garantiza que la condición fundamental $\sum_j \hat{w}_j = 1$ se preserve, pero introduce correlación negativa entre los pesos perturbados (si uno crece, los demás deben encoger relativamente al renormalizar).

\subsection{Procedimiento Monte Carlo}

\begin{enumerate}[leftmargin=*]
    \item Para $s = 1, 2, \ldots, S$ (con $S = 1000$ muestras):
    \begin{enumerate}
        \item Generar $\epsilon_j^{(s)} \sim \mathcal{U}(-\epsilon_{\max}, +\epsilon_{\max})$ independientes para $j = 1, \ldots, n$.
        \item Calcular $w'_j{}^{(s)} = w_j (1 + \epsilon_j^{(s)})$.
        \item Renormalizar: $\hat{w}_j^{(s)} = w'_j{}^{(s)} / \sum_k w'_k{}^{(s)}$.
        \item Aplicar TOPSIS completo con $\hat{\mathbf{w}}^{(s)}$.
        \item Registrar $CC_i^{(s)}$ y $\text{rank}_i^{(s)}$ para cada alternativa $i$.
    \end{enumerate}
    \item Computar estadísticas agregadas sobre las $S$ muestras:
    \begin{itemize}
        \item Frecuencia con que cada alternativa ocupa Rank \#1: $f_i^{(1)} = \#\{s : \text{rank}_i^{(s)} = 1\}/S$.
        \item Frecuencia con que está en Top-2: $f_i^{(\text{top-2})} = \#\{s : \text{rank}_i^{(s)} \leq 2\}/S$.
        \item Media y desviación estándar de $CC_i$ y de $\text{rank}_i$.
    \end{itemize}
\end{enumerate}

\subsection{Configuración del experimento}

\begin{table}[H]
\centering
\caption{Parámetros del experimento Monte Carlo.}
\renewcommand{\arraystretch}{1.3}
\begin{tabular}{lc}
\toprule
\textbf{Parámetro} & \textbf{Valor} \\
\midrule
Número de muestras $S$ & 1000 \\
Niveles de perturbación $\epsilon_{\max}$ & $\{0.10, 0.20\}$ \\
Matrices evaluadas & sin pen., con pen. \\
Total de configuraciones & $2 \times 2 = 4$ \\
Total de evaluaciones TOPSIS & $4 \times 1000 = 4000$ \\
Semilla pseudo-aleatoria & 42 \\
\bottomrule
\end{tabular}
\end{table}

\newpage
%==================================================================
\section{Resultados Monte Carlo}
%==================================================================

\subsection{Configuración 1: $\epsilon = \pm 10\%$, sin penalización}

\subsubsection*{Frecuencias de ranking ($\epsilon = \pm 10\%$, sin penalización)}
\begin{table}[H]
\centering
\caption{Estadísticos sobre 1000 muestras Monte Carlo ($\epsilon = \pm 10\%$, sin penalización).}
\renewcommand{\arraystretch}{1.3}
\begin{tabular}{lcccccc}
\toprule
\rowcolor{primary}
\textcolor{white}{\textbf{Mét.}} & \textcolor{white}{\textbf{Rank \#1 (\%)}} & \textcolor{white}{\textbf{Top-2 (\%)}} & \textcolor{white}{\textbf{Mean rank}} & \textcolor{white}{\textbf{SD rank}} & \textcolor{white}{\textbf{Mean CC}} & \textcolor{white}{\textbf{SD CC}} \\
\midrule
M0 & 0.0 & 0.0 & 3.00 & 0.00 & 0.6480 & 0.0147 \\
M1 & 0.0 & 0.0 & 5.00 & 0.00 & 0.1524 & 0.0070 \\
M2 & 0.0 & 0.0 & 4.00 & 0.00 & 0.5255 & 0.0124 \\
M3 & 49.0 & 100.0 & 1.51 & 0.50 & 0.8030 & 0.0105 \\
M4 & 51.0 & 100.0 & 1.49 & 0.50 & 0.8033 & 0.0084 \\
\bottomrule
\end{tabular}
\end{table}

\subsubsection*{Interpretación}

Con perturbación moderada y sin penalización de fallos:
\begin{itemize}
    \item M4 ocupa Rank \#1 en 51.0\% de las muestras.
    \item M3 ocupa Rank \#1 en 49.0\% de las muestras.
    \item Ambos están en Top-2 en 100.0\% y 100.0\% respectivamente.
    \item M0, M1 y M2 nunca alcanzan Rank \#1 ni Top-2 (frecuencias 0\%).
\end{itemize}

La diferencia entre M4 y M3 sin penalización es muy estrecha. La elección entre ambos depende sensiblemente de los pesos exactos.

\subsection{Configuración 2: $\epsilon = \pm 10\%$, con penalización}

\subsubsection*{Frecuencias de ranking ($\epsilon = \pm 10\%$, con penalización)}
\begin{table}[H]
\centering
\caption{Estadísticos sobre 1000 muestras Monte Carlo ($\epsilon = \pm 10\%$, con penalización).}
\renewcommand{\arraystretch}{1.3}
\begin{tabular}{lcccccc}
\toprule
\rowcolor{primary}
\textcolor{white}{\textbf{Mét.}} & \textcolor{white}{\textbf{Rank \#1 (\%)}} & \textcolor{white}{\textbf{Top-2 (\%)}} & \textcolor{white}{\textbf{Mean rank}} & \textcolor{white}{\textbf{SD rank}} & \textcolor{white}{\textbf{Mean CC}} & \textcolor{white}{\textbf{SD CC}} \\
\midrule
M0 & 0.0 & 0.0 & 3.00 & 0.00 & 0.6626 & 0.0128 \\
M1 & 0.0 & 0.0 & 5.00 & 0.00 & 0.0765 & 0.0040 \\
M2 & 0.0 & 0.0 & 4.00 & 0.00 & 0.2268 & 0.0046 \\
M3 & 21.9 & 100.0 & 1.78 & 0.41 & 0.7967 & 0.0103 \\
M4 & 78.1 & 100.0 & 1.22 & 0.41 & 0.8009 & 0.0076 \\
\bottomrule
\end{tabular}
\end{table}

\subsubsection*{Interpretación}

Con penalización de fallos, M4 domina claramente:
\begin{itemize}
    \item M4 ocupa Rank \#1 en 78.1\% de las muestras (vs 51.0\% sin penalización).
    \item M3 cae a 21.9\% en Rank \#1.
    \item Ambos siguen en Top-2 en 100\% de las muestras.
\end{itemize}

La penalización de fallos amplifica la ventaja de M4 sobre M3, ya que M3 también tiene tasa de éxito 100\% pero pierde frente a M4 en suavidad cinemática y eficiencia tras la penalización indirecta.

\subsection{Configuración 3: $\epsilon = \pm 20\%$, sin penalización}

\subsubsection*{Frecuencias de ranking ($\epsilon = \pm 20\%$, sin penalización)}
\begin{table}[H]
\centering
\caption{Estadísticos sobre 1000 muestras Monte Carlo ($\epsilon = \pm 20\%$, sin penalización).}
\renewcommand{\arraystretch}{1.3}
\begin{tabular}{lcccccc}
\toprule
\rowcolor{primary}
\textcolor{white}{\textbf{Mét.}} & \textcolor{white}{\textbf{Rank \#1 (\%)}} & \textcolor{white}{\textbf{Top-2 (\%)}} & \textcolor{white}{\textbf{Mean rank}} & \textcolor{white}{\textbf{SD rank}} & \textcolor{white}{\textbf{Mean CC}} & \textcolor{white}{\textbf{SD CC}} \\
\midrule
M0 & 0.0 & 0.0 & 3.00 & 0.00 & 0.6489 & 0.0275 \\
M1 & 0.0 & 0.0 & 5.00 & 0.00 & 0.1523 & 0.0132 \\
M2 & 0.0 & 0.0 & 4.00 & 0.00 & 0.5252 & 0.0246 \\
M3 & 48.8 & 100.0 & 1.51 & 0.50 & 0.8031 & 0.0210 \\
M4 & 51.2 & 100.0 & 1.49 & 0.50 & 0.8031 & 0.0162 \\
\bottomrule
\end{tabular}
\end{table}

\subsection{Configuración 4: $\epsilon = \pm 20\%$, con penalización}

\subsubsection*{Frecuencias de ranking ($\epsilon = \pm 20\%$, con penalización)}
\begin{table}[H]
\centering
\caption{Estadísticos sobre 1000 muestras Monte Carlo ($\epsilon = \pm 20\%$, con penalización).}
\renewcommand{\arraystretch}{1.3}
\begin{tabular}{lcccccc}
\toprule
\rowcolor{primary}
\textcolor{white}{\textbf{Mét.}} & \textcolor{white}{\textbf{Rank \#1 (\%)}} & \textcolor{white}{\textbf{Top-2 (\%)}} & \textcolor{white}{\textbf{Mean rank}} & \textcolor{white}{\textbf{SD rank}} & \textcolor{white}{\textbf{Mean CC}} & \textcolor{white}{\textbf{SD CC}} \\
\midrule
M0 & 0.0 & 0.0 & 3.00 & 0.00 & 0.6621 & 0.0253 \\
M1 & 0.0 & 0.0 & 5.00 & 0.00 & 0.0765 & 0.0077 \\
M2 & 0.0 & 0.0 & 4.00 & 0.00 & 0.2267 & 0.0093 \\
M3 & 35.8 & 100.0 & 1.64 & 0.48 & 0.7974 & 0.0205 \\
M4 & 64.2 & 100.0 & 1.36 & 0.48 & 0.8007 & 0.0152 \\
\bottomrule
\end{tabular}
\end{table}

\newpage
%==================================================================
\section{Síntesis comparativa}
%==================================================================

\subsection{Frecuencias de Rank \#1}

\begin{table}[H]
\centering
\caption{Frecuencia de Rank \#1 (\%) por configuración.}
\renewcommand{\arraystretch}{1.3}
\begin{tabular}{lcccc}
\toprule
\rowcolor{primary}
\textcolor{white}{\textbf{Método}} &
\textcolor{white}{\textbf{$\epsilon=10\%$ s.p.}} &
\textcolor{white}{\textbf{$\epsilon=10\%$ c.p.}} &
\textcolor{white}{\textbf{$\epsilon=20\%$ s.p.}} &
\textcolor{white}{\textbf{$\epsilon=20\%$ c.p.}} \\
\midrule
M0 & 0.0 & 0.0 & 0.0 & 0.0 \\
M1 & 0.0 & 0.0 & 0.0 & 0.0 \\
M2 & 0.0 & 0.0 & 0.0 & 0.0 \\
M3 & 49.0 & 21.9 & 48.8 & 35.8 \\
M4 & 51.0 & 78.1 & 51.2 & 64.2 \\
\bottomrule
\end{tabular}
\end{table}

\textbf{Notación:} s.p. = sin penalización; c.p. = con penalización.

\subsection{Frecuencias de Top-2}

\begin{table}[H]
\centering
\caption{Frecuencia de aparición en Top-2 (\%) por configuración.}
\renewcommand{\arraystretch}{1.3}
\begin{tabular}{lcccc}
\toprule
\rowcolor{primary}
\textcolor{white}{\textbf{Método}} &
\textcolor{white}{\textbf{$\epsilon=10\%$ s.p.}} &
\textcolor{white}{\textbf{$\epsilon=10\%$ c.p.}} &
\textcolor{white}{\textbf{$\epsilon=20\%$ s.p.}} &
\textcolor{white}{\textbf{$\epsilon=20\%$ c.p.}} \\
\midrule
M0 & 0.0 & 0.0 & 0.0 & 0.0 \\
M1 & 0.0 & 0.0 & 0.0 & 0.0 \\
M2 & 0.0 & 0.0 & 0.0 & 0.0 \\
M3 & 100.0 & 100.0 & 100.0 & 100.0 \\
M4 & 100.0 & 100.0 & 100.0 & 100.0 \\
\bottomrule
\end{tabular}
\end{table}

\subsection{Comportamiento del rango medio}

\begin{table}[H]
\centering
\caption{Rango medio (1=mejor, 5=peor) por configuración.}
\renewcommand{\arraystretch}{1.3}
\begin{tabular}{lcccc}
\toprule
\rowcolor{primary}
\textcolor{white}{\textbf{Método}} &
\textcolor{white}{\textbf{$\epsilon=10\%$ s.p.}} &
\textcolor{white}{\textbf{$\epsilon=10\%$ c.p.}} &
\textcolor{white}{\textbf{$\epsilon=20\%$ s.p.}} &
\textcolor{white}{\textbf{$\epsilon=20\%$ c.p.}} \\
\midrule
M0 & 3.00 & 3.00 & 3.00 & 3.00 \\
M1 & 5.00 & 5.00 & 5.00 & 5.00 \\
M2 & 4.00 & 4.00 & 4.00 & 4.00 \\
M3 & 1.51 & 1.78 & 1.51 & 1.64 \\
M4 & 1.49 & 1.22 & 1.49 & 1.36 \\
\bottomrule
\end{tabular}
\end{table}

%==================================================================
\section{Conclusiones del análisis de sensibilidad}
%==================================================================

\begin{enumerate}[leftmargin=*]
    \item \textbf{Robustez Top-2:} M3 y M4 ocupan Top-2 en \emph{el 100\% de las 4000 evaluaciones} realizadas (1000 muestras × 4 configuraciones). Ningún otro método aparece nunca en Top-2 bajo ninguna combinación de pesos perturbados. Esta es la conclusión más fuerte del estudio multicriterio.
    
    \item \textbf{Dominancia condicional de M4:} M4 domina claramente cuando se penalizan los fallos ($78.1\%$ con $\epsilon=10\%$, $64.2\%$ con $\epsilon=20\%$). Sin penalización, la competencia M4 vs M3 es muy estrecha ($\approx 51/49$), lo cual refleja que ambos métodos socialmente activos continuos producen rankings muy similares en las métricas básicas.
    
    \item \textbf{El efecto $\epsilon$:} aumentar la perturbación de $\pm 10\%$ a $\pm 20\%$ no cambia las conclusiones cualitativas. M4 sigue dominando con penalización (de $78.1\%$ baja a $64.2\%$), y la ventaja sobre M3 sigue presente.
    
    \item \textbf{Estabilidad de M0:} M0 (NavMesh puro) presenta rango medio estable cerca de 3 en todas las configuraciones, ocupando consistentemente la posición intermedia.
    
    \item \textbf{Posiciones inferiores robustas para M1 y M2:} M2 está siempre en posición 4 y M1 siempre en posición 5 (con penalización aún más marcada). El bloqueo operacional en E4 los condena al fondo del ranking en cualquier configuración de pesos razonable.
    
    \item \textbf{Implicación práctica:} La afirmación \emph{"M4 es la mejor o segunda mejor opción para navegación social en gemelos digitales agrícolas"} es defendible con confianza estadística altísima. La afirmación más fuerte \emph{"M4 es estrictamente mejor que M3"} es defendible cuando las tasas de fallo se consideran explícitamente (Análisis B); sin penalización, ambos métodos son prácticamente equivalentes.
\end{enumerate}



\part{Sintesis, conclusiones y trabajo futuro}

\chapter{Sintesis integradora y discusion}
\label{chap:sintesis}

\section{Resultados clave por hipotesis}

\begin{table}[H]
\centering
\caption{Sintesis del veredicto estadistico por hipotesis.}
\renewcommand{\arraystretch}{1.4}
\small
\begin{tabular}{lp{4cm}p{4cm}l}
\toprule
\rowcolor{primary}
\textcolor{white}{\textbf{ID}} & \textcolor{white}{\textbf{Hipotesis}} & \textcolor{white}{\textbf{Evidencia}} & \textcolor{white}{\textbf{Veredicto}} \\
\midrule
H1 & M0 invade zona intima & K-W $T_{intima}$, $p<0.001$ & Confirmada \\
H2 & M1/M2 mejoran $d^\star$ en E2/E3 & M-W parcial & Parcial \\
H3 & M3 mayor $d^\star$ global & $\bar{d}^\star_{M3}=1.004$ m & Confirmada \\
H4 & M4 mejora $\sigma_v$ vs todos & K-W $H=164.97$, Cliff large & Confirmada (fuerte) \\
H5 & M4 reduce $N_{stop}$ vs M1/M2 & Cliff $\delta < -0.98$ & Confirmada \\
H6 & M4 mejora $\eta_s$ en E4 & Fisher $p=7.14\times 10^{-4}$ & Confirmada (fuerte) \\
\bottomrule
\end{tabular}
\end{table}

\section{Posicionamiento del aporte}

A la luz de las cinco brechas identificadas en el estado del arte, las contribuciones del estudio se sintetizan en cuatro puntos:

\begin{enumerate}[leftmargin=*]
    \item \textbf{Comparativa exhaustiva de cinco arquitecturas de control} en un gemelo digital agricola Unity-MATLAB con 184 corridas.
    \item \textbf{Evaluacion inferencial completa} con tests no parametricos rigurosos (Kruskal-Wallis, Mann-Whitney con Holm-Bonferroni y Cliff's $\delta$, Fisher exacto, $\chi^2$ con Cramer's V).
    \item \textbf{Ranking multicriterio TOPSIS con analisis de sensibilidad Monte Carlo} (1000 muestras x 4 configuraciones), demostrando robustez del Top-2 (M3, M4) en 100\% de las 4000 evaluaciones.
    \item \textbf{Posicionamiento honesto del aporte}: M4 no es novedoso en sus componentes individuales, sino en la \emph{integracion sistematica} de NavMesh global, ORCA local y supervisor social continuo con correccion de bloqueo operacional ($\alpha_{\min} > 0$).
\end{enumerate}

\chapter{Limitaciones y trabajo futuro}
\label{chap:limitaciones}

\section{Limitaciones reconocidas}

\begin{itemize}[leftmargin=*]
    \item \textbf{Toda la evaluacion es en simulacion.} La transferencia a hardware fisico requiere caracterizacion adicional bajo perturbaciones reales.
    \item \textbf{Agentes humanos simulados con patrones predefinidos.} No incluyen comportamientos cooperativos o adversariales explicitos.
    \item \textbf{Parametrizacion del simulador.} Se inspira en literatura agricola pero no constituye mediciones de campo.
    \item \textbf{Cinco metodos representativos.} Otros metodos como TEB humano-aware o MPC con Social-LSTM no se incluyeron.
    \item \textbf{Contexto agricola estructurado.} Las conclusiones aplican al ambito de vinedos en hileras.
\end{itemize}

\section{Trabajo futuro}

\begin{itemize}[leftmargin=*]
    \item Transferencia a hardware fisico con analisis Sim2Real explicito.
    \item Integracion con TEB humano-aware y MPC predictivo.
    \item Estudio de aceptacion humana con experimentos de campo (SUS, NASA-TLX).
    \item Extension a entornos multi-robot cooperativos.
    \item Analisis de seguridad funcional bajo ISO 18497-2 con sensorica certificada.
    \item Fijacion explicita de semillas pseudo-aleatorias para replicacion bit-exacta.
\end{itemize}

\appendix

\chapter{Material suplementario}

\section{Archivos de datos disponibles}

\begin{itemize}[leftmargin=*]
    \item \texttt{master\_log\_clean.csv}: 184 filas x 26 columnas.
    \item \texttt{M\{x\}\_E\{y\}\_run\{n\}.csv}: 184 archivos con telemetria a 10 Hz.
    \item \texttt{master\_calculos.xlsx}: Excel master con 9 hojas y formulas.
    \item Scripts Python que reproducen los tests estadisticos.
    \item Scene Unity y scripts del controlador.
\end{itemize}

\section{Software y versiones}

\begin{itemize}
    \item Unity 2022.3 LTS con NavMesh Components.
    \item MATLAB R2023b con Statistics and Machine Learning Toolbox.
    \item Python 3.12 con NumPy 1.26, SciPy 1.14, Pandas 2.2.
    \item LaTeX TeX Live 2023.
\end{itemize}

\chapter{Reproducibilidad de los calculos}

Todos los calculos estadisticos del presente trabajo se verifican cruzadamente entre:
\begin{enumerate}
    \item \texttt{scipy.stats} (Python).
    \item Calculos manuales paso a paso documentados en los Capitulos \ref{chap:descriptiva}-\ref{chap:montecarlo}.
\end{enumerate}

Las diferencias maximas entre implementaciones son del orden de $10^{-10}$. La semilla pseudo-aleatoria del analisis Monte Carlo es 42.

\backmatter

\chapter*{Referencias bibliograficas}
\addcontentsline{toc}{chapter}{Referencias bibliograficas}

\begin{enumerate}[label={[\arabic*]}, leftmargin=*]

\item \label{Hall1966} E.~T. Hall, \emph{The Hidden Dimension}. New York: Doubleday, 1966.

\item \label{Aiello1987} J.~R. Aiello, ``Human spatial behavior,'' in \emph{Handbook of Environmental Psychology}, D. Stokols and I. Altman, Eds. New York: Wiley, 1987, pp. 389--504.

\item \label{AielloAiello1974} J.~R. Aiello and T. D. Aiello, ``The development of personal space: Proxemic behavior of children 6 through 16,'' \emph{Human Ecology}, vol. 2, no. 3, pp. 177--189, 1974.

\item \label{Mead2017} R. Mead and M.~J. Matarić, ``Autonomous human-robot proxemics: Socially aware navigation based on interaction potential,'' \emph{Autonomous Robots}, vol. 41, no. 5, pp. 1189--1201, 2017.

\item \label{Pacchierotti2005} E. Pacchierotti, H.~I. Christensen, and P. Jensfelt, ``Human-robot embodied interaction in hallway settings: A pilot user study,'' in \emph{IEEE International Workshop on Robot and Human Interactive Communication (ROMAN)}, 2005, pp. 164--171.

\item \label{Walters2005} M.~L. Walters et al., ``The influence of subjects' personality traits on personal spatial zones in a human-robot interaction experiment,'' in \emph{IEEE International Workshop on Robot and Human Interactive Communication (ROMAN)}, 2005, pp. 347--352.

\item \label{RiosMartinez2015} J. Rios-Martinez, A. Spalanzani, and C. Laugier, ``From proxemics theory to socially-aware navigation: A survey,'' \emph{International Journal of Social Robotics}, vol. 7, no. 2, pp. 137--153, 2015.

\item \label{Neggers2024} M.~M.~E. Neggers et al., ``Determining shape and size of personal space of a human when passed by a robot,'' \emph{International Journal of Social Robotics}, vol. 14, no. 2, pp. 561--572, 2022.

\item \label{Neggers2022} M.~M.~E. Neggers, R.~H. Cuijpers, P. A.~M. Ruijten, and W.~A. IJsselsteijn, ``Determining shape and size of personal space of a human when passed by a robot using virtual reality,'' \emph{Frontiers in Virtual Reality}, vol. 3, 768184, 2022.

\item \label{Wengefeld2022} T. Wengefeld, B. Schuetz, G. Girdziunaite, A. Scheidig, and H.-M. Gross, ``The MORPHIA project: First results of a long-term user study in an elderly care scenario from robotic point of view,'' in \emph{ISR Europe 2022}, pp. 1--8, 2022.

\item \label{Singamaneni2024} P.~T. Singamaneni, P. Bachiller-Burgos, L.~J. Manso, A. Garrell, A. Sanfeliu, A. Spalanzani, and R. Alami, ``A survey on socially aware robot navigation: Taxonomy and future challenges,'' \emph{International Journal of Robotics Research}, vol. 43, no. 12, 2024.

\item \label{Mavrogiannis2023} C. Mavrogiannis et al., ``Core challenges of social robot navigation: A survey,'' \emph{ACM Transactions on Human-Robot Interaction}, vol. 12, no. 3, pp. 1--39, 2023.

\item \label{Karwowski2024} J. Karwowski et al., ``Bridging requirements, planning, and evaluation: A review of social robot navigation,'' \emph{Sensors}, vol. 24, no. 9, 2794, 2024.

\item \label{Kruse2013} T. Kruse, A.~K. Pandey, R. Alami, and A. Kirsch, ``Human-aware robot navigation: A survey,'' \emph{Robotics and Autonomous Systems}, vol. 61, no. 12, pp. 1726--1743, 2013.

\item \label{Mirsky2024} R. Mirsky, X. Xiao, J. Hart, and P. Stone, ``Conflict avoidance in social navigation---A survey,'' \emph{ACM Transactions on Human-Robot Interaction}, 2024.

\item \label{Charalampous2017} K. Charalampous, I. Kostavelis, and A. Gasteratos, ``Recent trends in social aware robot navigation: A survey,'' \emph{Robotics and Autonomous Systems}, vol. 93, pp. 85--104, 2017.

\item \label{Lasota2017} P.~A. Lasota, T. Fong, and J.~A. Shah, ``A survey of methods for safe human-robot interaction,'' \emph{Foundations and Trends in Robotics}, vol. 5, no. 4, pp. 261--349, 2017.

\item \label{Dijkstra1959} E.~W. Dijkstra, ``A note on two problems in connexion with graphs,'' \emph{Numerische Mathematik}, vol. 1, no. 1, pp. 269--271, 1959.

\item \label{Hart1968} P.~E. Hart, N.~J. Nilsson, and B. Raphael, ``A formal basis for the heuristic determination of minimum cost paths,'' \emph{IEEE Transactions on Systems Science and Cybernetics}, vol. 4, no. 2, pp. 100--107, 1968.

\item \label{Daniel2010} K. Daniel, A. Nash, S. Koenig, and A. Felner, ``Theta*: Any-angle path planning on grids,'' \emph{Journal of Artificial Intelligence Research}, vol. 39, pp. 533--579, 2010.

\item \label{Harabor2011} D. Harabor and A. Grastien, ``Online graph pruning for pathfinding on grid maps,'' in \emph{AAAI Conference on Artificial Intelligence}, 2011.

\item \label{Dolgov2010} D. Dolgov, S. Thrun, M. Montemerlo, and J. Diebel, ``Path planning for autonomous vehicles in unknown semi-structured environments,'' \emph{International Journal of Robotics Research}, vol. 29, no. 5, pp. 485--501, 2010.

\item \label{Kavraki1996} L.~E. Kavraki, P. Svestka, J.-C. Latombe, and M.~H. Overmars, ``Probabilistic roadmaps for path planning in high-dimensional configuration spaces,'' \emph{IEEE Transactions on Robotics and Automation}, vol. 12, no. 4, pp. 566--580, 1996.

\item \label{LaValle1998} S.~M. LaValle, ``Rapidly-exploring random trees: A new tool for path planning,'' Technical Report, Iowa State University, 1998.

\item \label{LaValle2001} S.~M. LaValle and J.~J. Kuffner, ``Randomized kinodynamic planning,'' \emph{International Journal of Robotics Research}, vol. 20, no. 5, pp. 378--400, 2001.

\item \label{Karaman2011} S. Karaman and E. Frazzoli, ``Sampling-based algorithms for optimal motion planning,'' \emph{International Journal of Robotics Research}, vol. 30, no. 7, pp. 846--894, 2011.

\item \label{Gammell2014} J.~D. Gammell, S.~S. Srinivasa, and T.~D. Barfoot, ``Informed RRT*: Optimal sampling-based path planning focused via direct sampling of an admissible ellipsoidal heuristic,'' in \emph{IROS}, 2014, pp. 2997--3004.

\item \label{Gammell2015} J.~D. Gammell, S.~S. Srinivasa, and T.~D. Barfoot, ``Batch informed trees (BIT*): Sampling-based optimal planning via the heuristically guided search of implicit random geometric graphs,'' in \emph{ICRA}, 2015, pp. 3067--3074.

\item \label{Snook2000} G. Snook, ``Simplified 3D movement and pathfinding using navigation meshes,'' in \emph{Game Programming Gems}, M. DeLoura, Ed. Charles River Media, 2000, pp. 288--304.

\item \label{Mononen2013} M. Mononen, ``Recast \& Detour: Navigation mesh toolset for games,'' Open-source library, 2013. [Online]. Available: \url{https://github.com/recastnavigation/recastnavigation}

\item \label{SanchezIbanez2021} J.~R. Sánchez-Ibáñez, C.~J. Pérez-del-Pulgar, and A. García-Cerezo, ``Path planning for autonomous mobile robots: A review,'' \emph{Sensors}, vol. 21, no. 23, 7898, 2021.

\item \label{Mohanan2018} M.~G. Mohanan and A. Salgoankar, ``A survey of robotic motion planning in dynamic environments,'' \emph{Robotics and Autonomous Systems}, vol. 100, pp. 171--185, 2018.

\item \label{Khatib1986} O. Khatib, ``Real-time obstacle avoidance for manipulators and mobile robots,'' \emph{International Journal of Robotics Research}, vol. 5, no. 1, pp. 90--98, 1986.

\item \label{Connolly1990} C.~I. Connolly, J.~B. Burns, and R. Weiss, ``Path planning using Laplace's equation,'' in \emph{ICRA}, 1990, pp. 2102--2106.

\item \label{Rimon1992} E. Rimon and D.~E. Koditschek, ``Exact robot navigation using artificial potential functions,'' \emph{IEEE Transactions on Robotics and Automation}, vol. 8, no. 5, pp. 501--518, 1992.

\item \label{Fox1997} D. Fox, W. Burgard, and S. Thrun, ``The dynamic window approach to collision avoidance,'' \emph{IEEE Robotics \& Automation Magazine}, vol. 4, no. 1, pp. 23--33, 1997.

\item \label{Liang2024} S. Liang et al., ``Dynamic adaptive dynamic window approach,'' \emph{IEEE Transactions on Robotics}, 2024.

\item \label{Rosmann2012} C. Rösmann, W. Feiten, T. Wösch, F. Hoffmann, and T. Bertram, ``Trajectory modification considering dynamic constraints of autonomous robots,'' in \emph{ROBOTIK}, 2012, pp. 74--79.

\item \label{Rosmann2017} C. Rösmann, F. Hoffmann, and T. Bertram, ``Integrated online trajectory planning and optimization in distinctive topologies,'' \emph{Robotics and Autonomous Systems}, vol. 88, pp. 142--153, 2017.

\item \label{Bermudez2024} G. Bermudez et al., ``Comparative analyses of ROS local planners for quadrupedal locomotion: A study in real and simulated environments,'' in \emph{CLAWAR 2024}, Springer LNNS, 2024.

\item \label{Yuan2022} H. Yuan, H. Li, Y. Zhang, S. Du, L. Yu, and X. Wang, ``Comparison and improvement of local planners on ROS for narrow passages,'' in \emph{HDIS}, 2022, pp. 125--130.

\item \label{Kollmitz2021} M. Kollmitz, T. Koller, J. Boedecker, and W. Burgard, ``Human-aware navigation planner for diverse human-robot contexts,'' in \emph{IROS}, 2021, pp. 6635--6642.

\item \label{Fiorini1993} P. Fiorini and Z. Shiller, ``Motion planning in dynamic environments using the relative velocity paradigm,'' in \emph{ICRA}, 1993, pp. 560--565.

\item \label{Fiorini1998} P. Fiorini and Z. Shiller, ``Motion planning in dynamic environments using velocity obstacles,'' \emph{International Journal of Robotics Research}, vol. 17, no. 7, pp. 760--772, 1998.

\item \label{VandenBerg2008} J. van den Berg, M. Lin, and D. Manocha, ``Reciprocal velocity obstacles for real-time multi-agent navigation,'' in \emph{ICRA}, 2008, pp. 1928--1935.

\item \label{VandenBerg2011} J. van den Berg, S.~J. Guy, M. Lin, and D. Manocha, ``Reciprocal n-body collision avoidance,'' in \emph{Robotics Research}, Springer Tracts in Advanced Robotics, vol. 70, 2011, pp. 3--19.

\item \label{Snape2011} J. Snape, J. van den Berg, S.~J. Guy, and D. Manocha, ``The hybrid reciprocal velocity obstacle,'' \emph{IEEE Transactions on Robotics}, vol. 27, no. 4, pp. 696--706, 2011.

\item \label{Alonso2018} J. Alonso-Mora, P. Beardsley, and R. Siegwart, ``Cooperative collision avoidance for nonholonomic robots,'' \emph{IEEE Transactions on Robotics}, vol. 34, no. 2, pp. 404--420, 2018.

\item \label{London2025} J. London, ``Improved obstacle avoidance for autonomous robots with ORCA-FLC,'' \emph{arXiv preprint} arXiv:2508.06722, 2025.

\item \label{Helbing1995} D. Helbing and P. Molnar, ``Social force model for pedestrian dynamics,'' \emph{Physical Review E}, vol. 51, no. 5, pp. 4282--4286, 1995.

\item \label{Helbing2000} D. Helbing, I. Farkas, and T. Vicsek, ``Simulating dynamical features of escape panic,'' \emph{Nature}, vol. 407, pp. 487--490, 2000.

\item \label{Helbing2002} D. Helbing, I. Farkas, P. Molnar, and T. Vicsek, ``Simulation of pedestrian crowds in normal and evacuation situations,'' in \emph{Pedestrian and Evacuation Dynamics}, M. Schreckenberg and S.~D. Sharma, Eds. Berlin: Springer, 2002, pp. 21--58.

\item \label{Moussaid2009} M. Moussaïd, D. Helbing, S. Garnier, A. Johansson, M. Combe, and G. Theraulaz, ``Experimental study of the behavioural mechanisms underlying self-organization in human crowds,'' \emph{Proceedings of the Royal Society B}, vol. 276, no. 1668, pp. 2755--2762, 2009.

\item \label{Moussaid2010} M. Moussaïd, N. Perozo, S. Garnier, D. Helbing, and G. Theraulaz, ``The walking behaviour of pedestrian social groups and its impact on crowd dynamics,'' \emph{PLOS ONE}, vol. 5, no. 4, e10047, 2010.

\item \label{Pellegrini2009} S. Pellegrini, A. Ess, K. Schindler, and L. van Gool, ``You'll never walk alone: Modeling social behavior for multi-target tracking,'' in \emph{ICCV}, 2009, pp. 261--268.

\item \label{Ferrer2013} G. Ferrer, A. Garrell, and A. Sanfeliu, ``Robot companion: A social-force based approach with human awareness-navigation in crowded environments,'' in \emph{IROS}, 2013, pp. 1688--1694.

\item \label{Ferrer2014} G. Ferrer and A. Sanfeliu, ``Proactive kinodynamic planning using the extended social force model and human motion prediction in urban environments,'' in \emph{IROS}, 2014, pp. 1730--1735.

\item \label{Luber2010} M. Luber, J.~A. Stork, G.~D. Tipaldi, and K.~O. Arras, ``People tracking with human motion predictions from social forces,'' in \emph{ICRA}, 2010, pp. 464--469.

\item \label{Pradalier2003} C. Pradalier, J. Hermosillo, C. Koike, C. Braillon, P. Bessière, and C. Laugier, ``The CyCab: A car-like robot navigating autonomously and safely among pedestrians,'' \emph{Robotics and Autonomous Systems}, vol. 50, no. 1, pp. 51--67, 2005.

\item \label{Sisbot2007} E.~A. Sisbot, L.~F. Marin-Urias, R. Alami, and T. Siméon, ``A human aware mobile robot motion planner,'' \emph{IEEE Transactions on Robotics}, vol. 23, no. 5, pp. 874--883, 2007.

\item \label{Kirby2009} R. Kirby, R. Simmons, and J. Forlizzi, ``Companion: A constraint-optimizing method for person-acceptable navigation,'' in \emph{ROMAN}, 2009, pp. 607--612.

\item \label{Lichtenthaler2012} C. Lichtenthäler, T. Lorenz, and A. Kirsch, ``Influence of legibility on perceived safety in a virtual human-robot path crossing task,'' in \emph{ROMAN}, 2012, pp. 676--681.

\item \label{Trautman2013} P. Trautman and A. Krause, ``Unfreezing the robot: Navigation in dense, interacting crowds,'' in \emph{IROS}, 2010, pp. 797--803.

\item \label{Trautman2015} P. Trautman, J. Ma, R.~M. Murray, and A. Krause, ``Robot navigation in dense human crowds: Statistical models and experimental studies of human-robot cooperation,'' \emph{International Journal of Robotics Research}, vol. 34, no. 3, pp. 335--356, 2015.

\item \label{Alahi2016} A. Alahi, K. Goel, V. Ramanathan, A. Robicquet, L. Fei-Fei, and S. Savarese, ``Social LSTM: Human trajectory prediction in crowded spaces,'' in \emph{CVPR}, 2016, pp. 961--971.

\item \label{Gupta2018} A. Gupta, J. Johnson, L. Fei-Fei, S. Savarese, and A. Alahi, ``Social GAN: Socially acceptable trajectories with generative adversarial networks,'' in \emph{CVPR}, 2018, pp. 2255--2264.

\item \label{Salzmann2020} T. Salzmann, B. Ivanovic, P. Chakravarty, and M. Pavone, ``Trajectron++: Dynamically-feasible trajectory forecasting with heterogeneous data,'' in \emph{ECCV}, 2020, pp. 683--700.

\item \label{Mao2024} W. Mao et al., ``Diffusion-based trajectory prediction for pedestrians in dynamic environments,'' \emph{Pattern Recognition}, 2024.

\item \label{Schulz2003} D. Schulz, W. Burgard, D. Fox, and A.~B. Cremers, ``People tracking with mobile robots using sample-based joint probabilistic data association filters,'' \emph{International Journal of Robotics Research}, vol. 22, no. 2, pp. 99--116, 2003.

\item \label{Brito2019} B. Brito, B. Floor, L. Ferranti, and J. Alonso-Mora, ``Model predictive contouring control for collision avoidance in unstructured dynamic environments,'' \emph{IEEE Robotics and Automation Letters}, vol. 4, no. 4, pp. 4459--4466, 2019.

\item \label{Schoels2020} T. Schoels, P. Rutquist, L. Palmieri, A. Zanelli, K.~O. Arras, and M. Diehl, ``CIAO*: MPC-based safe motion planning in predictable dynamic environments,'' \emph{IFAC-PapersOnLine}, vol. 53, no. 2, pp. 6555--6562, 2020.

\item \label{Kamenev2023} A. Kamenev et al., ``Predictive social navigation in crowded environments using model predictive control,'' \emph{arXiv preprint} arXiv:2309.16838, 2023.

\item \label{Shamsah2024} A. Shamsah, K. Agarwal, N. Katta, A. Raju, S. Kousik, and Y. Zhao, ``Socially acceptable bipedal robot navigation via Social Zonotope Network model predictive control,'' \emph{arXiv preprint} arXiv:2406.17151, 2024.

\item \label{Chen2017} Y.~F. Chen, M. Liu, M. Everett, and J.~P. How, ``Decentralized non-communicating multiagent collision avoidance with deep reinforcement learning,'' in \emph{ICRA}, 2017, pp. 285--292.

\item \label{Chen2019} Y.~F. Chen, M. Everett, M. Liu, and J.~P. How, ``Socially aware motion planning with deep reinforcement learning,'' in \emph{IROS}, 2017, pp. 1343--1350.

\item \label{Everett2018} M. Everett, Y.~F. Chen, and J.~P. How, ``Motion planning among dynamic, decision-making agents with deep reinforcement learning,'' in \emph{IROS}, 2018, pp. 3052--3059.

\item \label{Liu2023} S. Liu, P. Chang, Z. Huang, N. Chakraborty, K. Hong, W. Liang, and K. Driggs-Campbell, ``Intention aware robot crowd navigation with attention-based interaction graph,'' in \emph{ICRA}, 2023.

\item \label{Chen2023} C. Chen, Y. Liu, S. Kreiss, and A. Alahi, ``Crowd-robot interaction: Crowd-aware robot navigation with attention-based deep reinforcement learning,'' in \emph{ICRA}, 2019, pp. 6015--6022.

\item \label{Vasconez2019} J.~P. Vasconez, G.~A. Kantor, and F.~A. Auat-Cheein, ``Human-robot interaction in agriculture: A survey and current challenges,'' \emph{Biosystems Engineering}, vol. 179, pp. 35--48, 2019.

\item \label{Tsolakis2024} N. Tsolakis et al., ``Agricultural robotics: Trends, applications, and current challenges,'' \emph{Computers and Electronics in Agriculture}, 2024.

\item \label{Macenski2020} S. Macenski et al., ``The Marathon 2: A navigation system,'' in \emph{IROS}, 2020, pp. 2718--2725.

\item \label{Sani2024} E. Sani, A. Sgorbissa, and S. Carpin, ``Improving the ROS 2 navigation stack with real-time local costmap updates for agricultural applications,'' \emph{arXiv preprint} arXiv:2407.18535, 2024.

\item \label{Reina2017} G. Reina, A. Milella, R. Rouveure, M. Nielsen, R. Worst, and M.~R. Blas, ``Ambient awareness for agricultural robotic vehicles,'' \emph{Biosystems Engineering}, vol. 146, pp. 114--132, 2017.

\item \label{ISO18497} ISO 18497-1:2024, \emph{Agricultural Machinery and Tractors --- Safety of Highly Automated Agricultural Machines --- Part 1: Principles for Design}. International Organization for Standardization, Geneva, 2024.

\item \label{Bispo2023} R. Bispo, J. Heimann, and L. Wendel, ``Towards safe robotic agricultural applications: Safe navigation system design for a robotic grass-mowing application through the risk management method,'' \emph{Robotics}, vol. 12, no. 3, 63, 2023.

\item \label{Roozemond2024} P. Roozemond et al., ``Safety of automated agricultural machineries: A systematic literature review,'' \emph{Safety}, vol. 9, no. 1, 13, 2023.

\item \label{Spyrosoft2024} Spyrosoft, ``Why agricultural robotics fails without functional safety engineering,'' Industry Whitepaper, 2024.

\item \label{Grieves2014} M. Grieves, ``Digital twin: Manufacturing excellence through virtual factory replication,'' Whitepaper, 2014.

\item \label{Tao2018} F. Tao, J. Cheng, Q. Qi, M. Zhang, H. Zhang, and F. Sui, ``Digital twin-driven product design, manufacturing and service with big data,'' \emph{International Journal of Advanced Manufacturing Technology}, vol. 94, pp. 3563--3576, 2018.

\item \label{Tao2019} F. Tao, Q. Qi, L. Wang, and A.~Y.~C. Nee, ``Digital twins and cyber-physical systems toward smart manufacturing and industry 4.0: Correlation and comparison,'' \emph{Engineering}, vol. 5, no. 4, pp. 653--661, 2019.

\item \label{Han2025} Y. Han, S. Park, and J. Kim, ``Digital twin-driven control platform for industrial robotics,'' \emph{Sensors}, vol. 25, no. 13, 4153, 2025.

\item \label{Chang2025} W. Chang, W. Sun, P. Chen, and H. Xu, ``Construction and validation of a digital twin-driven virtual-reality fusion control platform for industrial robots,'' \emph{Sensors}, vol. 25, no. 13, 4153, 2025.

\item \label{Tobin2017} J. Tobin, R. Fong, A. Ray, J. Schneider, W. Zaremba, and P. Abbeel, ``Domain randomization for transferring deep neural networks from simulation to the real world,'' in \emph{IROS}, 2017, pp. 23--30.

\item \label{GuptaDA2018} A. Gupta et al., ``Robot learning in homes: Improving generalization and reducing dataset bias,'' in \emph{NeurIPS}, 2018.

\item \label{Yan2025} W. Yan et al., ``From simulation to field validation: A digital twin-driven Sim2real transfer approach for strawberry fruit detection and sizing,'' \emph{AgriEngineering}, vol. 7, no. 3, 81, 2025.

\item \label{Mann1947} H.~B. Mann and D.~R. Whitney, ``On a test of whether one of two random variables is stochastically larger than the other,'' \emph{Annals of Mathematical Statistics}, vol. 18, no. 1, pp. 50--60, 1947.

\item \label{Kruskal1952} W.~H. Kruskal and W.~A. Wallis, ``Use of ranks in one-criterion variance analysis,'' \emph{Journal of the American Statistical Association}, vol. 47, no. 260, pp. 583--621, 1952.

\item \label{leCessie2020} S. le Cessie, J.~J. Goeman, and O.~M. Dekkers, ``Who is afraid of non-normal data? Choosing between parametric and non-parametric tests,'' \emph{European Journal of Endocrinology}, vol. 182, no. 2, E1--E3, 2020.

\item \label{Holm1979} S. Holm, ``A simple sequentially rejective multiple test procedure,'' \emph{Scandinavian Journal of Statistics}, vol. 6, no. 2, pp. 65--70, 1979.

\item \label{Benjamini1995} Y. Benjamini and Y. Hochberg, ``Controlling the false discovery rate: A practical and powerful approach to multiple testing,'' \emph{Journal of the Royal Statistical Society B}, vol. 57, no. 1, pp. 289--300, 1995.

\item \label{Dunn1964} O.~J. Dunn, ``Multiple comparisons using rank sums,'' \emph{Technometrics}, vol. 6, no. 3, pp. 241--252, 1964.

\item \label{Cliff1993} N. Cliff, ``Dominance statistics: Ordinal analyses to answer ordinal questions,'' \emph{Psychological Bulletin}, vol. 114, no. 3, pp. 494--509, 1993.

\item \label{Cliff1996} N. Cliff, \emph{Ordinal Methods for Behavioral Data Analysis}. Mahwah, NJ: Lawrence Erlbaum Associates, 1996.

\item \label{Romano2006} J. Romano, J.~D. Kromrey, J. Coraggio, and J. Skowronek, ``Appropriate statistics for ordinal level data: Should we really be using t-test and Cohen's d for evaluating group differences on the NSSE and other surveys?,'' in \emph{Annual Meeting of the Florida Association of Institutional Research}, 2006.

\item \label{Fisher1922} R.~A. Fisher, ``On the interpretation of $\chi^2$ from contingency tables, and the calculation of P,'' \emph{Journal of the Royal Statistical Society}, vol. 85, no. 1, pp. 87--94, 1922.

\item \label{Pearson1900} K. Pearson, ``On the criterion that a given system of deviations from the probable in the case of a correlated system of variables is such that it can be reasonably supposed to have arisen from random sampling,'' \emph{Philosophical Magazine}, vol. 50, no. 302, pp. 157--175, 1900.

\item \label{Cramer1946} H. Cramér, \emph{Mathematical Methods of Statistics}. Princeton: Princeton University Press, 1946.

\item \label{Rea1992} L.~M. Rea and R.~A. Parker, \emph{Designing and Conducting Survey Research}. San Francisco: Jossey-Bass, 1992.

\item \label{VirtanenScipy2020} P. Virtanen et al., ``SciPy 1.0: Fundamental algorithms for scientific computing in Python,'' \emph{Nature Methods}, vol. 17, pp. 261--272, 2020.

\item \label{HwangYoon1981} C.-L. Hwang and K. Yoon, \emph{Multiple Attribute Decision Making: Methods and Applications}. New York: Springer-Verlag, 1981.

\item \label{ChenHwang1992} S.-J. Chen and C.-L. Hwang, \emph{Fuzzy Multiple Attribute Decision Making: Methods and Applications}. Berlin: Springer, 1992.

\item \label{Onder2013} E. Onder and S. Dag, ``A hybrid AHP-TOPSIS model for supplier selection in the airline industry,'' \emph{International Journal of Operations Research}, vol. 10, no. 2, pp. 56--66, 2013.

\item \label{Wu2020} K.-Y. Wu, ``Applying the fuzzy Delphi method to analyze the evaluation indexes for service quality after restaurant ECRM problem-solving,'' \emph{International Journal of Quality \& Reliability Management}, 2020.

\item \label{ElHadrami2023} A. El Hadrami et al., ``New distributed-TOPSIS approach for multi-criteria decision-making problems in a big data context,'' \emph{Journal of Big Data}, vol. 10, 97, 2023.

\item \label{Bhargava2024} A. Bhargava et al., ``A novel TOPSIS framework for multi-criteria decision making with random hypergraphs,'' \emph{Symmetry}, vol. 16, no. 12, 1602, 2024.

\item \label{Olson2004} D.~L. Olson, ``Comparison of weights in TOPSIS models,'' \emph{Mathematical and Computer Modelling}, vol. 40, no. 7--8, pp. 721--727, 2004.

\item \label{Saltelli2008} A. Saltelli et al., \emph{Global Sensitivity Analysis: The Primer}. Chichester: Wiley, 2008.

\item \label{Yazdani2017} M. Yazdani, P. Zarate, A. Coulibaly, and E.~K. Zavadskas, ``A group decision making support system in logistics and supply chain management,'' \emph{Expert Systems with Applications}, vol. 88, pp. 376--392, 2017.

\item \label{Roy1991} B. Roy, ``The outranking approach and the foundations of ELECTRE methods,'' \emph{Theory and Decision}, vol. 31, no. 1, pp. 49--73, 1991.

\item \label{Brans1986} J.-P. Brans, P. Vincke, and B. Mareschal, ``How to select and how to rank projects: The PROMETHEE method,'' \emph{European Journal of Operational Research}, vol. 24, no. 2, pp. 228--238, 1986.

\item \label{Opricovic2004} S. Opricovic and G.-H. Tzeng, ``Compromise solution by MCDM methods: A comparative analysis of VIKOR and TOPSIS,'' \emph{European Journal of Operational Research}, vol. 156, no. 2, pp. 445--455, 2004.

\item \label{Velasquez2013} M. Velasquez and P.~T. Hester, ``An analysis of multi-criteria decision making methods,'' \emph{International Journal of Operations Research}, vol. 10, no. 2, pp. 56--66, 2013.

\item \label{Mathew2023} M. Mathew and H. Tang, ``A comprehensive guide to the TOPSIS method for multi-criteria decision making,'' \emph{Sustainable Social Development}, 2023.

\item \label{Kastner2024} L. Kästner et al., ``Arena-Rosnav 3.0: A comprehensive development and benchmarking platform for navigation,'' \emph{RSS}, 2024.

\item \label{Tsoi2022} N. Tsoi, A. Xiang, P. Yu, S.~S. Sohn, M. Schwager, and M. Vázquez, ``SEAN 2.0: Formalizing and generating social situations for robot navigation,'' \emph{IEEE Robotics and Automation Letters}, vol. 7, no. 4, pp. 11047--11054, 2022.

\item \label{Lerner2007} A. Lerner, Y. Chrysanthou, and D. Lischinski, ``Crowds by example,'' \emph{Computer Graphics Forum}, vol. 26, no. 3, pp. 655--664, 2007.

\item \label{ThomazLeite2016} A. Thomaz, G. Hoffman, and M. Cakmak, ``Computational human-robot interaction,'' \emph{Foundations and Trends in Robotics}, vol. 4, no. 2--3, pp. 105--223, 2016.

\item \label{Robicquet2016} A. Robicquet, A. Sadeghian, A. Alahi, and S. Savarese, ``Learning social etiquette: Human trajectory understanding in crowded scenes,'' in \emph{ECCV}, 2016, pp. 549--565.

\end{enumerate}

\end{document}
